\documentclass[11pt,a4paper]{article}
\usepackage[margin=26mm,headheight=15pt]{geometry}
\usepackage[no-math]{fontspec}
\setmainfont{Latin Modern Roman}
\setsansfont{Latin Modern Sans}
\setmonofont{Latin Modern Mono}
\usepackage{amsmath,amssymb,amsthm,mathtools,mathrsfs}
\usepackage{booktabs,longtable,array,tabularx}
\newcolumntype{L}[1]{>{\raggedright\arraybackslash}p{#1}}
\newcolumntype{Y}{>{\raggedright\arraybackslash}X}
\usepackage{enumitem,microtype,xcolor}
\usepackage[round,authoryear]{natbib}
\usepackage{xurl}
\usepackage{hyperref}
\hypersetup{colorlinks=true,linkcolor=blue!45!black,citecolor=blue!45!black,urlcolor=blue!45!black,pdftitle={The Teichmuller tower: arithmetic images, normalizers, and stratified completions},pdfauthor={GPT-6 Astra (OpenAI, ChatGPT Pro mode)}}
\usepackage{bookmark,fancyhdr}
\pagestyle{fancy}\fancyhf{}
\fancyhead[L]{\small LM3\_TEICHMULLER\_TOWER}
\fancyhead[R]{\small Source and proof audit}
\fancyfoot[C]{\thepage}
\setlength{\parindent}{1.3em}
\setlength{\parskip}{3pt}
\setlength{\emergencystretch}{3em}
\setcounter{tocdepth}{2}
\newcommand{\statement}[3]{\par\medskip\phantomsection\label{st:#1}{\noindent\bfseries\raggedright Statement #1 (#2). #3.\par}\nobreak\noindent\ignorespaces}
\newcommand{\Q}{\mathbb Q}\newcommand{\Z}{\mathbb Z}\newcommand{\C}{\mathbb C}\newcommand{\F}{\mathbb F}
\newcommand{\GQ}{G_{\mathbb Q}}
\newcommand{\wh}[1]{\widehat{#1}}
\DeclareMathOperator{\Inn}{Inn}\DeclareMathOperator{\Aut}{Aut}\DeclareMathOperator{\Out}{Out}\DeclareMathOperator{\Gal}{Gal}
\DeclareMathOperator{\Hom}{Hom}\DeclareMathOperator{\Isom}{Isom}\DeclareMathOperator{\Fun}{Fun}
\DeclareMathOperator{\Stab}{Stab}\DeclareMathOperator{\Norm}{N}\DeclareMathOperator{\im}{im}
\DeclareMathOperator{\Spec}{Spec}\DeclareMathOperator{\Sec}{Sec}
\newcommand{\GT}{\widehat{\mathrm{GT}}}\newcommand{\GRT}{\mathrm{GRT}}
\newcommand{\grt}{\mathfrak{grt}}\newcommand{\et}{\mathrm{\acute et}}
\newcommand{\etale}{\'etale}
\newcommand{\op}{\mathrm{op}}\newcommand{\cts}{\mathrm{cts}}
\newcommand{\pro}{\mathrm{pro}}\newcommand{\id}{\mathrm{id}}
\newcommand{\into}{\hookrightarrow}\newcommand{\onto}{\twoheadrightarrow}
\newcommand{\Man}[3]{\emph{Manuscript:} \S#1, folio #2, scan #3.}
\begin{document}
\begin{center}
{\Large\bfseries The Teichm\"uller tower:\\[3pt] arithmetic images, normalizers, and stratified completions}\par
\vspace{6pt}
{\small Development of \emph{La Longue Marche}, \S\S19, 21, 26--30, 32--36}\par
{\small AI-produced by GPT-6 Astra (OpenAI, ChatGPT Pro mode)}\par
{\small 11 November 2026. No priority is claimed.}
\end{center}
\noindent\textbf{Notation fixed before any comparison.}
We write $n$ for the manuscript's $\nu$. A type is hyperbolic if $2g-2+n>0$.
All profinite homomorphisms are continuous; quotients of profinite groups are by closed normal subgroups. An outer homomorphism to $H$ is a homomorphism modulo conjugation by $H$, unless a smaller conjugating group is explicitly specified. An identification obtained from a geometric base point is invariant only up to the indicated inner conjugations. We never infer injectivity of a profinite completion map from injectivity of the discrete map.
\begin{center}\small
\begin{tabularx}{\textwidth}{@{}L{.19\textwidth}L{.32\textwidth}Y@{}}
\toprule
Our notation & Manuscript notation and scan & Meaning\\\midrule
$A_{g,n}$, $A^!_{g,n}$ & $A_{g,\nu}$, $A^!_{g,\nu}$; 204--205 & Homeomorphisms preserving the puncture set, respectively each puncture.\\
$M^{\pm}_{g,n}$, $M^+_{g,n}$ & $\Gamma_{g,\nu}$, $\Gamma^+_{g,\nu}$, later $\mathcal T$, $\mathcal T^+$; 205, 226, 328 & Extended mapping class group; orientation-preserving subgroup; punctures may be permuted.\\
$P_{g,n}$ & $\Gamma^{!+}_{g,\nu}$; 205, 226 & Pure orientation-preserving mapping class group. Part F uses $G_n=P_{0,n}$.\\
$\Pi_{g,n}$, $\pi_{g,n}$ & $\Pi_{g,\nu}$, $\widehat\Pi_{g,\nu}$; 221, 328 & Discrete surface fundamental group and its full profinite completion, with peripheral data.\\
$O_{g,n}$ & $\widehat{\widehat{\mathcal T}}$; 321--324 & $\Out_{\mathrm{lac}}(\pi_{g,n})$, not a second application of completion.\\
$\Sigma_{g,n}$ & $\widehat{\mathcal T}^{+}=\Sigma$; 323--324 & Geometric image in $O_{g,n}$; distinguished from the abstract completion $\widehat{M^+_{g,n}}$.\\
$N^{\mathrm{norm}}$, $\mathfrak G^{\mathrm{norm}}$ & $\mathcal N_\Sigma$, $\boldsymbol\Gamma_\Sigma$; 324--326 & $N_{O}(\Sigma)$ and $N_{O}(\Sigma)/\Sigma$, \S26 version.\\
$N^{\mathrm{arith}}$, $\mathfrak G^{\mathrm{arith}}$ & $N_{g,\nu}$, $\Gamma_{g,\nu}$; 370--371 & Arithmetic monodromy image and its quotient by $\Sigma$, \S29 version.\\
$\mathscr M_{g,[n]}$ & $M_{g,\nu}$; 326--329 & Moduli stack of curves with an unordered reduced degree-$n$ puncture divisor. $\mathscr M_{g,n}$ has labelled punctures.\\
$C_n$, $T_n=C_n^{\op}$ & No manuscript symbol; STRATA F.7 & Exit category, respectively degeneration category. Completion retains these variances.\\
$G_K$, $\chi$ & $\mathbb I\Gamma_K$, $\chi$; 327, 323 & $\Gal(\overline K/K)$ and the cyclotomic multiplier.\\\bottomrule
\end{tabularx}
\end{center}
\noindent\small
Later sections may use locally declared synonymous letters when quoting a source. The labels \textbf{Known}, \textbf{Proved}, \textbf{Refuted}, and \textbf{Open} classify the exact numbered assertion. ``Open'' also marks an exact assertion left unverified by this audit; we distinguish that from an established open problem in its accompanying text. A conditional implication can be \textbf{Proved} while its antecedent remains an unrefereed claim. ``Verified'' in the bibliography means that the cited source was opened, not that its entire proof was independently certified.\normalsize

\clearpage
\begin{abstract}
We distinguish the normalizer quotients introduced in \S26 from the arithmetic-image quotients introduced in \S29. The distinction changes the mathematical content of the conjecture fondamentale: faithfulness modulo geometric monodromy identifies the latter quotients with the absolute Galois group, whereas a description of all compatible profinite tower automorphisms by Galois automorphisms is a separate problem. We give typed comparison maps, audit the claims concerning change of type and reconstruction, and construct a profinite category retaining the boundary groups and their crossing bimodules. We distinguish verified theorems, proofs given here, counterexamples, and unresolved statements, including the exact dependence on the attached unrefereed manuscripts of 2026. A statement register, an editorial reading register, reproducible computations, and a verification report accompany this paper.
\end{abstract}
\tableofcontents
\clearpage
\section{A. The tower and the manuscript's changing definitions}
\subsection{A.1. Surfaces, punctures, and the two finite characters}

\statement{A.1}{Proved}{The permutation and orientation conventions}
Let $X_g$ be a closed connected orientable topological surface and $S_n\subset X_g$ a specified $n$-element set. Give its homeomorphism group the compact-open topology on a homeomorphism and its inverse, and put $U_{g,n}=X_g\setminus S_n$. Then
\[
1\longrightarrow A^!_{g,n}\longrightarrow A_{g,n}
 \longrightarrow\mathfrak S_n\longrightarrow1,
\qquad
1\longrightarrow\pi_0(A^!_{g,n})\longrightarrow M^{\pm}_{g,n}
 \longrightarrow\mathfrak S_n\longrightarrow1.
\]
The orientation character $\epsilon:M^{\pm}_{g,n}\to\{\pm1\}$ is surjective and its kernel is $M^+_{g,n}$. The intersection of the permutation and orientation kernels is $P_{g,n}$. In particular the manuscript's unadorned group is neither pure nor orientation-preserving by definition.
\Man{19}{102r--v}{204--205}; the positive superscript is specified in \S19, folio 113, scan 226.
\begin{proof}
Restriction to the finite set gives the first homomorphism. Every permutation is realized by an orientation-preserving homeomorphism: one realizes transpositions by interchanging two small discs along an embedded arc avoiding the other punctures, and composes. Its kernel is precisely the pointwise stabilizer. A path of homeomorphisms induces a constant permutation, so the identity components of the two groups coincide. Taking component groups therefore gives the second sequence, including surjectivity. Degree on the oriented fundamental class gives $\epsilon$. A reflection model for a genus-$g$ surface, followed by a point-moving homeomorphism, realizes degree $-1$ preserving the prescribed punctures individually. Thus both signs occur even in the pure group. No splitting of the permutation sequence is asserted.

Every homeomorphism of $U_{g,n}$ extends uniquely to its finite end compactification $X_g$: the ends are exactly the deleted punctures, and a homeomorphism permutes ends. This proves the abstract isomorphism in manuscript (5). The extension and restriction respect the indicated homeomorphism topologies: a compact core together with one smaller annular neighbourhood of each end controls the image and inverse image of the complementary discs. This provides uniform control on the compactification from compact-open control on the punctured surface.
\end{proof}

\statement{A.2}{Known}{Surface groups with peripheral structure}
Choose an orientation, a base point $b\in U_{g,n}$, and paths to small positively oriented puncture circles. Then
\[
\Pi_{g,n}=\left\langle a_1,b_1,\ldots,a_g,b_g,c_1,\ldots,c_n\ \middle|\
 \prod_{i=1}^g[a_i,b_i]\prod_{j=1}^nc_j=1\right\rangle.
\]
For $n>0$ this is free of rank $2g+n-1$. For hyperbolic type it is centre-free and residually finite. Its peripheral datum consists of the conjugacy classes of the cyclic subgroups $\langle c_j\rangle$, with the coherent orientation datum; for $n=0$ we retain the orientation line in $H_2(\Pi_{g,0},\mathbb Z)$. The Dehn--Nielsen--Baer identification is
\[
 M^{\pm}_{g,n}\xrightarrow{\sim}\Out_{\mathrm{lac}}(\Pi_{g,n}).
\]
The target preserves the peripheral classes as a set; restricting to fixed classes and multiplier $+1$ gives $P_{g,n}$.
\Man{19}{111--113}{221--226}; see \citet[Theorems 8.1 and 8.8]{AF_FarbMargalit2012}.
\begin{proof}[Verification and scope]
The presentation follows by cutting into a polygon and applying van Kampen. Eliminating $c_n$ proves the free-rank assertion. For a closed hyperbolic surface the universal covering is the disc and the deck group is torsion-free cocompact Fuchsian; its centre is trivial. Free groups of rank at least two have trivial centre by reduced words. Residual finiteness follows from the faithful finitely generated characteristic-zero linear realizations and reduction modulo finite residue fields. The identification with peripheral-preserving outer automorphisms is exactly the punctured or closed Dehn--Nielsen--Baer theorem cited above; this is an imported theorem, not an assertion that an arbitrary group automorphism preserves omitted punctures. Changing $b$ or paths conjugates the resulting fundamental-group maps. Completing yields $\Pi_{g,n}\hookrightarrow\pi_{g,n}$; it does not yet prove faithfulness of $\widehat{M^+_{g,n}}\to\Out(\pi_{g,n})$.
\end{proof}

\statement{A.3}{Known}{Filling one point and the precise Birman group}
For a hyperbolic target $U_{g,n}$ and one extra point $b$, let $M^{\pm}_{g,n+1,*}$ be the subgroup fixing that extra point, while allowing the old punctures to be permuted. There are exact sequences of discrete groups
\[
1\longrightarrow\Pi_{g,n}\xrightarrow{\mathrm{push}}
 M^{\pm}_{g,n+1,*}\xrightarrow{\mathrm{forget}}M^{\pm}_{g,n}\longrightarrow1,
\qquad
1\longrightarrow\Pi_{g,n}\longrightarrow P_{g,n+1}\longrightarrow P_{g,n}\longrightarrow1.
\]
There is no forgetting homomorphism on the full group that permits the extra point itself to be permuted. On fundamental groups the inclusion of opens instead gives the surjection
$\Pi_{g,n+1}\to\Pi_{g,n}$ killing the normal closure of the last peripheral loop. These are different maps between different groups.
\Man{19}{103--111}{206--221}; \citet[Theorem 4.6]{AF_FarbMargalit2012}.
\begin{proof}[Source and comparison maps]
Evaluation at $b$ is a locally trivial fibration of homeomorphism groups, by isotopy extension. Its homotopy sequence, and contractibility of the identity component in hyperbolic type, give the first exact sequence. The connecting map pushes $b$ along a based loop. Its conjugation action on $\Pi_{g,n}$ is the point-pushing action and its injectivity uses the centre-free hyperbolic hypothesis. Restricting to homeomorphisms fixing the old punctures gives the second sequence. Surjectivity follows by isotoping the image of $b$ back to $b$ through an arc missing the punctures. For the surface-group map, the filled disc attaches along the last boundary circle, so van Kampen gives exactly the indicated normal quotient. Section C distinguishes their profinite analogues and their kernels.
\end{proof}

\statement{A.4}{Proved}{A covering gives a correspondence with specified lifts}
Let $p:U'\to U$ be a connected finite unramified covering of orientable finite-type surfaces, with $U$ hyperbolic. A choice of points above $b$ identifies $\pi_1(U')$ with a finite-index subgroup $H<\Pi=\pi_1(U)$. Put $L_H=\Stab_{M^{\pm}(U)}([H])$, where $[H]$ denotes conjugacy in $\Pi$. The group of isotopy classes of pairs of commuting lifts $(f',f)$ fits into
\[
1\longrightarrow N_\Pi(H)/H\longrightarrow\mathrm{Lift}(p)
 \longrightarrow L_H\longrightarrow1,
\qquad \mathrm{Lift}(p)\longrightarrow M^{\pm}(U').
\]
Thus the general covering operation in the tower is a correspondence, not a canonical homomorphism from all of $M^{\pm}(U)$ to $M^{\pm}(U')$. It preserves the sign of orientation and the induced peripheral subgroups. The degree $d=[\Pi:H]$ satisfies
$2g'-2+n'=d(2g-2+n)$.
\Man{19 and 28}{102--111 and 177--180}{204--221 and 350--356}.
\begin{proof}
A base homeomorphism lifts if and only if its outer action preserves $[H]$, by the covering-space lifting criterion. Any two lifts differ by a deck transformation, and the deck group is $N_\Pi(H)/H$. Lifting an isotopy gives the asserted exact sequence on isotopy classes. A deck transformation isotopic to the identity on a hyperbolic surface is trivial: it has finite order, whereas a nontrivial periodic orientation-preserving mapping class cannot be represented by an identity-isotopic map; equivalently use its faithful action on the hyperbolic universal cover and the centre-free surface group. This also verifies that the kernel has not been accidentally collapsed. Forgetting the lower homeomorphism gives the displayed map to the upstairs mapping class group. The restriction of a representative automorphism to $H$ changes, when the lift changes, by the deck action. Even for a characteristic subgroup this is not generally an inner automorphism of $H$. The correspondence records that ambiguity. Peripheral annuli lift to peripheral annuli, and Euler characteristic multiplies under a finite covering; both assertions follow by lifting a finite cell decomposition. The same subgroup inclusion completes to an open embedding $\widehat H\hookrightarrow\widehat\Pi$, since a finite-index subgroup inherits the full profinite topology.
\end{proof}

\subsection{A.2. Teichm\"uller space and the modular topos}

\statement{A.5}{Known}{The modular object retains isotropy}
For a stable labelled type, the analytic moduli stack is
$[\mathcal T_{g,n}/P_{g,n}]$, with $\mathcal T_{g,n}$ contractible; the unordered stack is $[\mathcal T_{g,n}/M^+_{g,n}]$. Consequently its geometric \emph{stack} fundamental group is respectively $\widehat P_{g,n}$ or $\widehat{M^+_{g,n}}$. The manuscript's $M_{g,n}$ in (19)--(25), scans 326--329, corresponds to the unordered version: (20) explicitly has $\widehat{\mathcal T}^+_{g,n}$, not $\widehat{\mathcal T}^{!+}_{g,n}$. In particular
$\mathscr M_{0,3}=\Spec\mathbb Q$ but $\mathscr M_{0,[3]}\simeq B\mathfrak S_3$.
\Man{21 and 26}{128--131 and 165--166}{255--261 and 326--329}.
\begin{proof}[Verification]
The complex structure on an oriented surface gives the usual Teichm\"uller parameter; quotient by the identity component of the diffeomorphism group is the manuscript's $E_g/A^\circ_{g,n}$. Teichm\"uller contractibility, proved by Fenchel--Nielsen coordinates in \citet[Theorem 10.6 and \S10.6.3]{AF_FarbMargalit2012}, and the mapping class action give the stated quotient stack, retaining automorphism groups at its points. Analytic Riemann existence for these Deligne--Mumford stacks identifies finite covering spaces with finite \etale covers, hence gives their completed fundamental groups. One can prove this stack version by applying ordinary Riemann existence to an \etale atlas and its twofold and threefold intersections, then descending finite covers. This is the stack extension of the comparison used in \citet[\S13.9]{AF_BGH2020}. A labelled three-pointed projective line has one isomorphism class and no automorphisms; allowing permutations gives precisely $\mathfrak S_3$. A coarse-space quotient would erase this group, so it cannot satisfy the manuscript's sequence (20). The convention in STRATA F.16 is correct only with this unordered/labelled distinction retained explicitly.
\end{proof}

\statement{A.6}{Known}{Arithmetic fundamental-group and universal-monodromy sequences}
Fix $\overline{\mathbb Q}\hookrightarrow\mathbb C$ and compatible geometric base points on the unordered moduli stack and its universal punctured curve $\mathscr U$. For hyperbolic type,
\[
1\longrightarrow\widehat{M^+_{g,n}}\longrightarrow
 A^{\mathrm{arith}}_{g,n}:=\pi_1^{\et}(\mathscr M_{g,[n]})
 \longrightarrow G_{\mathbb Q}\longrightarrow1,
\]
\[
1\longrightarrow\pi_{g,n}\longrightarrow\pi_1^{\et}(\mathscr U)
 \longrightarrow A^{\mathrm{arith}}_{g,n}\longrightarrow1.
\]
Conjugation on the normal geometric fibre defines a commutative diagram to
\[
1\longrightarrow\pi_{g,n}\longrightarrow
 \Aut_{\mathrm{lac}}(\pi_{g,n})\longrightarrow O_{g,n}\longrightarrow1.
\]
The last sequence uses $Z(\pi_{g,n})=1$. The induced map
$\rho_{g,n}:A^{\mathrm{arith}}_{g,n}\to O_{g,n}$ is genuine after these choices; its change under base-point transport is conjugation in $O_{g,n}$. Its geometric restriction need not be identified with an injection before the congruence problem is settled.
\Man{26}{165--166}{326--329}.
\begin{proof}[Source and exactness]
The first sequence is the homotopy sequence of a geometrically connected stack over $\mathbb Q$. For the second sequence, the universal punctured curve is smooth with a smooth proper compactification and relative normal-crossings puncture divisor; its geometric fibres have the fixed hyperbolic type. The homotopy sequence is exact and the geometric fibre embeds; analytically it is the Birman sequence, and the profinite fibre injection is the assertion justified in Section C. Passing to an \etale atlas and descent gives the same sequence for the unordered stack. Normality makes conjugation a homomorphism into automorphisms, and centre-freeness identifies inner automorphisms with the fibre. Images of peripheral loops are determined by the puncture divisor, with the cyclotomic multiplier; this proves that the maps land in the indicated groups. Restricting to the algebraic closure recovers the completed topological monodromy. The labelled sequences are displayed explicitly in \citet[Corollary 6.4]{HoshiMochizuki2011}; the finite permutation quotient gives the unordered versions just proved. These are precisely (19)--(25) with all base points and orientation conventions stated.
\end{proof}

\subsection{A.3. Completion, discretification, and arithmetization}

\statement{A.7}{Proved}{A completion is not its image}
The discrete action induces a continuous surjection
\[
\widehat{M^+_{g,n}}\twoheadrightarrow
\Sigma_{g,n}:=\overline{\rho_{g,n}(M^+_{g,n})}\subset O_{g,n}.
\]
The kernel is the congruence kernel. The double hat in the manuscript denotes the ambient automorphism group $O_{g,n}$, not the profinite completion of $\widehat{M^+_{g,n}}$. In all formulae for stabilizers or normalizers inside $O_{g,n}$, it is the image $\Sigma$ that can be used without the congruence hypothesis.
\Man{26}{162r--163v}{320--323}.
\begin{proof}
The group $\pi_{g,n}$ is topologically finitely generated. Characteristic open subgroups are cofinal among its open normal subgroups, so $\Aut(\pi_{g,n})$ and $\Out(\pi_{g,n})$ have their usual profinite topologies. The discrete mapping class action is continuous for the topology pulled back from the finite outer automorphism groups, hence extends from its full profinite completion. Compactness makes its image closed, and density identifies that image with the closure of the discrete action. The universal property alone says nothing about the kernel. The manuscript explicitly writes on scan 320 that the natural surjection is being assumed injective. Section B supplies the exact known ranges and implications.
\end{proof}

\statement{A.8}{Proved}{The quotient sets behind the terminology}
Let $\pi$ be a profinite group with the loop structure of $\widehat\Pi_{g,n}$. Write
$I(\pi)=\Isom_{\mathrm{lac}}(\widehat\Pi_{g,n},\pi)$ and
$I^{\mathrm{out}}(\pi)=I(\pi)/\Inn(\pi)$.
Actual discretifications are the images of $\Pi_{g,n}$ under elements of $I(\pi)$, with their transported loop structures. Their based parameter set is
\[
\mathrm{Discr}(\pi)=I(\pi)/\Aut_{\mathrm{lac}}(\Pi_{g,n}).
\]
After taking inner-conjugacy classes, the denominator is $M^{\pm}_{g,n}$. Pre-discretifications replace this discrete denominator by its closure. If
$\Sigma^{\pm}=\overline{M^{\pm}_{g,n}}\subset O_{g,n}$, the outer quotient sets are
\[
\mathrm{Pre}(\pi)=I^{\mathrm{out}}(\pi)/\Sigma^{\pm},
\qquad
\mathrm{Pre}^{+}(\pi)=I^{\mathrm{out}}(\pi)/\Sigma.
\]
The maps from discretifications to pre-discretifications are quotient maps, not identifications of dense subgroups. Oriented pre-discretifications map two-to-one to unoriented ones for hyperbolic type.
\Man{26}{162--163}{320--323}; \S29, folios 188--189, scans 371--374.
\begin{proof}
Two identifications give the same embedded discrete group and transported structure exactly when their quotient restricts to an automorphism of the standard discrete group. An automorphism of the discrete group extends uniquely to its completion, which proves the first orbit description. Taking target inner conjugacy gives the outer description. The closure of the based discrete automorphism group contains $\Inn(\pi)$, because the discrete surface group is dense in $\pi$; thus the pre-discretification quotient is already outer. The sign character on the discrete mapping class group has image $\{\pm1\}$, and remains so on its closure. Its positive subgroup is closed of index two. This proves the last assertion without imposing the congruence property. It also explains the harmless change of notation from classes of discretifications in \S26 to pre-discretifications in \S29.
\end{proof}

\statement{A.9}{Proved}{Two different groups called $N$, and two different conjectures}
For a chosen geometric subgroup $\Sigma\subset O$, define
\[
N^{\mathrm{norm}}=N_O(\Sigma),\qquad
\mathfrak G^{\mathrm{norm}}=N^{\mathrm{norm}}/\Sigma.
\]
These are the \S26 definitions. In \S29 the manuscript defines instead
\[
N^{\mathrm{arith}}=\rho(A^{\mathrm{arith}}),\qquad
\mathfrak G^{\mathrm{arith}}=N^{\mathrm{arith}}/\Sigma.
\]
There is an inclusion $N^{\mathrm{arith}}\subseteq N^{\mathrm{norm}}$, and a canonical continuous surjection
$G_{\mathbb Q}\twoheadrightarrow\mathfrak G^{\mathrm{arith}}$ once the algebraic closure and the arithmetic marking have been fixed. The statement that this surjection is injective does not assert that $N^{\mathrm{arith}}=N^{\mathrm{norm}}$.
\Man{26 and 29}{164--165 and 188}{324--326 and 370--371}.
\begin{proof}
The geometric subgroup of the arithmetic fundamental group is normal. Taking its image under $\rho$ proves the inclusion in the normalizer. The image is compact and therefore closed. Quotienting the surjective map from $A^{\mathrm{arith}}$ to its image by the geometric subgroup gives the displayed surjection. No element outside the image is produced by this quotient operation. Hence replacing the arithmetic image by the full normalizer is an additional assertion, and changes an injectivity problem into a larger realization problem. Sections D and E settle the arithmetic injectivity and retain the other question separately.
\end{proof}

\noindent The dictionary now has an economical form:
\begin{center}\small
\begin{tabularx}{\textwidth}{@{}L{.27\textwidth}L{.27\textwidth}Y@{}}
\toprule Manuscript object & Precise modern object & Qualification\\\midrule
$\widehat{\mathcal T}^+$ & $\widehat{M^+}$ or its image $\Sigma$ & These agree exactly when the relevant congruence kernel vanishes.\\
$\widehat{\widehat{\mathcal T}}$ & $\Out_{\mathrm{lac}}(\pi)$ & Full continuous loop-preserving outer automorphisms.\\
$\operatorname{Discr\acute et}$ & Dense discrete surface models & Actual models and their inner-conjugacy classes are distinguished.\\
$\operatorname{Pr\acute ediscr\acute et}^+$ & $I^{\mathrm{out}}/\Sigma$ & A homogeneous space, not a group without additional normality.\\
Strict geometric structure & Conjugate of $\Sigma$ & Orbit parametrized by $O/N^{\mathrm{norm}}$.\\
Arithmetization $a$ & Class in $I^{\mathrm{out}}/N^{\mathrm{arith}}$ & Its oriented pre-discretifications form a $\mathfrak G^{\mathrm{arith}}$-torsor.\\
Arithmetic Teichm\"uller group & $\pi_1^{\et}(\mathscr M_{g,[n]})$ & Its image $N^{\mathrm{arith}}$ can be a congruence quotient.\\\bottomrule
\end{tabularx}
\end{center}

\statement{A.10}{Proved}{The arithmetization fibre is a torsor}
The map $I^{\mathrm{out}}/\Sigma\to I^{\mathrm{out}}/N^{\mathrm{arith}}$ has fibres that are torsors for $N^{\mathrm{arith}}/\Sigma$. This assertion needs normality of $\Sigma$ in the arithmetic image, already proved, but not normality in all of $O$.
\Man{29}{189--190}{373--376}.
\begin{proof}
For a representative $f$ of an arithmetization, its fibre consists of the right cosets $fn\Sigma$, for $n\in N^{\mathrm{arith}}$. Multiplication identifies it with a simply transitive $N^{\mathrm{arith}}/\Sigma$-set. Replacing $f$ changes this identification by translation, which is precisely why the fibre is canonically a torsor and has no canonical origin. After the faithfulness theorem in D, the acting group is $G_{\mathbb Q}$, subject to the fixed arithmetic identification.
\end{proof}

\subsection{A.4. Exceptional types and editorial limits}

\statement{A.11}{Proved}{The non-hyperbolic types do not satisfy the hyperbolic assertions}
The following table refers to punctures, not boundary circles fixed pointwise:
\begin{center}\small
\begin{tabular}{@{}llll@{}}\toprule
Type & $\Pi$ & $M^+$ & $P$\\\midrule
$(0,0)$ & $1$ & $1$ & $1$\\
$(0,1)$ & $1$ & $1$ & $1$\\
$(0,2)$ & $\mathbb Z$ & $\mathbb Z/2$ & $1$\\
$(0,3)$ & $F_2$ & $\mathfrak S_3$ & $1$\\
$(1,0)$ & $\mathbb Z^2$ & $\mathrm{SL}_2(\mathbb Z)$ & $\mathrm{SL}_2(\mathbb Z)$\\
$(1,1)$ & $F_2$ & $\mathrm{SL}_2(\mathbb Z)$ & $\mathrm{SL}_2(\mathbb Z)$\\\bottomrule
\end{tabular}
\end{center}
The extended groups are obtained by including degree $-1$ maps; for $(0,2)$ the result is $(\mathbb Z/2)^2$, and for both torus types it is $\mathrm{GL}_2(\mathbb Z)$. The surface group $F_2$ in the last line has peripheral element $[a,b]^{-1}$; the $F_2$ in type $(0,3)$ has three peripheral classes. They are not the same group with loops. Forgetting the marked point on a torus gives an isomorphism of mapping class groups, not a Birman sequence with an injected $\mathbb Z^2$ kernel.
\Man{19}{111--115}{222--230}; \S21, folios 129--131, scans 257--261.
\begin{proof}
A sphere with three ordered points has trivial orientation-preserving mapping class group, by cutting into discs; allowing their permutations gives $\mathfrak S_3$. Forgetting points gives the first three entries, with the remaining transposition for two unordered points. Complex conjugation fixes the real punctures and supplies the orientation-reversing involution; for two punctures it is independent of the permutation. A torus homeomorphism acts on its first homology by an integral matrix, and the linear maps realize every such matrix. Cutting along a basis proves that the action on homology detects the mapping class. A point can be fixed by translation, and translations are isotopic to the identity, so the pointed and unpointed groups coincide. For the cylinder, write $\rho_{\rm perm}$ for the permutation sign and $\rho_{\pi}$ for the sign on $\pi_1\simeq\Z$. Then $\rho_{\pi}=\rho_{\rm perm}\epsilon$. Indeed $z\mapsto z^{-1}$ has permutation sign $-1$ and orientation sign $+1$, while $z\mapsto\overline z$ has signs $+1,-1$; both induce $-1$ on $\pi_1$. Thus the identification of the two homomorphisms called $\rho$ on scan 225 is false if read literally, and the determinant-one image assertion for the rank-one case on scan 227 is false as well. The orientation-preserving cylinder group acts on $\Z$ by both signs. These computations also exhibit why the fundamental group of the identity component, which is nontrivial in the torus case, destroys the hyperbolic point-push injection. Their finite permutation and matrix checks are reproduced in \texttt{code/af\_checks.py}. For $(0,0)$ and $(0,1)$ the group alone has no orientation information; the manuscript adds an orientation line and the type, as stated explicitly at scans 229--230.
\end{proof}

\statement{A.12}{Proved}{Fields, real structures, and characteristics are part of the data}
The complex Teichm\"uller tower in this paper consists of orientable surfaces. A non-orientable surface is not the underlying surface of a smooth complex algebraic curve. For a real curve the relevant structure is an extension by $G_{\mathbb R}\simeq\mathbb Z/2$, together with its outer involution on the geometric group; it is not obtained by replacing the curve with the topological quotient by complex conjugation. All arithmetic comparison statements here are in characteristic zero. At primes $2$ and $3$, prime-to-characteristic or tame specialization cannot silently replace the full profinite group, and stack isotropy of orders $2$ and $3$ must be retained.
\begin{proof}
Holomorphic coordinate changes preserve the real orientation, proving the first assertion. The arithmetic fundamental-group sequence over $\mathbb R$ has geometric kernel and quotient $G_{\mathbb R}$; conjugation reverses the complex orientation and has cyclotomic value $-1$. Its quotient surface may have boundary or be non-orientable and is a different topological object. In characteristic $p$, wild inertia is pro-$p$ and is killed by prime-to-$p$ completion. Moreover elliptic curves in characteristic zero already have stabilizers of orders $2$, $4$, and $6$ at suitable points. Thus erasing all $2$- or $3$-primary isotropy cannot preserve the stack group. The $p$-adic statements cited in D have their own all-prime hypotheses; no comparison made here transports them to a curve over a finite field. These are scope restrictions, not an unproved assertion of a non-orientable analogue.
\end{proof}

\noindent\textbf{Reading record.}
On scan 204 the sequence of marked points is written with a positive-index convention, while (1) and the margin use indices beginning with zero. We use an arbitrary nested enumeration: the mathematics supports a reindexing, not a settled palaeographic correction. On scan 320 the assumption of injectivity is visible in the image and must not be suppressed. On scan 328 the transcription's folio alternatives $166|168$ do not alter any argument; the numerical sequence favours 166, without settling the reading. The phrase ``classes of discretifications'' at scan 321 is explicitly broadened by the insertion about pre-discretifications. At scan 261 the transcribed chain with changing genus cannot be literally correct by dimension: the intended zero-genus chain $\mathcal T_{0,3}\to\mathcal T_{0,2}\to\mathcal T_{0,1}\to\mathcal T_{0,0}$ fits the following sentence. That scan image was not supplied, so we record mathematical support only, not a visual emendation.


\subsection{A.5. The universal lift in \S21}

\statement{A.13}{Proved}{The three lift extensions and the mixed fundamental group}
Assume hyperbolic type and choose the additional base point $b$. Write $A=A^+_{g,n}$, $A_0$ for its identity component, $M=\pi_0A$, and $M_*=\pi_0(A^+_{g,n+1,*})$. The topological pullback
\[
 \widetilde A=A\times_M M_*
\]
has the three exact sequences
\[
\begin{split}
1&\longrightarrow\Pi_{g,n}\longrightarrow\widetilde A
                 \longrightarrow A\longrightarrow1,\\
1&\longrightarrow A_0\longrightarrow\widetilde A
                 \longrightarrow M_*\longrightarrow1,\\
1&\longrightarrow A_0\times\Pi_{g,n}\longrightarrow\widetilde A
                 \longrightarrow M\longrightarrow1.
\end{split}                                                   \tag{A.1}
\]
The first extension splits canonically over the subgroup $A_*$ fixing $b$, by lifting maps to fix a chosen point of the universal cover. If $E=EA$ is a universal principal $A$-space (for example the Milnor construction), the associated bundle
$\mathcal X=(E\times U_{g,n})/A_0$ has base $E/A_0$ contractible up to homotopy and fibre $U_{g,n}$. Its residual $M$-action has mixed fundamental group
\[
 \pi_1(\mathcal X\mathbin{/\mkern-6mu/}M)\simeq M_*,
 \qquad
 1\to\pi_1(U_{g,n})\to M_*\to M\to1.                         \tag{A.2}
\]
Here the double slash is the homotopy quotient. These identifications use compatible base points; changing them gives inner conjugation. The same pullback construction applies to the extended groups.
\Man{21}{126--132}{251--263}, especially (74)--(75), scans 252--254.
\begin{proof}
A.3 gives $1\to\Pi\to M_*\to M\to1$. Pull it back along $A\to M$. The first two kernels in (A.1) follow immediately. Over the identity of $M$ the pullback is the \emph{direct} product $A_0\times\Pi$: its multiplication is componentwise. This proves the third sequence and supports the reading $\times$, rather than a nontrivial semidirect product, in (75). Every $f\in A_*$ has a unique lift fixing the chosen lift of $b$; composition of such lifts has the same property. This gives the canonical splitting over $A_*$.

Hyperbolicity implies contractibility of $A_0$, as in the proof of the Birman theorem. Thus $E/A_0$ has the homotopy type of $BA_0$ and is contractible. The associated bundle has fibre $U$, so its fibre inclusion is a homotopy equivalence. Quotienting the residual action gives
$\mathcal X/\!/M\simeq E\times_AU$. Evaluation identifies the homogeneous space $U$ with $A/A_*$; local sections are supplied by isotopy extension. Therefore $E\times_AU\simeq E/A_*\simeq BA_*$. Its fundamental group is $\pi_0A_*=M_*$. The fibration over $BM$ gives exactly (A.2), including the point-pushing kernel. This is the mixed-group construction sketched in \S21, not a genuine action of $M$ on a single marked surface chosen without additional data.
\end{proof}


\clearpage\section{Congruence, normalizers, and discretizations}\label{part:B}

We use the conventions of Part A.  Thus $P_{g,n}$ is the pure,
orientation-preserving mapping class group, $M^+_{g,n}$ allows permutations
of the punctures, and $M^{\pm}_{g,n}$ also allows reversal of orientation.
We write $\Pi_{g,n}$ for the discrete surface group and
$\pi_{g,n}=\widehat{\Pi}_{g,n}$ for its full profinite completion.  Unless
an exception is explicitly stated, $2g-2+n>0$.  All automorphisms of
profinite groups are continuous.  In
\[
 O_{g,n}=\operatorname{Out}_{\mathrm{lac}}(\pi_{g,n})
\]
the peripheral conjugacy classes may be permuted and their generators
may be raised to a common profinite unit.  For a proper surface the
corresponding multiplier is its action on the orientation class.  We set
\[
 \check P_{g,n}=\overline{P_{g,n}},\qquad
 \Sigma_{g,n}=\overline{M^+_{g,n}},\qquad
 H_{g,n}=\overline{M^{\pm}_{g,n}},
 \qquad N^{\mathrm{norm}}_{g,n}=N_{O_{g,n}}(\Sigma_{g,n}).
\]
These are images or closures in $O_{g,n}$, not, by definition, abstract
profinite completions.  This distinction is indispensable in reading
\S26, folios 162--165, scans 320--326.  The arithmetic image
$N^{\mathrm{arith}}$ introduced in \S29, folio 188, scan 371 is a
different group; its behavior under filling is treated in Part C.

\subsection{The congruence problem and the size of the ambient group}

\statement{B.1}{Proved}{The precise congruence question}
There is a canonical continuous epimorphism
\[
 c_{g,n}:\widehat P_{g,n}\longrightarrow\check P_{g,n},
 \qquad C_{g,n}=\ker(c_{g,n}).
\]
The following assertions are equivalent: $C_{g,n}=1$; the profinite and
congruence topologies on $P_{g,n}$ coincide; and every subgroup of finite
index in $P_{g,n}$ contains the kernel of
$P_{g,n}\to\operatorname{Out}(\Pi_{g,n}/K)$ for some characteristic
subgroup $K$ of finite index in $\Pi_{g,n}$.  None of these assertions
asserts normality of $\check P_{g,n}$, $\Sigma_{g,n}$, or $H_{g,n}$ in
$O_{g,n}$.

\begin{proof}
The kernels in the statement form a neighborhood basis for the induced
congruence topology.  Indeed, a finitely generated group has only finitely
many subgroups of a prescribed finite index; intersecting all subgroups
of index at most a prescribed integer gives cofinal characteristic
subgroups.  The completion for this topology is exactly the closure of
the image in $\operatorname{Out}(\pi_{g,n})$, since the latter has the
topology determined by its characteristic finite quotients.  The
universal property of profinite completion gives $c_{g,n}$, whose image
is closed and contains a dense subgroup, hence is the whole closure.
If its kernel is trivial, compactness makes it an isomorphism of
topological groups; pulling back open subgroups gives cofinality of the
congruence neighborhoods.  The converse follows from the same inverse
limit description.  A kernel calculation for a map into an ambient
group makes no assertion about conjugation by elements outside its
image.  Even for finite groups, an injective map $C_2\hookrightarrow
\mathfrak S_3$ need not have normal image.  Thus the replacement of all
the normality questions in scans 322--326 by the congruence problem, as
suggested in Lead 4, is invalid.
\end{proof}

\statement{B.2}{Known}{The established congruence cases}
For full profinite completion, $C_{g,n}=1$ when $g=0$, $n\geq3$; when
$g=1$, $n\geq1$; and when $g=2$, $n\geq0$.  The same faithfulness
holds for the finite extensions allowing permutations and orientation
reversal.  In particular, all genus-zero groups used in Part F have
their full profinite topology induced by their surface-group action.

The first two assertions are Asada's results, in the precise forms
recorded and proved in \citet[Proposition 2.6 and Theorem 2.8]{Boggi2009BC}.
The same attached paper proves the hyperelliptic result for punctured
surfaces, hence the genus-two result for $n\geq1$, in its Theorem 3.5.
The closed genus-two case requires the subsequent
\citet[Theorem 1.4(ii), with $\mathscr C=\mathscr F$]{Boggi2020BC}:
the hyperelliptic mapping class group equals the full mapping class
group in genus two.  The distinction matters: the 2009 paper expressly
restricts its main theorem to open surfaces and records a gap in an
earlier proof.  The more recent primary source
\citet[\S1.1 and Corollary B]{Boggi2026BC} confirms the genus-at-most-two
range and establishes additional centralizer cases, without solving
the problem in arbitrary genus.

To justify the finite-extension assertion, the pure orientation subgroup
is congruence-open: permutation of the finitely many peripheral classes
is continuous, and orientation reversal is distinguished from
orientation preservation by the multiplier modulo $3$.  Within the
geometric image the multiplier takes only the values $1$ and $-1$.
The kernel of the completed action on the finite extension therefore
lies in $\widehat P_{g,n}$, where it is trivial.  This is not the ordinary
congruence-subgroup property of $\operatorname{SL}_2(\mathbb Z)$ acting
on $\mathbb Z^2$.  For the punctured torus the decisive representation
acts on $\widehat F_2$ and remembers more than homology.

\statement{B.3}{Open}{The remaining full congruence problem}
For $g\geq3$, determine whether
\[
 \widehat P_{g,n}\longrightarrow
 \operatorname{Out}(\pi_{g,n})
\]
is injective.  The smallest genus and number of punctures not covered
by B.2 is $(g,n)=(3,0)$.  Part C proves
$C_{g,n+1}\simeq C_{g,n}$ whenever the target type is hyperbolic, so
adding punctures does not remove this unresolved case.  The exact
missing step is cofinality of characteristic-quotient levels among
all finite-index subgroups; neither the normality of a geometric image
in an arithmetic image nor faithfulness of the Galois action supplies
it.  This is a literature-open problem, not merely a gap in the
manuscript.  Our status check includes the 2026 primary source cited in
B.2.

\statement{B.4}{Proved}{Properness of the geometric image and the exceptional types}
For every hyperbolic type,
\[
 H_{g,n}\ \ne\ O_{g,n}.
\]
The assertion ``never an isomorphism'' in \S26, folio 163, scan 322
is correct in this range.  Read as an assertion about bare
fundamental groups outside that range, it is false; adding an
independent orientation line changes the data and must be specified.

\begin{proof}
The continuous multiplier
$\lambda:O_{g,n}\to\widehat{\mathbb Z}^{\times}$ takes values in
$\{1,-1\}$ on $H_{g,n}$.  Choose a smooth marked curve of the
specified type over some number field $L$; all required markings can
be made rational after a finite extension.  The arithmetic outer
action has multiplier the cyclotomic character of $G_L$.  Its image
is an open, hence infinite, subgroup of
$\widehat{\mathbb Z}^{\times}$, since
$[L\cap\mathbb Q(\mu_\infty):\mathbb Q]<\infty$.  Consequently an
arithmetic element with multiplier outside $\{1,-1\}$ lies in
$O_{g,n}\setminus H_{g,n}$.  No faithfulness theorem and no congruence
conjecture are needed for this argument.

For $(0,0)$ and $(0,1)$ the geometric fundamental group is trivial,
so its automorphism group and every image in it are trivial.  For
$(0,2)$ it is $\widehat{\mathbb Z}$, and the geometric image is
$\{1,-1\}\subset\widehat{\mathbb Z}^{\times}$; forgetting the two
labels does not change this image.  For $(1,0)$ it is
$\widehat{\mathbb Z}^2$, and the ambient group is
$\operatorname{GL}_2(\widehat{\mathbb Z})$.  The extended geometric
image is the subgroup with determinant in $\{1,-1\}$, whereas the
pure oriented image is $\operatorname{SL}_2(\widehat{\mathbb Z})$.
Here the action factors through the usual homological congruence
completion, and its kernel on the abstract profinite completion of
$\operatorname{SL}_2(\mathbb Z)$ is nontrivial.  Thus the punctured
torus and the unpunctured torus must not be merged.
The auxiliary orientation line that the manuscript adds in the
trivial-group cases, noted in A.11, is not contained in the bare
fundamental group; the preceding triviality statement does not
erase that extra datum.

For completeness, this last nontriviality has a finite witness.
The presentation $\operatorname{PSL}_2(\mathbb Z)=C_2*C_3$
maps onto $\operatorname{SL}_2(\mathbb F_8)$ by sending its
generators to
\[
 \begin{pmatrix}1&t\\0&1\end{pmatrix},\qquad
 \begin{pmatrix}0&1\\1&1\end{pmatrix},
 \quad\mathbb F_8=\mathbb F_2[t]/(t^3+t+1).
\]
The code verifies their orders, that they generate all 504
determinant-one matrices, and that the group is simple: its
conjugacy-class sizes are $1,56,56,56,56,63,72,72,72$, and no
proper union containing the identity has cardinality dividing
504.  A finite quotient of
$\operatorname{SL}_2(\widehat{\mathbb Z})$ has only cyclic
simple factors and groups $\operatorname{PSL}_2(\mathbb F_p)$
for primes $p\geq5$.  This follows by reduction modulo $p$
in each factor; the reduction kernel is pro-$p$, and the
groups for $p=2,3$ are solvable.  None of those projective
groups has order 504: the order formula gives 168 at $p=7$
and 660 at $p=11$, with no intervening prime.  The displayed
simple quotient therefore cannot be a congruence quotient,
which proves the claimed nontrivial kernel.
\end{proof}

\statement{B.5}{Refuted}{Normality of the full extended geometric group}
The assertion that $H_{g,n}$ is normal in $O_{g,n}$ for every
hyperbolic type is false already at $(0,3)$.

\begin{proof}
The three-punctured sphere has
$P_{0,3}=1$, $M^+_{0,3}=\mathfrak S_3$, and
$M^{\pm}_{0,3}=\mathfrak S_3\times C_2$.
For the rational curve
$\mathbb P^1\setminus\{0,1,\infty\}$ the outer Galois action is
faithful, as proved in D.16 from Belyi's theorem.  Its elements fix
the three peripheral labels.  The six rational fractional linear
transformations permuting $0,1,\infty$ commute with this action in the
outer automorphism group.  Hence
\[
 \mathfrak S_3\times G_{\mathbb Q}\hookrightarrow O_{0,3}.
\]
The intersection of the two factors is trivial because the first
factor acts faithfully on the labels and the second fixes them.
If $c\in G_{\mathbb Q}$ is complex conjugation, its image is the
orientation-reversing class fixing those labels.  Thus
$H_{0,3}=\mathfrak S_3\times\langle c\rangle$ inside this subgroup.
Normality of $H_{0,3}$ would imply normality of
$\langle c\rangle$ in $G_{\mathbb Q}$.

The splitting field of $t^3-2$ has Galois group $\mathfrak S_3$:
Eisenstein's criterion gives irreducibility and its discriminant is
$-108$, which is not a square.  Complex conjugation acts as a
transposition on its three roots.  Conjugating this transposition by
a three-cycle gives another transposition outside its order-two
subgroup.  Lifting the three-cycle to $G_{\mathbb Q}$ contradicts
normality.  The finite calculation, including the discriminant and
the conjugation, is reproduced by \texttt{code/bc\_counterexamples.py}.
\end{proof}

This answers the normality question attached to the extended group in
scan 322.  It does not answer the different question about
$\Sigma_{g,n}=\overline{M^+_{g,n}}$.  Nor does an author's annotation
``ultra-faux'' by itself provide a mathematical counterexample to a
different self-normalizer assertion: the conjecture so annotated is
recorded as withdrawn, and we do not silently substitute B.5 for it.

\statement{B.6}{Proved}{Conjugacy of geometric structures and their normalizers}
Fix an identification of a surface-group completion with a profinite
group with loops $\pi$, and let $\Sigma\subset O=\operatorname{Out}_{\mathrm{lac}}(\pi)$
be its oriented geometric image.  The conjugates of $\Sigma$ are
parameterized by $O/N_O(\Sigma)$.  The natural map
\[
 O/\Sigma\longrightarrow O/N_O(\Sigma)
\]
has fibers which are torsors under $N_O(\Sigma)/\Sigma$, after
transport by a representative.  In particular, uniqueness of this
conjugate subgroup means $N_O(\Sigma)=O$, equivalently normality of
$\Sigma$ in $O$.

\begin{proof}
Two conjugates $a\Sigma a^{-1}$ and $b\Sigma b^{-1}$ agree exactly
when $b^{-1}a$ normalizes $\Sigma$.  For a fixed representative $a$,
the fiber over $aN_O(\Sigma)$ is $aN_O(\Sigma)/\Sigma$.
This is the asserted torsor.  All subgroups used for the displayed
quotients are closed: the normalizer of a closed subgroup of a
profinite group is closed.  These are therefore profinite coset spaces,
although the first need not be a group.  This proves the assertions
in \S26, folios 164--165, scans 324--326, with the normalizer definition
kept explicit.
\end{proof}

\statement{B.7}{Open}{The unrestricted normalizer comparison}
The unrestricted questions
\[
 \Sigma_{g,n}\triangleleft O_{g,n},\qquad
 N^{\mathrm{norm}}_{g,n}=O_{g,n}
\]
are equivalent to each other by B.6 and are separate from B.3.
We do not resolve this unrestricted comparison.  At the smallest
test type $(0,3)$ it asks whether every inertia-preserving outer
automorphism normalizes the geometric $\mathfrak S_3$.
E.7 identifies the normalizer quotient using the additional
symmetry condition; that identification does not show that every
element of $O_{0,3}$ satisfies the condition.  Our searches of the
primary sources on quasi-special symmetric automorphisms and
procongruence automorphisms did not supply either this implication or
a counterexample.  Here the label Open records the precise unresolved
comparison in this audit; unlike B.3, we do not assert that each
individual case has been proved open by a literature classification.
The earlier self-normalizer conjecture for $H$, explicitly withdrawn
in scan 322, is not used as a premise.

\subsection{Actual discretizations under a finite covering}

A based discretization of $\pi$ will mean a dense subgroup
$D\subset\pi$, equipped with the oriented surface peripheral
structure, such that completion induces an isomorphism
$\widehat D\simeq\pi$ preserving that structure.  In particular it is
more than an arbitrary dense subgroup.  Forgetting orientation gives
the corresponding unoriented notion.  Quotienting by inner
automorphisms of $\pi$ gives strict prediscretizations; quotienting a
marking torsor by the closure of the geometric automorphism group
gives the manuscript's non-strict prediscretizations.  The adjective
``strict'' must not be dropped in these comparisons.

\statement{B.8}{Proved}{A surface lattice is self-normalizing in its completion}
If $D$ is a nonabelian free group or the fundamental group of a closed
orientable surface of genus at least two, then
\[
 N_{\widehat D}(D)=D.
\]

\begin{proof}
Free groups and closed orientable surface groups are conjugacy
separable and have Grossman's Property A: every automorphism taking
each element to a conjugate of itself is inner.  For the exact
conjugacy-separability statements we use
\citet[Theorems 3.2 and 4.1; Property A is stated in \S1 immediately
before Theorem 1.7]{Paris2009BC}.  These theorems are stronger, giving
conjugacy $p$-separability.  If $h\in\widehat D$ normalizes $D$,
conjugation by $h$ induces an automorphism $\alpha$ of $D$.  For every
$d\in D$, $\alpha(d)$ and $d$ are conjugate in $\widehat D$; hence
they are conjugate in $D$.  Property A gives
$\alpha=\operatorname{int}(d_0)$ for some $d_0\in D$.
Then $d_0^{-1}h$ centralizes the dense subgroup $D$ and therefore
centralizes $\widehat D$.  The latter is center-free.  For example,
this follows from the slimness theorem in D.14, or
from the centralizer calculation for two independent simple loops
in \citet[Proposition 3.6]{BoggiZalesskii2019BC}.
Consequently $h=d_0$.  The reverse inclusion is immediate.
\end{proof}

\statement{B.9}{Proved}{Injectivity of restriction to a finite cover}
Let $\pi$ have hyperbolic orientable surface type and let
$V\subset\pi$ be an open subgroup with its induced peripheral
structure.  The map
\[
 \operatorname{Discr}(\pi)\longrightarrow\operatorname{Discr}(V),
 \qquad D\longmapsto D\cap V
\]
is well defined and injective, for based discretizations, and the
same holds with orientation retained.  No congruence-subgroup
property is assumed.  This proves the corollary in \S28, folio 178,
scans 352--353.

\begin{proof}
The intersection $D\cap V$ is of finite index in $D$.  The topology
induced on a finite-index subgroup by the full profinite topology is
its full profinite topology: the core in $D$ of a subgroup of finite
index in $D\cap V$ is of finite index in $D$.  Therefore
$\widehat{D\cap V}\simeq V$.  The ordinary finite covering of
surfaces determines the asserted peripheral and orientation data.

First suppose that $V$ is normal.  If
$D_1\cap V=D_2\cap V=\Lambda$, both $D_1$ and $D_2$ normalize
$\Lambda$.  B.8 applied inside $V=\widehat\Lambda$ gives
\[
 N_\pi(\Lambda)\cap V=\Lambda.
\]
Each $D_i$ maps onto the finite group $\pi/V$, since it is dense.
For $a\in N_\pi(\Lambda)$ choose $d\in D_i$ with the same image
in $\pi/V$.  Then $ad^{-1}\in\Lambda\subset D_i$, proving
$N_\pi(\Lambda)=D_i$.  Thus $D_1=D_2$.
For arbitrary $V$, replace it by its core
$W=\bigcap_{a\in\pi}aVa^{-1}$.  This is open and normal, and
$D_1\cap W=D_2\cap W$ follows by intersecting the given equality
with $W$.  The normal case applies.  Orientation is unchanged under
these restrictions.
\end{proof}

\statement{B.10}{Proved}{The image criterion for a normal finite cover}
Suppose $V\triangleleft\pi$ is open, $G=\pi/V$, and
$\Lambda\subset V$ is an oriented discretization.  It extends to an
oriented discretization of $\pi$ if and only if the faithful deck
outer action
\[
 G\longrightarrow\operatorname{Out}_{\mathrm{lac}}(V)
\]
factors through the image of
$\operatorname{Out}_{\mathrm{lac}}^+(\Lambda)$ in
$\operatorname{Out}_{\mathrm{lac}}(V)$.  The extension, when it
exists, is unique as a subgroup of $\pi$.

\begin{proof}
Necessity is immediate by conjugating $D\cap V$ inside an extending
lattice $D$.  For sufficiency, the discrete outer action on $\Lambda$
is unique because the completion map on its outer automorphism group
is injective, by the same conjugacy-separability and Property A
argument as B.8.  The deck action is faithful: an element of $\pi$
acting on $V$ by an inner automorphism differs from an element of
$V$ by an element centralizing $V$, and $C_\pi(V)=1$.  The last
equality is the slimness of a hyperbolic profinite surface group.

Since $Z(\Lambda)=1$, the specified outer action determines a unique
extension $1\to\Lambda\to E\to G\to1$.  Explicitly, it is the
pullback of $\operatorname{Aut}(\Lambda)\to
\operatorname{Out}(\Lambda)$ along that action.  Completing
automorphisms maps it into the analogous pullback for $V$; the
latter is $\pi$, since $Z(V)=1$.  The resulting map $E\to\pi$ is
injective, induces the given inclusion on $\Lambda$, and is the
identity on $G$.  Its image is dense.  Because $\Lambda$ has finite
index in $E$, completion of this extension is exact, and the
comparison with $1\to V\to\pi\to G\to1$ gives
$\widehat E\simeq\pi$.

The profinite group $\pi$ is torsion-free.  A precise source is
\citet[Lemma 3.1(ii)]{BoggiZalesskii2019BC}, applied to a Fuchsian
surface group without orbifold points: there are no decomposition
groups in which a nontrivial finite subgroup could lie.
Thus $E$ is torsion-free.  The finite extension is allowable in the
Fuchsian sense: its outer action preserves the peripheral classes
and orientation.  \citet[Theorem 7, p.~258]{Kerckhoff1983BC} says
that an allowable finite extension of a Fuchsian group is Fuchsian
when its outer action is injective.  Hence $E$ is the fundamental
group of an orientable surface of finite type, and the covering
defined by $\Lambda$ is a genuine unramified surface covering.

It remains to check the loops, rather than merely the underlying
group.  Each peripheral cyclic subgroup of $E$ meets $\Lambda$ in
a finite-index peripheral subgroup.  Its closure is the centralizer
in $\pi$ of the corresponding peripheral subgroup of $V$.
Conversely, the original peripheral groups of $\pi$ are obtained
from those of $V$ by exactly this centralizer operation.
Both assertions follow from
\citet[Proposition 3.6]{BoggiZalesskii2019BC}: a nonzero power of
an algebraically simple element has normalizer, and hence
centralizer, equal to the maximal procyclic group it generates.
The two peripheral structures therefore agree.  The allowable
orientation agrees on the finite cover and hence on the base.
Uniqueness is B.9.
\end{proof}

In scan 352 the displayed quotient for the deck group is transcribed
as $\pi'/\pi$, although the surrounding text has $\pi'\subset\pi$.
The image inspection confirms this apparent slip.  B.10 requires
$G=\pi/\pi'$; it does not assert that the handwriting has thereby
been editorially settled.  The same page only sketches the
loop-compatibility step, supplied above using the centralizer theorem.

\statement{B.11}{Refuted}{Surjectivity for raw discretizations of a finite cover}
Restriction of based discretizations need not be surjective.  A
counterexample is the degree-two covering
\[
 \mathbb P^1(\mathbb C)\setminus\{0,\infty,1,-1\}
 \longrightarrow
 \mathbb P^1(\mathbb C)\setminus\{0,\infty,1\},
 \qquad z\longmapsto z^2.
\]

\begin{proof}
Let $V=\widehat\Lambda$ be the subgroup defined by this cover in
$\pi=\widehat F_2$, and let $s$ be the deck involution in
$\operatorname{Out}_{\mathrm{lac}}(V)$.  It fixes the punctures
$0,\infty$ and interchanges $1,-1$.  In the discrete pure mapping
class group $P_{0,4}\simeq F_2$, conjugation by $s$ is nontrivial:
it takes the twist about a curve separating $\{0,1\}$ from the
other two punctures to the distinct twist about a curve separating
$\{0,-1\}$.  The permutation of these two partitions is checked in
\texttt{code/bc\_counterexamples.py}.

The centralizer of $s$ in $P_{0,4}$ has infinite index.  Otherwise
the induced automorphism would fix a subgroup of finite index.
Each element of $F_2$ has a positive power in that subgroup, and
the unique-root property in a free group would force the
automorphism to fix every element, a contradiction.  Let
$B=\overline{P_{0,4}}\subset\operatorname{Out}_{\mathrm{lac}}(V)$.
Its closed subgroup $C_B(s)$ is not open, since openness would
make its intersection with the dense group $P_{0,4}$ of finite
index.  Thus $B/C_B(s)$ is uncountable: a compact Hausdorff space
cannot be a countable union of closed cosets with empty interior.

The discrete group $\operatorname{Out}_{\mathrm{lac}}^+(\Lambda)$
is countable.  We may therefore choose $u\in B$ for which
$u^{-1}su$ is not in that discrete group.  For definiteness this
choice can be specified by enumerating the discrete group and
choosing successively nested basic clopen neighborhoods which
avoid the closed conjugator set for the next element; require
also containment in the next characteristic finite-quotient
neighborhood.  Each avoidance is possible because the nonempty
conjugator sets are cosets of the nonopen centralizer.  Compactness
and the cofinal neighborhood condition determine an element $u$.
Choose a genuine automorphism representing $u$, and put
$\Lambda'=u(\Lambda)$.  The deck outer action does not factor
through $\operatorname{Out}_{\mathrm{lac}}^+(\Lambda')$.
By B.10, $\Lambda'$ does not extend to $\pi$.

The counterexample is necessarily profinite: all discretizations
are dense and their difference cannot be certified by distinguishing
their images in a single finite quotient.  The nested construction
above specifies the required element, rather than relying on a
purported finite-quotient witness.  Notice also that $u\in B$:
$\Lambda$ and $\Lambda'$ have the same non-strict oriented
prediscretization.  This counterexample does not disprove the
prediscretization assertion in scan 355.
\end{proof}

\subsection{Passing to prediscretizations and the remaining comparisons}

\statement{B.12}{Proved}{Restriction descends to both prediscretization quotients}
For every finite covering $V\subset\pi$ as in B.9, restriction of
discretizations descends to strict prediscretizations and to the
non-strict oriented and unoriented prediscretizations.  In
particular, the well-definedness doubts in \S28, folios 179--181,
scans 355--358 have an affirmative answer.

\begin{proof}
First let $D_2=aD_1a^{-1}$, $a\in\pi$.  Density gives
$\pi=VD_1$: every coset of the open subgroup $V$ meets $D_1$.
Write $a=vd$ with $v\in V$, $d\in D_1$.  Then
$D_2=vD_1v^{-1}$, and
$D_2\cap V=v(D_1\cap V)v^{-1}$.  This is precisely the correction
already made in scan 357.

For the closure quotient, fix $D=D_1$ and let $A_D$ be its discrete
loop-preserving automorphism group, with orientation restriction
if required, inside $\operatorname{Aut}_{\mathrm{lac}}(\pi)$.
Suppose $D_2=u(D)$ with $u\in\overline{A_D}$.  The group $\pi$ is
topologically finitely generated, so it has finitely many open
subgroups of index $[\pi:V]$.  The action on this finite set is
continuous.  Density of $A_D$ in its closure permits a choice
$a\in A_D$ with $a(V)=u^{-1}(V)$.  Thus $ua$ stabilizes $V$ and
$ua(D)=D_2$.

The intersection of $\overline{A_D}$ with the open stabilizer of
$V$ is the closure of $A_D\cap\operatorname{Stab}(V)$.
Restriction of automorphisms stabilizing $V$ to $V$ is continuous;
for instance, restriction is a continuous map between the closed
automorphism spaces in the spaces of maps on the corresponding
finitely generated profinite groups.  It takes the discrete
stabilizer into $\operatorname{Aut}_{\mathrm{lac}}(D\cap V)$.
Therefore the restrictions of $D$ and $u(D)$ differ by an element
in the closure of that latter group.  The same proof preserves
orientation, because orientation restricts compatibly along a
surface covering.  It proves descent, not bijectivity.
\end{proof}

\statement{B.13}{Proved}{The exact coset criteria behind conditions (a)--(e)}
Let $f:A\to A'$ be a continuous homomorphism of profinite groups,
and let $B\subset A$, $B'\subset A'$ be closed subgroups with
$f(B)\subset B'$.  The map of right-coset spaces
$A/B\to A'/B'$ is injective if and only if
\[
 f^{-1}(B')=B,
\]
and is surjective if and only if $A'=f(A)B'$.  If in addition
$f$ is surjective and $\ker(f)\subset B$, then
\[
 B=f^{-1}(B')\quad\Longrightarrow\quad
 N_A(B)=f^{-1}(N_{A'}(B')),
 \qquad N_A(B)/B\simeq N_{A'}(B')/B'.
\]

\begin{proof}
The fiber over the identity coset is $f^{-1}(B')/B$, and translating
by an element of $A$ gives all other nonempty fibers.  The image
consists exactly of cosets meeting $f(A)$.  For the last assertion,
if $B=f^{-1}(B')$ and $f$ is onto, an element conjugates $B$ onto
itself exactly when its image conjugates $B'$ onto itself.
The map on normalizers is therefore onto and its kernel is
$\ker(f)\subset B$; quotienting gives the isomorphism.
This proves the precise formal mechanism, and also specifies the
extra hypotheses, behind \S28, folio 181, scans 358--359.  In a
restriction map to a characteristic open subgroup one need not
have surjectivity of $f$; the displayed last assertion cannot be
used without verifying it or an appropriate replacement.
\end{proof}

\statement{B.14}{Open}{Bijectivity for arbitrary profinite prediscretizations}
For a characteristic open subgroup $V\subset\pi$, put
\[
 A=\operatorname{Aut}_{\mathrm{lac}}(\pi),\quad
 A'=\operatorname{Aut}_{\mathrm{lac}}(V),\quad
 B=\overline{\operatorname{Aut}_{\mathrm{lac}}^+(D)},\quad
 B'=\overline{\operatorname{Aut}_{\mathrm{lac}}^+(D\cap V)}.
\]
Restriction gives a continuous injection $f:A\to A'$, and B.12
gives $f(B)\subset B'$.  The unrestricted bijectivity proposed in
scan 355 requires both
\[
 f^{-1}(B')=B,\qquad A'=f(A)B'.
\]
We have not proved these two assertions for arbitrary
loop-preserving automorphisms.  A concrete characteristic test,
already indicated in \S28, folio 183, scan 363, is
\[
 \pi=\widehat F_2,\qquad
 V=\ker(\widehat F_2\to(\mathbb Z/2)^2),
\]
the degree-four covering of the three-punctured sphere of type
$(0,6)$.  Among noncharacteristic coverings, the degree-two cover
in B.11 is the first basic test of the set map of B.12; its raw
failure of surjectivity disappears in the closure quotient and
does not decide this test.  The congruence property is known for
both surfaces in these tests, so restating these conditions as
CSP does not supply the missing argument.  The analogous filling
conditions are made exact in C.8.  As in B.7, Open here records
the unresolved unrestricted statement in this audit, without
claiming a literature classification of every individual cover.
For completeness, injectivity of the restriction homomorphism
uses slimness, not CSP: if an automorphism is the identity on
$V$, then for each $a\in\pi$ the element $a^{-1}f_0(a)$
centralizes the normal open subgroup $V$, and hence is trivial
by D.14--D.15.  Here $f_0$ denotes that automorphism, to
distinguish it from the restriction homomorphism $f$.

\statement{B.15}{Refuted}{Raw discretization injectivity does not transfer to filling}
If the proposed injection of raw discretizations in item 3 of the
prompt is applied also to filling a puncture, it is false.  It
already fails for $(0,4)\to(0,3)$.

\begin{proof}
Let $D=F(x,y,z)$ be the fundamental group of a four-punctured
sphere with peripheral representatives $x,y,z,(xyz)^{-1}$,
and let $\pi=\widehat D$.  Filling the puncture represented by
$x$ gives
\[
 q:\pi\longrightarrow\widehat F(y,z).
\]
Choose $\alpha\in\widehat{\mathbb Z}=\prod_p\mathbb Z_p$ with
$\alpha_2=0$ and $\alpha_p=1$ for every odd prime.  This is not
an integer: an integer equal to zero in $\mathbb Z_2$ is zero,
whereas its $3$-adic component here is one.  Set $h=x^\alpha$.
The retraction $D\to\langle x\rangle$, extended to completions,
shows that
$D\cap\overline{\langle x\rangle}=\langle x\rangle$;
hence $h\notin D$.  B.8 gives $hDh^{-1}\ne D$.
Both are oriented discretizations, because inner automorphisms
preserve all peripheral conjugacy and orientation data.  But
$q(h)=1$, so
\[
 q(hDh^{-1})=q(D)=F(y,z).
\]
Thus the raw filling map is not injective.  The code records
compatible finite Chinese-remainder approximations to $\alpha$;
the proof of nonintegrality is the argument just given, not a
claim that finitely many residues certify it.  The two
discretizations are inner-conjugate, so this counterexample leaves
both versions of the prediscretization question untouched.
\end{proof}

All statements in this part concern orientable surfaces with the
specified orientation option.  Allowing orientation reversal does
not turn them into statements about fundamental groups of
nonorientable surfaces.  A real curve is handled by its arithmetic
extension over $\mathbb R$, whose quotient is $C_2$; the image of
complex conjugation is exactly why the full extended group in
B.5 behaves differently from the oriented geometric image.  The
full profinite results retain the primes $2$ and $3$; no argument
has inverted either prime.  The use of reduction modulo $3$ in
B.2 only distinguishes the two orientation signs.

\clearpage\section{Filling punctures and reconstructing the complete curve}\label{part:C}

The three rows in \S32, folio 202, scan 399 must be read separately.
They concern an arithmetic image, a geometric completion, and the
full group of loop-preserving outer automorphisms.  The first two
have exact filling sequences.  The unrestricted third row, marked
``??'', requires a further assertion about the full automorphism
group.  It is not a formal consequence of either of the first two.
Throughout this part a filling map from type $(g,n+1)$ has hyperbolic
target $(g,n)$ unless the exceptional case is explicitly discussed.

\subsection{The quotient and the two geometric exact sequences}

\statement{C.1}{Proved}{Filling is quotient by the closed normal inertia subgroup}
Let $U=X\setminus S$ be a smooth geometrically connected curve over
a field $K$ of characteristic zero, with smooth proper completion
$X$, and let $S_0\subset S$ be a Galois-stable subset.  Put
$U'=X\setminus(S\setminus S_0)$.  After choosing compatible
geometric base points, the inclusion $U\hookrightarrow U'$ induces
a continuous epimorphism
\[
 \pi_1^{\mathrm{\acute et}}(U)\longrightarrow
 \pi_1^{\mathrm{\acute et}}(U')
\]
whose kernel is the closed normal subgroup generated by the
geometric inertia groups at the points of $S_0$.  The same statement
holds for the geometric groups, and the quotient is compatible
with the maps to $G_K$.  In particular, the quotient classifies
precisely the finite \emph{\'etale} coverings extending across the
filled punctures.

\begin{proof}
Normalize $U'$ in the function field of a connected finite \'etale
cover of $U$.  Normalization is finite for curves of finite type
over a field.  The resulting finite morphism is \'etale at a point
of $S_0$ exactly when its inertia action on the associated finite
set is trivial; in characteristic zero no wild-inertia qualification
is needed.  Thus finite \'etale covers of $U'$ identify with the
full subcategory of finite continuous $\pi_1(U)$-sets on which all
these inertia groups act trivially.  The fundamental group of
this Galois category is the quotient by their closed normal
closure.  The same argument after base change to $\overline K$
gives the geometric assertion.  The structural maps to
$\operatorname{Spec}K$ commute with the inclusion, giving
compatibility with $G_K$.  A change of geometric base point or
connecting path changes the homomorphisms by the corresponding
inner automorphisms of the target.  This is the group-theoretic
filling operation used in \S28, folios 175--176, scans 347--348,
and \S32, folios 200--202, scans 395--400.
\end{proof}

\statement{C.2}{Proved}{The full profinite Birman sequence}
For $2g-2+n>0$, forgetting the last marked point gives a short exact
sequence of profinite groups
\[
 1\longrightarrow\pi_{g,n}
 \longrightarrow\widehat P_{g,n+1}
 \longrightarrow\widehat P_{g,n}\longrightarrow1.
 \tag{C.1}
\]
The first arrow is the completed point-pushing homomorphism.
Allowing permutations of the retained punctures replaces the two
mapping class groups by $M^+_{g,n+1,i}$ and $M^+_{g,n}$, where the
source fixes the puncture being forgotten.  The same version holds
with orientation reversal allowed.

\begin{proof}
The discrete Birman sequence is
$1\to\Pi_{g,n}\to P_{g,n+1}\to P_{g,n}\to1$.
We justify the potentially delicate passage to full completions.
The action by conjugation on its normal subgroup $\Pi_{g,n}$
gives $P_{g,n+1}\to\operatorname{Aut}(\pi_{g,n})$, hence a map
from its profinite completion.  The composite
\[
 \pi_{g,n}\longrightarrow\widehat P_{g,n+1}
 \longrightarrow\operatorname{Aut}(\pi_{g,n})
\]
is the inner-automorphism map, which is injective because
$\pi_{g,n}$ is center-free.  Thus the first arrow of (C.1) is
injective.  Completion is right exact in the following precise
sense: $\widehat P_{g,n+1}$ modulo the closure of the point-pushing
subgroup is $\widehat P_{g,n}$, because both classify the same
finite quotients trivial on that subgroup.  The image of
$\pi_{g,n}$ is compact and contains the dense discrete
point-pushing group, so it is exactly that closure.  This proves
exactness.  The finite permutation and orientation quotients give
the other versions by the identical argument.  The exact sequence
is also explicitly recorded in
\citet[\S2, p.~4]{Boggi2009BC}.
\end{proof}

The hyperbolic target hypothesis is essential.  Filling the sole
puncture of a torus induces the isomorphism
$P_{1,1}\simeq P_{1,0}\simeq\operatorname{SL}_2(\mathbb Z)$;
the alleged kernel $\pi_{1,0}=\widehat{\mathbb Z}^2$ would be
wrong.  Point-pushing is killed by translations isotopic to the
identity on the unmarked torus.  Similarly,
$P_{0,3}\to P_{0,2}$ is a map of trivial groups, not an extension
with kernel $\pi_{0,2}=\widehat{\mathbb Z}$.  For targets $(0,0)$
and $(0,1)$ the pure oriented groups and geometric fundamental
groups are trivial, so the corresponding trivial sequence carries
no anabelian content.  For hyperbolic types the \emph{surface}
group kernel is center-free, even in cases where a \emph{mapping
class} group has center: the elliptic involution at $(1,1)$ and the
hyperelliptic involution at $(2,0)$ are retained.  For the pure
procongruence groups the exact center statement is
\citet[Theorem 4.4]{Boggi2020BC}: the center equals the discrete
center, of order two in these two cases and trivial in the other
hyperbolic pure cases.  Additional centers in variants permitting
permutations must not be suppressed by this pure-group statement.

\statement{C.3}{Known}{The procongruence Birman sequence}
For every hyperbolic target type there is an exact sequence
\[
 1\longrightarrow\pi_{g,n}
 \longrightarrow\check P_{g,n+1}
 \longrightarrow\check P_{g,n}\longrightarrow1.
 \tag{C.2}
\]
It is compatible with (C.1), and the map on the left-hand kernels
is the identity under their natural identifications.  The
point-pushing subgroup therefore has its full profinite topology
inside the procongruence completion.

This is \citet[Corollary 4.7, with $\mathscr C=\mathscr F$,
and Proposition 2.4(ii)]{Boggi2020BC}.  Its Remark 4.8 specifies
the history: Matsumoto's theorem treated open targets, and
Hoshi--Mochizuki extended it to closed targets.  The earlier
attached \citet[Theorem 2.5]{Boggi2009BC} alone has the extra
condition $n\geq1$ and cannot be cited for filling the last
puncture of a higher-genus surface.  Corollary 4.7 also has a
pro-$p$ congruence form, with kernel the pro-$p$ surface group,
for every prime $p$, including $2$ and $3$.  This does not identify
an arbitrary pro-$p$ congruence completion with the pro-$p$
completion of the whole discrete mapping class group.

\statement{C.4}{Proved}{The congruence kernel is unchanged by puncturing}
The natural map
\[
 C_{g,n+1}\longrightarrow C_{g,n}
\]
is an isomorphism for every hyperbolic target $(g,n)$.

\begin{proof}
Compare (C.1) and (C.2).  If an element of $C_{g,n+1}$ maps to
the identity of $\widehat P_{g,n}$, it lies in $\pi_{g,n}$; this
subgroup maps injectively to the lower middle group, so the
element is trivial.  Conversely, take $a\in C_{g,n}$ and lift it
to $b\in\widehat P_{g,n+1}$.  Its image in
$\check P_{g,n+1}$ lies in the left-hand copy of $\pi_{g,n}$.
Multiply $b$ by the inverse of the corresponding element of the
upper copy.  The resulting lift lies in $C_{g,n+1}$.  The map is
therefore a continuous bijection between profinite groups and
hence a topological isomorphism.  This proves exactly which
congruence property transfers in either direction under filling;
it does not assume its truth.
\end{proof}

\subsection{The arithmetic image and the revised quotient}

Let $\mathcal M_{g,n}$ now denote the stack of smooth curves with
\emph{ordered} markings over $\mathbb Q$.  With compatible
geometric base points put
\[
 \mathcal A_{g,n}=\pi_1^{\mathrm{\acute et}}(\mathcal M_{g,n}),
 \quad
 \rho_{g,n}:\mathcal A_{g,n}\to O_{g,n},
 \quad
 N^{\mathrm{arith},!}_{g,n}=\operatorname{im}(\rho_{g,n}).
\]
Its geometric image is $\check P_{g,n}$.  The image is closed
because its source is profinite.  This is the ordered version of
the revised definition in \S29, folio 188, scan 371, not the
normalizer definition of \S26.

\statement{C.5}{Proved}{The revised arithmetic quotient is the absolute Galois group}
For every hyperbolic type, the structural quotient induces a
canonical isomorphism of profinite groups
\[
 G_{\mathbb Q}\ \xrightarrow{\ \sim\ }\
 N^{\mathrm{arith},!}_{g,n}/\check P_{g,n}.
 \tag{C.3}
\]
Moreover, the kernel of the universal monodromy representation
$\rho_{g,n}$ is exactly $C_{g,n}$, viewed in its geometric
subgroup $\widehat P_{g,n}$.  With unordered markings the
corresponding assertion is
$N^{\mathrm{arith}}_{g,n}/\Sigma_{g,n}\simeq G_{\mathbb Q}$.

\begin{proof}
The standard geometric/arithmetic exact sequence of the moduli
stack gives
\[
 1\to\widehat P_{g,n}\to\mathcal A_{g,n}
 \to G_{\mathbb Q}\to1.
\]
Consequently its quotient map to the right-hand side of (C.3)
is surjective by definition.  Let its kernel be $J$.
Choose a marked curve of this type over a finite extension
$L/\mathbb Q$.  Such a curve exists by choosing a geometric
curve and its finitely many markings over $\overline{\mathbb Q}$
and descending their finite defining data; no assertion about
arbitrarily many rational points over $\mathbb Q$ is needed.
The associated point of the moduli stack supplies a section over
$G_L$.  \citet[Corollary 6.4, pp.~66--67]{HoshiMochizuki2011}
says that the preimage in $G_L$ of the geometric monodromy image
is trivial.  Hence $J\cap G_L=1$, so $J$ is a finite normal
subgroup of $G_{\mathbb Q}$.

There is no nontrivial finite normal subgroup of $G_{\mathbb Q}$.
Indeed, a finite normal subgroup is centralized by an open
subgroup, whereas every open subgroup of $G_{\mathbb Q}$ is
slim and has trivial centralizer in $G_{\mathbb Q}$; see
\citet[Corollary 2.2 in the published version; Corollary 1.3
in the author preprint]{Mochizuki2006} and D.14.
Thus $J=1$.  The kernel of $\rho_{g,n}$ is now contained in
the geometric subgroup, where it is by definition $C_{g,n}$.
The unordered version follows by the finite permutation
extension; alternatively its pure subgroup and geometric image
give the same quotient, precisely as in \S32, folio 200,
scan 395, equations (1)--(3).
\end{proof}

The required arithmetic input is a published anabelian theorem.
No section conjecture and no conjectural identification of
$G_{\mathbb Q}$ with $\widehat{\mathrm{GT}}$ is used.
Corollary 6.5 of the same Hoshi--Mochizuki paper states the useful
consequence that total universal monodromy is faithful if and only
if its geometric restriction is faithful.  Thus C.5 does not
silently solve the high-genus congruence problem.

\statement{C.6}{Proved}{The arithmetic filling sequence and invariance of the quotient}
For every hyperbolic target there is a continuous short exact
sequence
\[
 1\longrightarrow\pi_{g,n}
 \longrightarrow N^{\mathrm{arith},!}_{g,n+1}
 \longrightarrow N^{\mathrm{arith},!}_{g,n}
 \longrightarrow1.
 \tag{C.4}
\]
The quotient map on the groups in (C.3) is the identity of
$G_{\mathbb Q}$.  The corresponding sequence with retained
markings unordered uses the stabilizer of the forgotten marking
in the source.  In particular, the faithfulness condition (f) and
the filling isomorphism of the revised arithmetic quotients in
\S32, folios 201--202, scans 397--400 hold for every hyperbolic
type.

\begin{proof}
The universal punctured curve over $\mathcal M_{g,n}$ is
$\mathcal M_{g,n+1}$.  The fundamental-group homomorphism fits
into the arithmetic exact sequence
\[
 1\to\pi_{g,n}\to\mathcal A_{g,n+1}
 \to\mathcal A_{g,n}\to1,
\]
whose geometric row is (C.1) and whose quotient row is the
identity of $G_{\mathbb Q}$.  One may obtain exactness directly
by this diagram: lift an element over a prescribed Galois
element, then adjust by geometric surjectivity; the kernel lies
in the geometric row.  Universal monodromy is compatible with
filling, so quotient this sequence by its monodromy kernels.
By C.5 they are $C_{g,n+1}$ and $C_{g,n}$, and by C.4 the map
between them is an isomorphism.  Their intersection with
$\pi_{g,n}$ is trivial by C.3.  It follows that the quotient
sequence is precisely (C.4).  The arithmetic quotients in the
resulting diagram are induced by the same structural morphism
to $\operatorname{Spec}\mathbb Q$, hence their comparison is
canonically the identity on $G_{\mathbb Q}$.

The same argument over a finite extension $L$ gives the identity
on $G_L$, compatibly with restriction from $G_{\mathbb Q}$.
Changing the base points conjugates the group homomorphisms by
geometric elements.  It has no effect on their arithmetic
quotients, which is the precise sense in which their isomorphism
is canonical.  For coverings defined over a number field the
arithmetic quotients likewise compare through that same field
group.  This last observation concerns arithmetic images; it
does not identify the full normalizers of arbitrary geometric
structures under every abstract finite covering.
\end{proof}

On the scan of page 399, the first sequence has $N$, rather than
the based automorphism group $M$, and its arrow is directed
towards the type with one fewer puncture.  Equation (C.4) is its correct
arithmetic interpretation.  On scan 395, equation (4) omits the
pure-marking superscript on the geometric denominator in its
left-hand quotient; equation (3) immediately above supplies the
mathematically consistent version.  We record both observations
as evidence for the editors, not as silent changes to the source.

\subsection{What the sequence marked ``??'' requires}

\statement{C.7}{Proved}{The canonical based map supplied by a distinguished cusp}
Fix a representative inertia subgroup $I_i\subset\pi_{g,n+1}$
and put $\pi'=\pi_{g,n+1}/\overline{\langle I_i\rangle}^{\mathrm{normal}}$
of type $(g,n)$.  There is a continuous homomorphism
\[
 \theta_i:O_{g,n+1,i}\longrightarrow
 \operatorname{Aut}_{\mathrm{lac}}(\pi'),
 \tag{C.5}
\]
where $O_{g,n+1,i}$ stabilizes the conjugacy class of $I_i$.
Its composite with the outer quotient is the usual filling map
$f_i:O_{g,n+1,i}\to O_{g,n}$.

\begin{proof}
Represent an outer automorphism by an actual automorphism
preserving $I_i$, which is possible by the definition of its
stabilizer.  It induces an automorphism of the quotient $\pi'$.
Two such representatives differ by conjugation by an element
of $N_{\pi_{g,n+1}}(I_i)=I_i$; the equality follows from
\citet[Proposition 3.6]{BoggiZalesskii2019BC}.  That conjugation
is killed in $\pi'$.  Thus the induced genuine automorphism is
independent of the representative.  Composition and continuity
follow by passing from the closed stabilizer of $I_i$ in the
automorphism group to its compact quotient.  Replacing $I_i$ by
a conjugate conjugates the construction by the corresponding
inner automorphism of $\pi'$.  This is the based ambiguity.
On the point-pushing subgroup
$\pi_{g,n}\subset\check P_{g,n+1}\subset O_{g,n+1,i}$,
the map is the inner-automorphism embedding of $\pi_{g,n}$ in
$\operatorname{Aut}(\pi')$, as is seen first in the discrete
Birman sequence and then by continuity.
\end{proof}

\statement{C.8}{Open}{The unrestricted filling assertion and its exact equivalent forms}
With the choices of C.7, the unrestricted third row of scan 399
would assert
\[
 1\longrightarrow\pi_{g,n}
 \longrightarrow O_{g,n+1,i}
 \xrightarrow{f_i}O_{g,n}\longrightarrow1.
 \tag{C.6?}
\]
This is equivalent to $\theta_i$ in (C.5) being an isomorphism.
It is also equivalent to bijectivity of the induced filling map
on the unrestricted oriented prediscretization spaces
\[
 O_{g,n+1,i}/\Sigma_{g,n+1,i}
 \longrightarrow O_{g,n}/\Sigma_{g,n}.
 \tag{C.7?}
\]
Here the source is canonically the corresponding marking
quotient with the chosen cusp filled: the geometric permutation
group permits the cusp to be fixed before taking the quotient.
For individually labeled punctures use the pure variants in
both numerator and denominator.  We have not proved (C.6?)
for the full $O$ groups.  Its smallest hyperbolic target instance
is $(0,4)\to(0,3)$.  The remaining conditions in this instance,
and in general, are explicitly
\[
 f_i(O_{g,n+1,i})=O_{g,n},\qquad
 \ker(f_i)=\pi_{g,n}.
\]

\begin{proof}[Proof of the stated equivalences]
Since $\pi'$ is center-free, the lower row of
\[
\begin{array}{ccccccccc}
1&\to&\pi_{g,n}&\to&O_{g,n+1,i}&\xrightarrow{f_i}&O_{g,n}&\to&1\\
 &&\Vert&&\downarrow\theta_i&&\Vert\\
1&\to&\operatorname{Inn}(\pi')&\to&
\operatorname{Aut}_{\mathrm{lac}}(\pi')&\to&O_{g,n}&\to&1
\end{array}
\]
is exact, with the left identification supplied by C.7.
Exactness of the upper row is consequently equivalent to
$\theta_i$ being an isomorphism, by elementary lifting and
kernel arguments.  C.3 with permutations gives
$f_i(\Sigma_{g,n+1,i})=\Sigma_{g,n}$ and kernel
$\pi_{g,n}$ on that geometric subgroup.  By B.13, (C.7?) is
bijective if and only if $f_i$ is onto and
$f_i^{-1}(\Sigma_{g,n})=\Sigma_{g,n+1,i}$.  Under geometric
surjectivity, the latter is equivalent to
$\ker(f_i)\subset\Sigma_{g,n+1,i}$, and then the known
geometric kernel makes it equivalent to the second condition
above.  This proves all the equivalences without assuming any
of them.
\end{proof}

Thus the two geometric rows and the arithmetic row have been
settled, while the unrestricted double-hat row has been reduced
to its exact missing content.  The map to genuine automorphisms
in scan 348 and the map to outer automorphisms in scan 399 are
not interchangeable: the former would be an isomorphism; the
latter would have the point-pushing kernel.  The manuscript's
``peut-\^etre enti\`erement faux'' in folio 176, scan 349 is a
question, not a counterexample.  As in B.7 and B.14, the label
Open here specifies what this audit has not established; no
claim is made that a comprehensive literature proof of
independence has been supplied.  In particular, we do not
rename this unresolved lifting problem the congruence subgroup
problem: both involved genus-zero congruence assertions are known.

\statement{C.9}{Known}{The established combinatorial cuspidalization theorem}
Let $X$ be a hyperbolic curve over an algebraically closed field
of characteristic zero, and let $X_m$ be its configuration space
of $m$ ordered distinct points.  Let $\Lambda$ be either all
primes or a singleton, and let $\pi_m$ be the maximal pro-$\Lambda$
quotient of $\pi_1(X_m)$.  Define
$\operatorname{Out}^{\mathrm{FC}}(\pi_m)$ by preservation of
the fiber subgroups of all configuration projections and of the
cuspidal inertia classes in their surface-group subquotients,
as in the cited definition.  Then forgetting one point gives
\[
 \operatorname{Out}^{\mathrm{FC}}(\pi_{m+1})
 \longrightarrow\operatorname{Out}^{\mathrm{FC}}(\pi_m)
\]
which is injective for all $m\geq1$, bijective for $m\geq3$ if
$X$ is affine, and bijective for $m\geq4$ if $X$ is proper.
The factor-permutation subgroup centralizes these FC-admissible
outer automorphisms.

These are the exact assertions of
\citet[Theorem 6.1, pp.~61--63]{HoshiMochizuki2011}, extending
\citet[Theorem 4.1]{Mochizuki2010BC}.  In the original definition,
F-admissibility preserves each fiber subgroup, rather than
merely the collection up to an unspecified permutation;
FC-admissibility adds cuspidal compatibility on the relevant
subquotients.  The image of an F-admissible automorphism from
one higher configuration level is FC-admissible, by
\citet[Theorem A and Corollary 1.9]{HoshiMochizukiTopicsBC}.
These extra data are a substantive hypothesis.  A single
filled surface group is a fiber in this construction, not the
whole configuration-space group, and an arbitrary member of
$O_{g,n+1,i}$ has not thereby been supplied with such a lift.
Consequently this theorem does not by itself remove the
question mark in (C.6?).  Its full and single-prime forms include
the primes $2$ and $3$ in characteristic zero; its positive
characteristic version requires the characteristic to lie
outside the chosen set of primes.

\subsection{What it means to reconstruct the entire \texorpdfstring{\'etale}{etale} topos}

\statement{C.10}{Refuted}{The finite-cover topos is not the full \texorpdfstring{\'etale}{etale} topos}
For a connected positive-dimensional smooth curve $X$, the
topos of continuous discrete $\pi_1^{\mathrm{\acute et}}(X)$-sets
need not, and in fact cannot, be its full small \'etale topos.
Thus C.1 alone does not reconstruct ``tout le topos \'etale''.

\begin{proof}
In the classifying topos of any group, a subobject of the
terminal object is an invariant subset of a one-element set.
There are exactly two such subobjects.  In the small \'etale
topos of a curve, every Zariski open subset represents a
subterminal object.  A nonempty proper open subset gives a
third one, distinct from the empty and the whole scheme.
An equivalence of topoi preserves subterminal objects, so the
topoi cannot be equivalent.  One particularly transparent case
is $X=\mathbb P^1_{\overline{\mathbb Q}}$: its geometric
fundamental group is trivial, while its \'etale topos still
contains all its distinct open subobjects.  The finite
two-versus-three subobject count is also recorded in
\texttt{code/bc\_counterexamples.py}.  The argument applies to
the classifying topos of all continuous discrete sets, so it
does not rely on confusing finite sets with all such sets.
\end{proof}

\statement{C.11}{Proved}{Reconstruction of the complete curve and its full \texorpdfstring{\'etale}{etale} topos}
Fix a sub-$p$-adic field $K$ of characteristic zero, in
particular a number field or a finite extension of
$\mathbb Q_p$.  On the essential image of hyperbolic curves
over $K$, the arithmetic extension
\[
 1\longrightarrow\pi_1(U_{\overline K})
 \longrightarrow\Pi_U=\pi_1(U)
 \longrightarrow G_K\longrightarrow1
 \tag{C.8}
\]
reconstructs the smooth proper completion $X$ of $U$ and its
entire small \'etale topos, functorially for $K$-isomorphisms.
The group datum includes the projection to the specified
$G_K$; isomorphisms of it are taken modulo inner automorphisms
from the geometric kernel.  The resulting topos is determined
up to the equivalence induced by the reconstructed
$K$-isomorphism, with the usual natural-isomorphism ambiguity.

\begin{proof}
The relative anabelian Isom theorem in D.4 follows from
\citet[Theorem A, equivalently Theorem 16.5 and its following
Remark]{Mochizuki1999}: for hyperbolic curves over a sub-$p$-adic
field, the map from dominant $K$-morphisms to continuous open
homomorphisms of arithmetic fundamental groups over $G_K$,
modulo geometric inner automorphisms of the target, is
bijective.  Applied to an isomorphism and its inverse, this
reconstructs a unique $K$-isomorphism $U\simeq U'$ inducing
the prescribed outer class.  Consequently the functor on
isomorphism groupoids is fully faithful, and an inverse on
its essential image reconstructs $U$ with precisely this
uniqueness.  Existence is a hypothesis on the input datum,
not a claim that every extension of $G_K$ is geometric.

A smooth geometrically connected curve has a unique smooth
proper completion up to unique isomorphism extending the
identity on $U$.  It may be constructed as the nonsingular
proper model of its one-variable function field $F=K(U)$.
An isomorphism of open curves extends uniquely to these
proper models.  It remains to show explicitly that this
reconstructs the whole topos rather than just its finite
coverings.

For every finite separable field extension $L/F$, let $X_L$
be the normalization of $X$ in $L$.  It is finite over $X$.
Consider the site whose objects are finite disjoint unions
of opens $V\subset X_L$ on which $V\to X$ is \'etale.
Morphisms are $X$-morphisms, equivalently the compatible
generic field embeddings and their extensions on the chosen
opens.  Covering families are jointly surjective \'etale
families.  All these data are determined by the recovered
curve $X$, and hence by (C.8).

This site is a basis for the small \'etale site of $X$.
Indeed, a connected separated \'etale scheme $Y$ of finite
presentation over $X$ is normal and integral, and its
function field is a finite separable extension $L/F$.
Zariski's main theorem identifies $Y$ with an open subscheme
of $X_L$.  One can use the precise factorization of
\citet[Lemmas 37.43.2 and 37.43.3, Tags 02LR and 05K0]{StacksBC}:
a separated quasi-finite morphism factors as an open
immersion followed by a finite morphism; normalizing and
using the same generic field gives the specified
normalization.  Every \'etale scheme over $X$ is covered
by separated affine \'etale opens, and the quasi-compact
ones are of finite presentation.  Hence these objects
form a covering basis even when a general object of the
small site is not separated or quasi-compact.  The
comparison lemma for sites now identifies its sheaf
topos with $\operatorname{Sh}(X_{\mathrm{\acute et}})$.
Transport of normalized extensions and opens along
$K$-isomorphisms proves functoriality.
\end{proof}

This theorem gives a complete affirmative formulation of the
reconstruction sought in \S32, folios 206--217, scans 407--430.
The exact phrase quoted in the prompt was not located as a
literal sentence in scans 395--420; the subsequent construction
of varying open groupoids and coverings occurs in scans
421--430.  We cite that development rather than manufacture
a quotation.  The indispensable input in C.11 is the
\emph{arithmetic} group over its field, and the proof uses
relative anabelian reconstruction before constructing the
site.  A bare geometric $\pi_{g,n}$ is the same for every
complex curve of that type and cannot recover its algebraic
moduli.  Nor does retaining only the finite covers unramified
over the original $U$ produce the displayed site: finite
extensions branched at newly removed points are also needed.
We do not claim that the manuscript's direct construction
from an independently characterized class of admissible
splittings has thereby been completed without anabelian
reconstruction.

\statement{C.12}{Proved}{The correct point-filling statement over a field}
For a smooth curve $U'$ over a field $K$, a $K$-rational point
$x\in U'(K)$, and $U=U'\setminus\{x\}$,
\[
 U(K)=U'(K)\setminus\{x\}.
\]
Under the anabelian point maps this gives the corresponding
bijection of \emph{geometric point-induced} conjugacy classes
of splittings, compatibly with filling.  If the deleted set
is merely Galois-stable, the correct statement uses its
$K$-rational points; a point of that set need not itself be
rational.  The statement does not assert that every group
section is point-induced.

\begin{proof}
A morphism $\operatorname{Spec}K\to U'$ factors through the
open complement precisely when its image is not $x$.
The equality follows.  The maps on arithmetic fundamental
groups commute with these morphisms, so the same equality
holds for their images in the set of section classes;
injectivity of the point map in the stated anabelian setting
is the theorem recorded in D.5.  Nothing in this argument
requires surjectivity of the section map.  This supplies
the field-valued form of the expectation in \S32, folios
211--212, scans 417--420, while keeping the word
``admissible'' tied to actual geometric points.
\end{proof}

\statement{C.13}{Refuted}{Deleting one section is not deleting one element of a section set}
The analogue of C.12 for sections over an arbitrary base is
false: distinct sections of a family can intersect.  Thus
the assertion in scan 418 cannot be imposed on arbitrary
profinite acting groups by treating all sections as a
set of disjoint geometric points.

\begin{proof}
Over $\mathbb Q$ put
$S=\mathbb P^1\setminus\{0,1,\infty\}$ and consider the
constant hyperbolic family $S\times S\to S$, projection to
the second factor.  The sections
\[
 s(t)=(t,t),\qquad e(t)=(2,t)
\]
are distinct and intersect at $t=2$.  Delete the image of
$e$ from the family.  Its fibers now have type $(0,4)$,
but $s$ does not lift to a section of the deleted family:
at $2$ its value is exactly the deleted point.  It is
therefore false that the new section set is the old one
minus the single section $e$.  All objects and morphisms
are defined over $\mathbb Q$, and the intersection is
checked by substituting $t=2$ in $t-2$ in the accompanying
code.  This also gives a geometric example for the acting
group $\pi_1(S)$, so the defect is not an artifact of
inventing unrelated group actions.

The manuscript itself recognizes this obstruction later:
\S32, folios 219--220, scans 433--435 withdraws the
overgeneralized formulation because sections of families
can meet.  We preserve that withdrawal.  C.12 over a
field and C.11 via the reconstructed curve are the
valid replacements needed for the present task.
\end{proof}

\clearpage\section{Arithmetic curves, inertia, sections, and faithfulness}
\label{sec:D}

Throughout this section a hyperbolic curve is a smooth geometrically connected
curve $U/K$ with smooth proper completion $X$, reduced complement $S=X\setminus U$,
and type $(g,n)=(g(X),\#S(\overline K))$ satisfying $2g-2+n>0$.
Unless a different coefficient set is explicitly specified, all fundamental groups
are full profinite \'etale groups.  We choose compatible geometric base points and
write
\[
 1\longrightarrow\Delta_U\longrightarrow\Pi_U
   \xrightarrow{p_U}G_K\longrightarrow1,
 \qquad \Delta_U=\pi_1(U_{\overline K}),\quad G_K=\operatorname{Gal}(\overline K/K).
\]
Changing a geometric base point changes maps by geometric inner conjugacy.
An inertia subgroup is the image of the geometric fundamental group of the
punctured strict henselian local scheme at a point of $S_{\overline K}$;
it is a conjugacy class of closed subgroups $I_s\simeq\widehat{\mathbb Z}(1)$.
We write $\operatorname{Sec}(U/K)$ for continuous sections of $p_U$, modulo
conjugation by $\Delta_U$, and $\rho_U:G_K\to\operatorname{Out}(\Delta_U)$ for
the outer action.  The topology on automorphism groups is the topology of
uniform agreement on finite characteristic quotients.

The distinction between an extension, a rational-point section of that extension,
and a section of a moduli-stack extension is essential in reading \S34.
We first settle the statements that do not depend on the three 2026 proof packs.

\subsection{What the arithmetic extension determines}

\statement{D.1}{Proved}{A centre-free extension is its outer-action pullback}
Let $\Delta$ be a topologically finitely generated centre-free profinite group
and let $1\to\Delta\to E\xrightarrow{p}G\to1$ be an extension of profinite groups.
If $\rho:G\to\operatorname{Out}(\Delta)$ is its associated outer action, then
\[
 E\longrightarrow
 \operatorname{Aut}(\Delta)\mathbin{\times}_{\operatorname{Out}(\Delta)}G,
 \qquad e\longmapsto(\operatorname{conj}_e|_\Delta,p(e))
\]
is an isomorphism of extensions.  Sections of $p$, modulo $\Delta$-conjugacy,
therefore correspond exactly to homomorphic lifts of $\rho$ to
$\operatorname{Aut}(\Delta)$, modulo $\operatorname{Inn}(\Delta)$-conjugacy.
This justifies the equivalence asserted in \S34, folio 245, scan 483.
\begin{proof}
The displayed map is a continuous homomorphism.  Its kernel lies in $\Delta$
and centralizes $\Delta$, hence is trivial.  Given $(a,g)$ in the fibre product,
choose $e\in E$ above $g$.  Equality of outer classes says
$a=\operatorname{conj}_d\operatorname{conj}_e$ for some $d\in\Delta$; thus
$(a,g)$ is the image of $de$.  A continuous bijection between compact Hausdorff
groups is an isomorphism.  Both assertions about sections follow from this
isomorphism.  Notice that $E\to\operatorname{Aut}(\Delta)$ by itself need not be
injective: it is the pair with $p$ that is always injective.
\end{proof}

\statement{D.2}{Refuted}{The literal open-normalizer criterion in scan 483}
It is false that the cuspidal inertia subgroups are the maximal closed subgroups
$L\subseteq\Delta_U$ for which $p_U(N_{\Pi_U}(L))$ is open in $G_K$.
It is also false that every nontrivial closed subgroup with this property is
contained in one inertia subgroup, as suggested in the margin of scan 483.
These failures already occur for $U=\mathbb P^1_{\mathbb Q}\setminus\{0,1,\infty\}$.
\begin{proof}
The subgroup $L=\Delta_U$ is normal in $\Pi_U$, so its normalizer is all of
$\Pi_U$ and its projected normalizer is $G_K$.  It is maximal among subgroups
with the stated property simply because it contains every subgroup under
consideration.  Here $\Delta_U\simeq\widehat F_2$ has quotient
$(\mathbb Z/2)^2$, whereas every inertia subgroup is procyclic.  Thus it cannot
be an inertia subgroup or be contained in one.  This is an example in the
arithmetic group of the curve, not merely an abstract profinite extension.

Even adding the word ``proper'' would not repair the marginal assertion.
The characteristic subgroup $[\Delta_U,\Delta_U]$ is proper, is normalized by
$\Pi_U$, and is not procyclic.  For the last assertion map the two free generators
onto $(12345)$ and $(123)$ in $A_5$.  They generate $A_5$: conjugates give
adjacent $3$-cycles, which generate all even permutations.  Since $A_5$ is
perfect, the image of $[\Delta_U,\Delta_U]$ is $A_5$, which is not cyclic.
The attached computation \texttt{code/d\_counterexamples.py} verifies these
finite calculations and records its output.

Finally, if the condition were instead imposed on subgroups of $\Pi_U$, every
rational-point section $s(G_K)$ would qualify because its normalizer contains
$s(G_K)$.  Such a subgroup projects isomorphically onto $G_K$ and is not
geometric inertia.  For instance $2\in U(\mathbb Q)$ gives such a section.
This explains why rational-point decomposition groups must not be silently
identified with cuspidal inertia.
\end{proof}
The scan actually says subgroups ``dans $\pi$''; no doubtful glyph supplies the
missing procyclic or cyclotomic hypotheses.  The alternatives
$\Gamma^+$ and $\Gamma$ in the projected-normalizer phrase make no difference:
$G_K$ is open in $G_{\mathbb Q}$.

\statement{D.3}{Known}{The cyclotomic recognition theorem for affine inertia}
Let $U/K$ be an affine hyperbolic curve over a number field.  A nontrivial
element $z\in\Delta_U$ belongs to a cuspidal inertia subgroup if and only if
there is a closed subgroup $N\subseteq\Pi_U$ such that $N$ normalizes
$\overline{\langle z\rangle}$, $p_U(N)$ is open in $G_K$, and
\[
 hzh^{-1}=z^{\chi_K(p_U(h))}\qquad(h\in N),
\]
where $\chi_K:G_K\to\widehat{\mathbb Z}^{\times}$ is the full cyclotomic
character.  In particular the action factors through $p_U(N)$; this condition
includes triviality of the conjugation action of $N\cap\Delta_U$.
The maximal procyclic subgroups having this property are precisely the
cuspidal inertia subgroups.  The collection is intrinsic to the arithmetic
extension together with its cyclotomic character.
\begin{proof}
The equivalence is \citet[Theorem (2.1.1)]{Nakamura1994}, applied to the full
class of finite groups.  The source assumes $n\geq1$, characteristic zero,
and a number field; all these hypotheses have been retained here.
For a cusp the branch-cycle description of its decomposition group gives
the displayed conjugation law.  Conversely Nakamura's weight argument identifies
$z$ as peripheral.  If a procyclic subgroup satisfying the condition is generated
topologically by $z$, it lies in the inertia subgroup supplied by the theorem.
The inertia subgroup itself satisfies the condition and is therefore the maximal
one.  For completeness, distinct maximal peripheral inertia subgroups intersect
trivially: if $1\ne x$ belongs to two of them, \citet[Proposition 3.6]{BoggiZalesskii2019BC}
gives $C_{\Delta_U}(x)=I$ for the first subgroup $I$, since a peripheral
generator is algebraically simple.  The second subgroup centralizes $x$, hence
is contained in $I$ and equals $I$ by maximality.  This also gives uniqueness
of the containing inertia subgroup for a nontrivial cyclotomic subgroup.
\end{proof}
We do not erase ``cyclotomic'' from this theorem.  The literal assertion refuted
in D.2 and the published assertion in D.3 have different hypotheses.  An
isomorphism of arithmetic extensions over $G_K$ preserves the latter hypotheses,
so it recovers the inertia conjugacy classes without their being supplied as
extra data.  The conclusion for arbitrary open homomorphisms follows more
cleanly from the Hom theorem below than by treating a normalizer calculation
as if it commuted with a noninjective homomorphism.

\statement{D.4}{Known}{The relative Hom theorem, including its exact ambiguity}
Let $K$ be sub-$p$-adic, meaning a subfield of a finitely generated extension
of $\mathbb Q_p$, and let $U,V$ be hyperbolic curves over $K$.  There is a
bijection
\[
 \operatorname{Hom}^{\mathrm{dom}}_K(U,V)
 \xrightarrow{\ \sim\ }
 \operatorname{Hom}^{\mathrm{open}}_{G_K}(\Pi_U,\Pi_V)/\operatorname{Inn}(\Delta_V).
\]
On the right the maps are continuous homomorphisms with open image whose
composite with $p_V$ is $p_U$.  In particular this applies over every number
field and every finite extension of $\mathbb Q_p$, including $p=2,3$.
A homomorphism in this set takes a cuspidal inertia subgroup either to the
trivial group or into a conjugate of a target cuspidal inertia subgroup.
\begin{proof}
The bijection is \citet[Theorem A; Theorem 16.5 and its following
Remark]{Mochizuki1999}.  We use the full profinite version.  The paper also has a
geometrically pro-$p$ version, whose arithmetic quotient is still the full
$G_K$; it is not a theorem about replacing the whole arithmetic group by its
maximal pro-$p$ quotient.

For clarity we verify the peripheral consequence, which is the assertion of
\S34, folios 245--245 verso, scans 483--484.  The associated nonconstant morphism
$f:U\to V$ extends uniquely to the smooth proper completions $\bar f:X\to Y$.
If a deleted point of $X$ maps to a point of $V$, its local inertia maps trivially
in $\pi_1(V)$.  If it maps to a deleted point of $Y$, choose local parameters
$t,u$ at the source and target.  The equality $\bar f^*u=t^e a$ with $a$ a local
unit shows that tame geometric inertia maps into the target inertia by
multiplication by the positive ramification index $e$.  Characteristic zero
makes all these horizontal inertia groups tame.  Thus ``respects loops''
allows a loop to be killed; it does not mean a permutation of all punctures.
\end{proof}
The relative Isom theorem follows formally: the two homomorphisms associated
to a group isomorphism and its inverse give morphisms whose composites are
identities by the injectivity of this bijection.  No separate compatibility
with a previously chosen loop system is required.

\statement{D.5}{Proved}{Admissible effective lifts are exactly the point sections}
Let $U/K$ be hyperbolic with $K$ sub-$p$-adic.  Call a lift of $\rho_U$
\emph{geometrically admissible} if it is obtained from a hyperbolic curve $V/K$,
a point $v\in V(K)$, and a $G_K$-equivariant outer isomorphism
$\Delta_V\simeq\Delta_U$ by transporting the action defined by $v$.
Then the admissible lifts modulo $\Delta_U$-conjugacy are in bijection with
$U(K)$.  This proves the deduction in \S34, scan 485, folio read
$246\mid245$, without assuming surjectivity of the section map.
\begin{proof}
By D.1 the given outer isomorphism produces an isomorphism
$\Pi_V\simeq\Pi_U$ over $G_K$, with the stated geometric ambiguity.
The Isom consequence of D.4 gives a unique $K$-isomorphism $f:V\to U$.
Functoriality carries the section at $v$ to the section at $f(v)$.
Conversely the point $u\in U(K)$ supplies the required admissible datum by
taking $V=U$, $v=u$, and the identity outer isomorphism.
Finally two points cannot give the same section class:
\citet[Theorem C = Theorem 19.1]{Mochizuki1999} proves injectivity of the
point-to-section map under precisely these field and hyperbolicity hypotheses.
Thus the two constructions give the asserted bijection.  Changing the chosen
representative of the outer isomorphism conjugates the lift by $\Delta_U$,
which is exactly the ambiguity already factored out.
\end{proof}
This proposition is about lifts that are already assumed admissible.  It does
not say that every lift is admissible.  For a proper curve of genus at least two
that additional assertion is the section conjecture; for an affine curve the
presence of cuspidal sections changes the assertion itself.

\statement{D.6}{Refuted}{All affine sections need not come from interior points}
For $U=\mathbb P^1_{\mathbb Q}\setminus\{0,1,\infty\}$ there is a continuous
section of $\Pi_U\to G_{\mathbb Q}$ which is not the section of any point of
$U(\mathbb Q)$.  Thus an unqualified affine version of ``every lift is the
lift of a rational point of $U$'' is false.
\begin{proof}
Take the rational tangential base point at $t=0$ with tangent vector $1$.
The associated fibre functor gives a section $s_0$.  More explicitly, embed the
local punctured disk in the Puiseux field
$\bigcup_m\overline{\mathbb Q}((t^{1/m}))$ and let $G_{\mathbb Q}$ act on
coefficients while fixing the compatible symbols $t^{1/m}$.  On the finite
covers of $U$ this gives the tangential section; changing paths conjugates it
by $\Delta_U$.

Under the morphism $t:U\to\mathbb G_m$, its Kummer class in
$H^1(G_{\mathbb Q},\mu_m)=\mathbb Q^\times/(\mathbb Q^\times)^m$ is zero for
every $m$, since each $t^{1/m}$ is fixed by the section.  If $s_0$ were the
section of $u\in U(\mathbb Q)$, the same classes would be the classes of $u$.
Hence $u$ would belong to $\bigcap_{m\geq1}(\mathbb Q^\times)^m$.
Every rational prime valuation of such a number is divisible by every $m$,
so all valuations vanish and $u=\pm1$.  Taking $m=2$ removes $-1$, giving
$u=1$, which is not in $U$.  This contradiction proves the assertion.
\end{proof}
The appropriate affine section conjecture includes rational cusps and their
cuspidal sections.  D.6 does not refute that conjecture.

\subsection{The section manuscript: an audited conditional chain}

\statement{D.7}{Open}{The 2026 proper $p$-adic section claim}
The assertion
\[
 X(k)\xrightarrow{\sim}\operatorname{Sec}(X/k)
 \quad(k/\mathbb Q_p\text{ finite},\ X/k\text{ proper},\ g(X)\geq2)
\]
for every prime $p$ is claimed in \citet[Theorem 1.1]{OpenAI019}.
The precise new input is its Theorem 1.2: for $F=k(X)$, a fixed separable
closure $F^{\mathrm{sep}}$, and the subfield $F^{\mathrm{et}}$ generated by all
connected finite \'etale covers of $X$ with all their $F$-embeddings,
\[
 F^{\mathrm{sep}}\subseteq\widehat{F^{\mathrm{et}}}
 \quad\hbox{inside }\widehat{F^{\mathrm{sep}}}
 \tag{D.7.1}
\]
for every fixed prolongation of every real-valued rank-one norm on $F$
extending the normalized $p$-adic norm.  The copies of the fields in this
inclusion are specified; an abstract embedding of a finite extension into
some unrelated completion would not suffice.
We retain the label Open for the unconditional new claim in this audit:
claimed, unrefereed, and not completely verified by us.  We have found no
counterexample.  The conditional implication D.8 is proved below.

\paragraph{What we checked in Sections 3--6 of the pack.}
We read these sections in full and followed their logical order, including
the distinction between horizontal ramification in characteristic zero and
wild reduction in characteristic $p$.
The patching construction in Proposition 3.2 requires a root on every
component above each overlap annulus.  With that hypothesis the pulled-back
branched cover splits on the overlap, so gluing trivial covers over the inner
disks produces a cover \'etale over the entire proper curve.  Proper GAGA and
invariance of finite \'etale covers under extension between algebraically
closed characteristic-zero fields are used after gluing.  Lemma 3.3 then
evaluates at the chosen completion-valued point.  We checked why this yields
a field embedding rather than merely an analytic map.  Lemma 3.4 uses cyclic
horizontal inertia over $\overline k((t))$: divisibility of the branch indices
by one integer is preserved under compositum and Galois closure, and an
upper cover \'etale over a lower cover introduces no further horizontal
ramification.

For Lemma 4.2 we checked the graph constructions separately.  A cycle in the
stable dual graph carries a tame cyclic cover of degree $\ell\ne p$, with
expansion factor $\ell$ on each lifted edge.  A separating edge is first
lifted to parallel edges using an unramified cover on the two sides; stability
ensures that each side has a cycle or a component of positive genus.
We checked that taking fibre products and then connected components retains
the required roots on all remaining sheets.  The external lifting result is
\citet[Theorem 7.7]{ABBR2015}; its tame-cover definition includes absence of
additional branch points when the marked puncture set is empty.
For Lemma 4.3 we opened \citet[Theorem A and Proposition 2.4(iv)]{MochizukiTsujimura2024}.
The latter gives a stable model of a finite \'etale cover dominating a prescribed
model.  Taking an irrational valuation level excludes vertices and smooth
disk tubes, so the inverse image of a sufficiently small rational-ended
collar lies in node annuli.  This is the specific model domination needed;
semistable reduction alone would not provide it.  The accessible primary
version is dated June 16, 2023; we did not silently identify it with the
March 30, 2024 revision cited by the pack.

In Section 5 the period lattice in a totally degenerate Jacobian is torsion-free.
Consequently local lifts of a multiplication-by-$m$ pullback can be adjusted
by periods and glued without a residual torsion ambiguity.  We checked that
Lemma 5.1 is intended to give a root of a torus character on the whole pullback,
not just one component.  Proposition 5.2 uses
$y^\ell=T(T-1)/(T-\alpha)$, with $\ell$ odd and $\ell\ne p$,
whose skeleton has two rational vertices joined by $\ell$ edges.
A cycle formed from two edges supplies a torus character with slope one on
the chosen edge.  If $u=cT(1+h)$ and $s=v(T)$ on an annulus of length $a$,
the Laurent coefficient inequalities give
$v(h)\geq\min(s,a-s)$.  The trimming condition
\[
 \min(s,a-s)>v_p(m)+1/(p-1)
\]
is sufficient for an analytic $m$th root of $1+h$, also when $p\mid m$.
We checked this numerical point; replacing it by a prime-to-$p$ root argument
would not prove Theorem 1.2.  We also checked the use of normality in Section 6:
a completion-valued embedding of a finite normal extension into the fixed
completion forces its specified image there to be the same normal subfield.

This is a substantial audit of the proof's mechanism, not a complete independent
verification of all analytic constructions.  In particular we have not fully
reproduced the asserted simultaneous stable-model and toric-uniformization
construction of Proposition 5.2, with its identification of the chosen
annular coordinate on every inverse component.  The opened Abel--Jacobi
skeleton comparison of \citet[Proposition 6.1 and Corollary 6.6]{BakerRabinoff2015}
explains the slope calculation but is not by itself the whole analytic
construction.  This is the remaining verification obligation at which we
stop promoting (D.7.1) to a theorem.  All later uses of (D.7.1) below are
explicitly conditional.  The separate attachment
\emph{\'Etale covers with a prescribed exterior sheet} is not cited by the
main section manuscript and is not used here; its theorem is claimed,
steps not checked by us.  Nor do we use Corollary 8.4, which depends on an
unattached Selmer-converse manuscript.

\statement{D.8}{Proved}{The completion assertion implies the proper local section theorem}
Assume (D.7.1) for a fixed proper hyperbolic curve $X/k$, with
$k/\mathbb Q_p$ finite.  Then $X(k)\to\operatorname{Sec}(X/k)$ is bijective.
Thus the group-theoretic and valuative passage in Section 7 of the pack does
not require an additional affine or birational-lifting conjecture.
\begin{proof}
Put $F=k(X)$, $F^{\mathrm{et}}$ as above, and prolong a rank-one valuation $w$
extending the $p$-adic valuation to $F^{\mathrm{sep}}$.
Restriction gives
\[
 D_{\bar w}(F^{\mathrm{sep}}/F)
   \longrightarrow D_{\widetilde w}(F^{\mathrm{et}}/F).
 \tag{D.8.1}
\]
It is surjective: lift an automorphism of $F^{\mathrm{et}}/F$ to
$F^{\mathrm{sep}}/F$, then correct the lift using the transitive action of
$\operatorname{Gal}(F^{\mathrm{sep}}/F^{\mathrm{et}})$ on prolongations.
The infinite-extension assertion follows from the finite assertions by
compactness.  It is injective because an element of its kernel is an isometry
fixing $F^{\mathrm{et}}$, hence fixes its completion, which contains the
specified $F^{\mathrm{sep}}$ by (D.7.1).  Normalization $w(p)=1$ ensures that
stabilizing the valuation ring preserves the absolute value itself.  The two
decomposition groups are closed, so (D.8.1) is an isomorphism of profinite groups.

No such rank-one decomposition group in $\Pi_X$ can contain the image of a
section.  Otherwise (D.8.1) lifts it to a birational section
$\widetilde s:G_k\to G_F$.  By
\citet[Proposition 2.4(2), Definition 1.4, and Observation 1.3]{Koenigsmann2005},
the entire image $\Sigma$ is also contained in a decomposition group of
$\operatorname{ord}_a$ for some $a\in X(k)$.
In the fixed field $M=(F^{\mathrm{sep}})^\Sigma$ the restrictions $v_1,v_2$
of these two selected prolongations are nontrivial rank-one henselian
valuations.  Henselianity follows because $G_M=\Sigma$ stabilizes the
prolongation, while transitivity gives uniqueness.  Algebraic extension
changes the value group by a torsion quotient and hence preserves its rank.
They satisfy $v_1(p)=1$ and $v_2|_{k^\times}=0$.

Here is the explicit contradiction, including the dyadic case.  Choose a
nonsquare $d\in k^\times$.  It is still a nonsquare in $M$, because
$\Sigma\to G_k$ is onto.  Let $t\in F$ be a uniformizer at $a$, choose an
integer $N>v_1(t)$, and set $z=p^N/t$.  Then $v_1(z)>0$ and $v_2(z)<0$.
For positive integers $r$ put
\[
 q_r=d+(1-d)\frac{z^r}{1+z^r}.
\]
The denominator is nonzero.  Direct valuation computations give
\[
 \begin{split}
 v_1(q_r/d-1)&=v_1((1-d)/d)+r v_1(z)\longrightarrow+\infty,\\
 v_2(q_r-1)&=v_2(d-1)-r v_2(z)\longrightarrow+\infty.
 \end{split}
\]
For large $r$ these values exceed $2v_1(2)$ and $2v_2(2)$, respectively.
Hensel's lemma at $1$ gives square roots of $q_r/d$ and $q_r$ in $M$;
the quotient makes $d$ a square.  The contradiction uses neither discrete
value groups nor $p\ne2$.

Now let $s:G_k\to\Pi_X$ be arbitrary.
\citet[Theorem 26 in the 2011 preprint]{PopStix2017} places its image in the
decomposition group of a valuation of $F$ extending the $p$-adic valuation.
The valuation transcendence inequality gives rational rank at most two,
hence ordered rank at most two.  Rank one has just been excluded.
For rank two, let $u$ be the rank-one coarsening.  Its prolongation is preserved
by the section: the relevant convex subgroup of the value group has a unique
extension under an algebraic field extension, because the quotient of value
groups is torsion.  If $u$ is nontrivial on $k$, the preceding rank-one
argument excludes it.  Otherwise properness gives a closed centre $a\in X$,
and a valuation dominating the discrete valuation ring $\mathcal O_{X,a}$
and trivial on $k$ is equivalent to $\operatorname{ord}_a$.
Its residue field is $k(a)$, and the remaining refinement is the unique
$p$-adic valuation of this finite extension.  The image on constants of its
horizontal decomposition group is $G_{k(a)}$.  Since $s$ projects onto $G_k$,
we obtain $k(a)=k$.  Horizontal inertia is trivial in the full \'etale tower
of the proper smooth curve, so this decomposition group is precisely the
point section at $a$.  This is also the type-$2h$ calculation in
\citet[Section 6.1]{PopStix2017}.  Surjectivity follows.

For completeness, injectivity can be checked without (D.7.1).  If $a,b\in X(k)$
give the same section, use $a$ for the Abel--Jacobi embedding into $J$.
The multiplication-by-$m$ torsor is finite \'etale for every $m\geq1$ in
characteristic zero.  Equality of section classes implies
$[b-a]\in mJ(k)$ for every $m$.  An open subgroup of the compact analytic
group $J(k)$ is $\mathbb Z_p^{g[k:\mathbb Q_p]}$.  If it has index $q$,
then $\bigcap_m mJ(k)\subseteq\bigcap_r p^r\mathbb Z_p^{g[k:\mathbb Q_p]}=0$,
using $m=qp^r$.  Thus $[b-a]=0$.  A nontrivial equivalence of the two distinct
degree-one divisors would give a degree-one map $X\to\mathbb P^1$;
this is impossible for positive genus.  Hence $a=b$.
\end{proof}

\statement{D.9}{Proved}{The exact conditional global descent statement}
Assume the proper local section theorem in D.7 over every finite completion
of a number field $K$.  For a proper hyperbolic curve $C/K$, define
\[
 C(\mathbb A_K)_\bullet=
 \prod_{v\nmid\infty} C(K_v)\times
 \prod_{v\text{ real}}\pi_0(C(K_v)).
\]
Complex places are omitted.  Let $C(\mathbb A_K)_\bullet^{\mathrm{f\text{-}cov}}$
be the locus surviving every torsor under a finite $K$-group scheme, with
survival meaning liftability to some global twist.  Then
\[
 \operatorname{loc}(\operatorname{Sec}(C/K))
   = C(\mathbb A_K)_\bullet^{\mathrm{f\text{-}cov}},
 \qquad
 C(K)\xrightarrow{\sim}\operatorname{Sec}(C/K)
 \ \Longleftrightarrow\
 C(\mathbb A_K)_\bullet^{\mathrm{f\text{-}cov}}=C(K).
 \tag{D.9.1}
\]
This is the precise conditional content of Proposition 8.1 and Theorem 8.2
of \citet{OpenAI019}.  The general equality on the right remains an additional
number-theoretic assertion, not a consequence of local point-theoreticity.
\begin{proof}
The local theorem gives a unique point at each finite place; the real section
theorem gives a connected component at each real place
\citep[Introduction, Theorem]{BrescianiVistoli2020}.  Thus localization is
well-defined.  Apply \citet[Theorem 2.1 and Remark 2.2(1), preprint
numbering]{HarariStix2012} with its normal subgroup $U=1$ and omitted set of
places $S=\varnothing$.  Its equivalence between condition (i) and condition
(iii) is exactly the equivalence between survival of all finite torsors and
interpolation of the local point sections by a global section, modulo the
full geometric group.  Evaluations are constant on real components and
complex places impose no condition.  We use the corrected convention of
\citet[Section 1]{Stoll2017Errata}; disconnected torsors are included.
This proves the first equality.

If the descent locus is $C(K)$, the local points of a section are the
localizations of one $a\in C(K)$.  Its localization has finite support
$\{a\}$, so \citet[Corollary 5]{Stix2015} makes the section geometric.
Properness excludes cuspidal sections.  Injectivity is D.5, or follows by
restriction to a single finite completion.  Conversely bijectivity of the
point-to-section map and the first equality immediately give the descent
equality.  These arguments also cover empty sets.
\end{proof}

\statement{D.10}{Proved}{What the modular-curve consequence says about a moduli section}
Assume the local theorem of D.7.  Let $N$ satisfy $g(X_1(N))\geq2$, and let
$s:G_{\mathbb Q}\to\pi_1(Y_1(N))$ be a section.  Its composite with
$\pi_1(Y_1(N))\to\pi_1(X_1(N))$ is geometrically conjugate to the section of
a unique rational cusp of $X_1(N)$.  Consequently the same conclusion applies
to a section of the arithmetic $N_{1,1}$ whose image is contained in the
open subgroup belonging to $Y_1(N)$, under the faithful stack identification
of Part E.  It does not assert that $s(G_{\mathbb Q})$ is contained in a
cuspidal decomposition subgroup of $\pi_1(Y_1(N))$.
\begin{proof}
The forgetful map from the fine modular curve $Y_1(N)$ to
$\mathcal M_{1,1,\mathbb Q}$ is finite \'etale in characteristic zero, so its
fundamental group is an open subgroup, up to conjugacy.  Containment of the
image of a section makes the same homomorphism a section of this open
subgroup's projection onto $G_{\mathbb Q}$.  Here $g(X_1(N))\geq2$ ensures
that $N$ is outside the small levels with stabilizers on $Y_1(N)$.
The compactification homomorphism is over $G_{\mathbb Q}$ and hence its
composite is a section of the proper curve.

\citet[Corollary 8.8]{Stoll2007}, with \citet{Stoll2017Errata}, proves
$X_1(N)(\mathbb A_{\mathbb Q})_\bullet^{\mathrm{f\text{-}ab}}=X_1(N)(\mathbb Q)$
when the genus is positive.  Since rational points are contained in the
finite-cover locus and that locus is contained in the finite-abelian locus,
D.9 applies.  This verifies the implication used in
\citet[Corollary 8.5]{OpenAI019} without its unrelated Corollary 8.4.
Finally \citet[Introduction, Theorem (7), p.~35]{Mazur1977} states directly
that all rational points of $X_1(N)$ are cusps whenever its genus is positive.
The equivalent torsion-order assertion is Theorem (7'); the full torsion
classification is Theorem (8).  Thus the proper section comes from a cusp,
and D.5 gives its uniqueness.

The kernel of compactification is generated normally by cusp inertia.
Identifying the image of $s$ after quotienting by this kernel does not identify
$s$ before quotienting.  D.7 is a theorem claimed for proper curves and
provides no such lifting assertion for $Y_1(N)$.  Nor does the conclusion
imply nonexistence of these splittings: a nonzero rational tangent vector at
any rational cusp gives a cuspidal section of $Y_1(N)$, exactly as in D.6.
\end{proof}

\subsection{Relative and absolute open homomorphisms}

\statement{D.11}{Proved}{An absolute open homomorphism preserves the geometric kernel}
Let $U/K$ and $V/K'$ be hyperbolic curves over number fields and let
$h:\Pi_U\to\Pi_V$ be a continuous homomorphism with open image, with no
compatibility assumed.  Then $h(\Delta_U)\subseteq\Delta_V$, and there is a
unique continuous homomorphism with open image
$\alpha:G_K\to G_{K'}$ such that $p_Vh=\alpha p_U$.
This conclusion does not depend on the 2026 open-Hom manuscript.
\begin{proof}
Set $J=p_V(h(\Pi_U))$.  The projection $p_V$ is an open map of profinite
groups, so $J$ is open in $G_{K'}$, and $J=G_L$ for a finite extension
$L/K'$ in the chosen algebraic closure.  The subgroup
$A=p_V(h(\Delta_U))$ is compact, hence closed, is topologically finitely
generated because $\Delta_U$ is, and is normal in $J$ because $\Delta_U$
is normal in $\Pi_U$.  By \citet[Theorem 2.4]{Mochizuki2006}, a number
field's absolute Galois group has no nontrivial topologically finitely
generated closed normal subgroup.  Thus $A=1$.  The universal property of
the quotient gives a unique continuous $\alpha$, and its image is $J$.
\end{proof}
The same argument is the mechanism behind
\citet[Proposition 3.3]{Mochizuki2006} for isomorphisms.  Merely calling a
geometric subgroup ``characteristic'' would not justify its preservation
by an arbitrary homomorphism; the proof above supplies the missing step.
Openness matters because it makes $J$ an absolute Galois group of a number
field, to which the theorem applies.

\statement{D.12}{Proved}{The conditional absolute Hom theorem and its known cases}
Assume \citet[Theorem 1.1]{OpenAI031} for absolute Galois groups of number
fields.  Then every $h$ as in D.11 is induced, up to target geometric inner
conjugacy after fixing the base-field map, by a unique field embedding
$j:K'\hookrightarrow K$ and a unique dominant $K$-morphism
\[
 f:U\longrightarrow V\otimes_{K',j}K.
\]
More invariantly there is a bijection
\[
 \coprod_{j:K'\hookrightarrow K}
 \operatorname{Hom}^{\mathrm{dom}}_K(U,V\otimes_{K',j}K)
 \ \simeq\
 \operatorname{Hom}^{\mathrm{open}}(\Pi_U,\Pi_V)\big/\operatorname{Inn}(\Pi_V).
 \tag{D.12.1}
\]
Without the 2026 claim the same conclusion is known whenever the induced
$\alpha$ of D.11 is injective, or is cyclotomic-compatible, and in particular
for isomorphisms $h$.  It is also known when the source field is $\mathbb Q$.
\begin{proof}
Apply the assumed theorem to $\alpha$.  It gives a unique compatible embedding
$\bar j:\overline {K'}\to\overline K$ with
$\alpha(\sigma)=\bar j^{-1}\sigma\bar j$ and $j=\bar j|_{K'}$ taking values
in $K$.  Between algebraic closures of number fields such an embedding is
onto: its image is algebraically closed and contains $\mathbb Q$, so it
contains every element algebraic over $\mathbb Q$.  Consequently $\alpha$
is injective as well as having open image.

Base change identifies the arithmetic group of $V\otimes_{K',j}K$ with
\[
 P=\Pi_V\mathbin{\times}_{G_{K'},\alpha}G_K.
\]
The map $\widetilde h(x)=(h(x),p_U(x))$ takes values in $P$ by D.11 and is
over $G_K$.  The first projection identifies $P$ with the open subgroup
$p_V^{-1}(\alpha(G_K))$ of $\Pi_V$, since $\alpha$ is injective.
Thus $\widetilde h$ has open image in $P$.  D.4 supplies the unique dominant
$f$ with the prescribed geometric conjugacy class.

We check the larger ambiguity in (D.12.1).  Conjugating $h$ by $e\in\Pi_V$
conjugates $\alpha$ by $p_V(e)\in G_{K'}$; the associated change of $\bar j$
does not alter its restriction to $K'$.  Hence $j$ is intrinsic to the absolute
outer class.  Once $\alpha$ is fixed, an $e$ conjugating one representative
to another over the same $\alpha$ has $p_V(e)$ centralizing the open subgroup
$\alpha(G_K)$.  This centralizer is trivial by
\citet[Corollary 2.2]{Mochizuki2006}.  The remaining conjugation is therefore
geometric and D.4 gives exactly the stated uniqueness of $f$.  Conversely a
pair $(j,f)$ gives such an open homomorphism: the relative map has open image
and the base-change subgroup is open in $\Pi_V$.

For the unconditional variants, an injective $\alpha$ is an isomorphism onto
an open subgroup $G_L$; the Neukirch--Uchida isomorphism theorem gives the
compatible field embedding.  One precise solvably-closed formulation is
\citet[Theorem, p.~359]{Uchida1979}, also reproduced as
\citet[Theorem 3.1]{Mochizuki2006}.  For cyclotomic compatibility use
\citet[Theorem 3.4]{Hoshi2025}.  If $h$ is an isomorphism, apply D.11 also
to $h^{-1}$, so that $\alpha$ is an isomorphism.  If $K=\mathbb Q$, use
\citet[Theorem 1]{Uchida1981}.  The rest of the proof is unchanged.
\end{proof}
For the original relative formulation in \S34, scans 483--484, the base
fields lie in one $\overline{\mathbb Q}$ and the maps are already over
$G_{\mathbb Q}$.  Its use of the Hom theorem is unconditional.  The additional
absolute assertion in D.12 is the part depending on the new claim; D.11 is
not an unresolved geometric-kernel gap.

\statement{D.13}{Open}{The 2026 unrestricted Galois open-Hom claim}
Let $E_i/F_i$ be solvably closed Galois extensions of number fields, meaning
that $E_i$ have no nontrivial abelian algebraic extensions.  The assertion
that every continuous homomorphism with open image
\[
 \alpha:\operatorname{Gal}(E_1/F_1)\longrightarrow
                \operatorname{Gal}(E_2/F_2)
\]
is induced by a unique compatible field embedding $E_2\hookrightarrow E_1$
is \citet[Theorem 1.1]{OpenAI031}, the conjecture on p.~595 of
\citet{Uchida1981}.  We record it as claimed, unrefereed, with its full
unconditional validity not completely verified by us.  D.12 is a proved
conditional consequence, not an assertion that the conjecture has thereby
been settled.

\paragraph{Proof audit of the open-Hom manuscript.}
We read Sections 1--6, including the full proofs of the new lemmas rather
than only the outline, and checked the following dependencies.
The reduction to $P=\operatorname{Gal}(\mathbb Q^{\mathrm{sol}}/\mathbb Q)$
is valid: a solvably closed algebraic field contains the roots of unity and
is closed under radicals, hence contains $\mathbb Q^{\mathrm{sol}}$.
For a number field $M\subset\mathbb Q^{\mathrm{sol}}$ this is also its maximal
prosolvable extension; a solvable normal closure of $M/\mathbb Q$ followed by
a compositum of conjugates proves the assertion.
We opened the precise local input in \citet[Section 1, pp.~595--597 and
Proposition 1]{Uchida1981}: for a fixed odd prime $\ell$, nontrivial
pro-$\ell$ inertia forces the nonabelian local case and hence an injective
map on that inertia into a unique target prime; all but finitely many target
primes occur.  The excluded torsion case is dyadic.  This result gives
cofinite coverage for each fixed $\ell$, not a bound uniform in $\ell$.
The claimed uniform bound later in the pack is therefore a real new step.
We also opened the uniqueness statement, Proposition 2 of Uchida, and
\citet[Theorem 3.4]{Hoshi2025}; the latter proves that full cyclotomic
compatibility is equivalent to geometricity.

Lemma 3.1 supplies the additive-character detection that drives the proof.
We followed its interpolation determinant and its separate translate/degree
zero estimate.  The relevant strict inequality is
\[
 n+n/r>n+1+n/m\qquad(1\leq r<n,\ m>n(n-1)).
\]
It compares the nonarchimedean determinant decay with the archimedean height
bound; no asserted Leopoldt rank is substituted.  In Proposition 3.2 the
logarithmic image of global units is the actual image, and generators of
powers of distinct prime ideals give disjoint divisor supports.  Thus the
claimed independence argument is compatible with class-field theory without
assuming Leopoldt's conjecture.  We checked the distinct roles of the
absolute-degree-one version, the Frobenius-average version, and the
nonzero-corestriction version.  Confusing them would invalidate the later
prime comparisons.

In Section 4 the map is first restricted to a normal open subgroup and then
extended to the normalizer of its kernel.  The regular-module enlargement
in Lemma 4.2 increases subgroup orders while keeping the bound on element
orders fixed.  Proposition 4.3 compares characters via a cyclic subgroup
and its proper subgroups.  We checked that its semisimple projection is
$U$-equivariant for the normal open subgroup $U$; the argument does not
require the larger intermediate subgroups to be normal in the whole
normalizer.  We followed the explicit source and target prime choices,
the residue-degree $f$, and the separation into those powers of Frobenius
that lie in $U$ and those that lie in the other kernel conjugates.
The result is surjectivity and logarithmic cyclotomic compatibility, not yet
full cyclotomic compatibility.

Section 5 passes to $\mathbb Q(i)$ and uses a sign character to show that
almost every split target prime has only finitely many possible source
partners, as the auxiliary prime varies.  Lemma 5.2 is an $S$-unit Kummer
statement; the finite class-group contribution vanishes in the relevant
inverse limit, so it is not a mistaken direct-sum assertion about all primes.
In Proposition 5.3 a finite cyclotomic twist is killed in a cyclic extension.
The extension degree can grow with $\ell$, but the multiplicity of the one
specified character in the $S$-unit representation is bounded by the number
of $S$-orbits.  The resulting bound $2[K:\mathbb Q]+2$ is independent of
$\ell$.  We checked why this independence is needed before applying the
simultaneous Kummer/Chebotarev separation of Lemma 5.4; the fixed-$\ell$
cofiniteness from Uchida cannot replace it.

Finally Section 6 compares character pullbacks at rational primes split
completely in an auxiliary finite field.  An infinite subset with one fixed
conjugate gives equality of the corresponding additive characters by
Proposition 3.2.  Faithfulness of the signed regular representation yields
equality modulo each finite quotient.  Compactness of the sets of possible
conjugating elements, with the finite-intersection property obtained by
passing to composita, supplies a single conjugation.  Extending equality
from a normal open subgroup uses the trivial centralizer of an open subgroup
of $P$.  The resulting full cyclotomic compatibility is the hypothesis of
Hoshi's theorem, which finishes the claimed proof.

These checks identify a coherent chain and expose where an apparently weaker
replacement would fail.  They do not constitute an independently reproduced
proof of every technical generality of the zero estimate in Lemma 3.1, the
character-separation argument of Proposition 4.3, and the uniform multiplicity
bound in Proposition 5.3.  We therefore do not treat those three new inputs
as established merely because their conclusions fit the reduction.  No
verified counterexample or contradictory hypothesis was found.  The attached
eight-lesson exposition is not evidence of independent verification; we do
not use it as an authority.  No Lean proof is supplied for this pack.

\subsection{Centres, restriction to covers, and the faithfulness questions}

\statement{D.14}{Known}{The centres and centralizers needed in \S35}
For a hyperbolic type $(g,n)$ in characteristic zero, the geometric profinite
surface group $\Delta_{g,n}$ is torsion-free and slim: every open subgroup has
trivial centralizer in $\Delta_{g,n}$.  In particular it and all its open
subgroups have trivial centre.  The absolute Galois group of a number field
is also slim, and hence has trivial centre and no nontrivial finite normal
subgroup.  These settle the two centre questions in \S35, folio 251 verso,
scan 496.  The second assertion does not assert that a number field's
absolute Galois group is torsion-free.
\begin{proof}
For $n>0$ the geometric group is free profinite of rank $2g+n-1\geq2$;
for $n=0$ it is the profinite completion of the closed oriented surface group
of genus $g\geq2$.
\citet[Corollaries (1.3.3)--(1.3.4)]{Nakamura1994} give centre-freeness
and the finite projective resolutions obtained from goodness of these groups.
Their cohomological dimensions at every prime are finite (one in the free
case, two in the closed surface case), so they have no elements of prime
order and are torsion-free.  Every open subgroup is again a hyperbolic
surface group of finite index, by the correspondence with finite covers.
Its Euler characteristic is its index times the original negative Euler
characteristic; thus the same centre assertion applies.

To derive slimness directly, suppose $c$ centralizes an open subgroup $H$.
Its normal core $H_0$ is open; some positive power $c^m$ belongs to $H_0$.
That power lies in the centre of $H_0$, hence is $1$.  Torsion-freeness gives
$c=1$.  This agrees with \citet[Lemma 3.2]{Mochizuki2006}.
For number fields use \citet[Corollary 2.2]{Mochizuki2006}.
If a finite subgroup $F$ is normal in a slim profinite group $G$, the kernel
of the conjugation map $G\to\operatorname{Aut}(F)$ is open and centralized
by $F$, forcing $F=1$.  Complex conjugation, when a real embedding exists,
is an element of order two of $G_K$; its subgroup is not normal.
\end{proof}

\statement{D.15}{Proved}{Rigidity on open subgroups and the two centralizer formulas}
If $G$ is slim and $u\in\operatorname{Aut}(G)$ is the identity on an open
subgroup $H$, then $u=1$.  If $u(H)=H$ and $u|_H$ is conjugation by an
element of $H$, then $u$ is that inner automorphism on $G$.
For a free basis element $x$ of $\widehat F_r$, or a standard nonseparating
simple generator $x$ of the profinite closed surface group of genus at least
two, and for every integer $m\geq1$,
\[
 C_G(x^m)=C_G(\overline{\langle x^m\rangle})
    =\overline{\langle x\rangle}.
 \tag{D.15.1}
\]
Thus Lemma 5 and equations (2), (4) in \S35, scans 498--501,
folios 252 verso--254, are valid with their stated hyperbolic hypotheses.
\begin{proof}
Let $H_0$ be the normal core of $H$ in $G$.  For $g\in G$ and $a\in H_0$
both $a$ and $gag^{-1}$ lie in $H$, so
\[
 u(g)a u(g)^{-1}=u(gag^{-1})=gag^{-1}.
\]
Consequently $g^{-1}u(g)$ centralizes the open subgroup $H_0$, hence is $1$.
This proves the first assertion.  Compose with the inverse of the specified
inner automorphism to obtain the second.

For the free basis case (D.15.1) is
\citet[Lemma (2.1.2)]{Nakamura1994}.
For the surface case, a standard generator is an algebraically simple element
in the sense of \citet[Definition 1.3]{BoggiZalesskii2019BC}.
Their Proposition 3.6 gives the stronger identity
$N_G(\overline{\langle x^m\rangle})=\overline{\langle x\rangle}$;
the centralizer contains the right-hand side and is contained in this
normalizer.  Both formulas follow.  Centralizing one element is equivalent
to centralizing the closed subgroup it topologically generates, so the
scan's alternatives $L_i^m$ and $\ell_i^m$ are mathematically equivalent in
this use.  The subgroup $\overline{\langle x\rangle}\simeq\widehat{\mathbb Z}$
has injective multiplication by every positive integer, which also completes
the power argument written across scans 500--501.
\end{proof}

\statement{D.16}{Known}{Faithful outer actions for every hyperbolic curve}
For every hyperbolic curve $U/K$ over a number field or a $p$-adic local
field, the full profinite outer representation
\[
 \rho_U:G_K\longrightarrow\operatorname{Out}(\Delta_U)
\]
is injective.  This includes affine and proper curves, type $(0,3)$,
type $(1,1)$, and all proper types $(g,0)$ with $g\geq2$.
It proves the expectation in \S35, scan 497, rather than only the weaker
assertion that a sufficiently small open subscheme is faithful.
\begin{proof}
The precise general result is
\citet[Corollary 6.3(i), with one configuration point]{HoshiMochizuki2011}.
Its full-prime-set hypothesis is essential.  The source explicitly identifies
the earlier affine assertion as
\citet[Theorem 2.1]{Matsumoto1996}; the 2011 result supplies the proper case
as well.

We also give the short Belyi argument for the particular theorem proved in
\S35, scans 491--495.  Suppose $\sigma\in G_{\mathbb Q}$ acts trivially in
$\operatorname{Out}(\widehat F_2)$ for the tripod.  Then it fixes the
isomorphism class of every connected finite \'etale cover of the tripod,
since these correspond to conjugacy classes of open subgroups of
$\widehat F_2$.  For any $j\in\overline{\mathbb Q}$ choose an elliptic curve
$E_j$ with that invariant.  By \citet[Section 4, Theorem 4]{Belyi1980} it
has a finite map to $\mathbb P^1$ branched only over $0,1,\infty$.
Deleting the inverse images of the three points gives a cover whose
$\sigma$-conjugate is isomorphic to it.  An isomorphism extends uniquely to
smooth proper completions, so $E_j\simeq E_j^\sigma$ and
$j=\sigma(j)$.  Since every algebraic $j$ was allowed, $\sigma=1$.
The unnumbered corollary on p.~256 of Belyi's English paper is exactly this
faithfulness consequence.  This argument proves injectivity, and has no
surjectivity conclusion about a separately defined normalizer or GT group.
\end{proof}
For fixed $\ell$, the analogous pro-$\ell$ assertion over a number field
is not a consequence of this theorem.  In fact for a curve with good
reduction, and a good finite place of residue characteristic different from
$\ell$, the geometrically pro-$\ell$ outer representation is unramified
there.  Nontrivial inertia of $G_K$ therefore lies in its kernel.
The full profinite statement cannot be replaced by a single-prime statement.

\statement{D.17}{Known}{Faithfulness even modulo geometric Teichm\"uller monodromy}
Let $2g-2+n>0$ and let $\mathcal M_{g,n}/\mathbb Q$ denote the stack with
ordered marked points.  Write
\[
 \Psi:\pi_1(\mathcal M_{g,n})\longrightarrow
                \operatorname{Out}(\Delta_{g,n}),\qquad
 N^{\mathrm{arith}}_{g,n}=\operatorname{im}(\Psi),\qquad
 \Sigma_{g,n}=\operatorname{im}(\pi_1(\mathcal M_{g,n,\overline{\mathbb Q}})).
\]
Then the canonical surjection
\[
 G_{\mathbb Q}\longrightarrow
       N^{\mathrm{arith}}_{g,n}/\Sigma_{g,n}
 \tag{D.17.1}
\]
is an isomorphism.  The same assertion holds with unordered marked points
and the corresponding unordered geometric monodromy group.
Here $\Sigma_{g,n}$ is the procongruence image of the orientation-preserving
mapping class group, not an assumed copy of its entire profinite completion.
This answers affirmatively the precise arithmetic-image question of \S35,
scans 502--505, including the proper types and $(1,1)$ singled out on scan 504
as possibly ``tout \`a fait faux''.
\begin{proof}
The geometric image is closed and normal because it is the continuous image
of the geometric kernel in the arithmetic stack extension.  Hence the
surjection (D.17.1) is well-defined without a rational base point on the stack.
Choose a finite extension $L/\mathbb Q$ and an $L$-point of
$\mathcal M_{g,n}$, equivalently a curve of the given type with ordered
$L$-rational marked points.  Such an $L$ exists by taking a closed point of
a nonempty finite type atlas.  We do not assume the existence of a
$\mathbb Q$-point for an arbitrary prescribed $(g,n)$.
The point gives a section over $G_L$ whose universal outer representation is
$\rho_{U/L}$.  By \citet[Corollary 6.4]{HoshiMochizuki2011},
\[
 \rho_{U/L}^{-1}(\Sigma_{g,n})=1.
\]
Thus the kernel $A$ of (D.17.1) intersects $G_L$ trivially.  It is finite,
normal in $G_{\mathbb Q}$, and therefore trivial by D.14.  This proves the
ordered assertion.

For unordered marks the ordered geometric image is a normal subgroup of
finite index in the unordered image.  Over $L$ the inverse image under
$\rho_{U/L}$ of the latter group is finite, because its map to the finite
quotient has trivial kernel by the ordered assertion of Corollary 6.4.
It is also normal in $G_L$, since geometric monodromy is normal in arithmetic
monodromy.  Slimness forces it to be trivial.  The same finite-normal-kernel
argument over $\mathbb Q$ completes the proof.
\end{proof}
The distinction from the congruence subgroup problem is exact.
\citet[Corollary 6.5]{HoshiMochizuki2011} says that the full universal
monodromy $\Psi$ is injective if and only if its geometric restriction is
injective.  Thus identifying $\Sigma_{g,n}$ with the profinite mapping class
group still requires the relevant congruence theorem.  The quotient
faithfulness (D.17.1) does not.  Equally, replacing $N^{\mathrm{arith}}$ by
a full normalizer is a different problem, retained separately in Part E.
Quotienting by orientation-reversing mapping classes would also change the
assertion; D.21 below gives the obstruction.

\statement{D.18}{Proved}{The direction of faithfulness under finite covers}
Let $k$ have characteristic zero, let $V\to U$ be a finite \'etale morphism
of geometrically connected hyperbolic curves over $k$, and assume that
$\rho_U:G_k\to\operatorname{Out}(\Delta_U)$ is faithful.  Then $\rho_V$
is faithful.  No rational point and no replacement of $k$ by a finite
extension are necessary.  This proves the finite-cover assertion in \S35,
folios 253--253 verso, scans 499--500, in its stated characteristic-zero scope.
\begin{proof}
Choose compatible base points.  Then $\Pi_V\subseteq\Pi_U$ is an open
subgroup, maps onto $G_k$ by geometric connectedness, and
$\Delta_V=\Delta_U\cap\Pi_V$ is open in $\Delta_U$.
For $\sigma\in\ker\rho_V$ choose $e\in\Pi_V$ above $\sigma$.
Its conjugation on $\Delta_V$ is conjugation by some $d\in\Delta_V$.
Thus $d^{-1}e$ centralizes $\Delta_V$.  Its conjugation on $\Delta_U$ is
an automorphism that is the identity on an open subgroup, hence is the
identity by D.14--D.15.  Therefore $\rho_U(\sigma)=1$, whence $\sigma=1$.
\end{proof}
This proof also explains the direction of the comparison map: an open
geometric subgroup detects an automorphism of the larger group once the
inner automorphism on the subgroup has been removed.  In the manuscript's
longer argument a finite residual kernel had to be eliminated; it is not
enough merely to observe that this kernel acts through a finite group.
Filling punctures is a surjection of geometric groups, so it is not the
same argument; the proper faithfulness needed after filling is supplied by
D.16 and D.17.

\statement{D.19}{Proved}{The tripod centralizer is the symmetric group}
For $U=\mathbb P^1_K\setminus\{0,1,\infty\}$ over a sub-$p$-adic field $K$,
\[
 C_{\operatorname{Out}(\Delta_U)}(\rho_U(G_K))
       =\operatorname{Aut}_K(U)\simeq\mathfrak S_3.
\]
In particular the automorphism assertion of \S34, scans 485--486, does not
require assuming a chosen arithmetic normalizer in addition to commuting
with the arithmetic action.
\begin{proof}
An outer automorphism $a$ commuting with $\rho_U$ induces, by D.1, an
isomorphism of the arithmetic extension over $G_K$.
D.4 and its Isom consequence identify it with a unique automorphism of $U$.
Conversely every such scheme automorphism commutes with the Galois action.
An automorphism extends to $\mathbb P^1$ and permutes $\{0,1,\infty\}$;
three distinct points determine a fractional linear transformation, and all
six permutations are defined over $\mathbb Q$.  This gives exactly
$\mathfrak S_3$.  Its injectivity on the outer group follows from the same
Isom theorem.  The alternatives $\mathfrak S_3\mid\mathcal S_3$ in scan 485
are two notations for this six-element group; the mathematics does not decide
the handwritten font.
\end{proof}

\subsection{Small types and the exact remaining qualifications}

\statement{D.20}{Refuted}{Dropping hyperbolicity from the faithfulness assertion}
The uniform assertion of D.16 becomes false if any of $(0,0)$, $(0,1)$,
$(0,2)$, $(1,0)$ is admitted.  The geometric centres and the counterexamples
are as follows:
\[
\begin{array}{c|c|c}
 (g,n)&\Delta_{g,n}&Z(\Delta_{g,n})\\\hline
 (0,0),(0,1)&1&1\\
 (0,2)&\widehat{\mathbb Z}&\widehat{\mathbb Z}\\
 (1,0)&\widehat{\mathbb Z}^{,2}&\widehat{\mathbb Z}^{,2}
\end{array}
\]
In contrast $(0,3)$ and $(1,1)$ are both hyperbolic with geometric group
$\widehat F_2$, but have different peripheral structures.
\begin{proof}
Over an algebraically closed field of characteristic zero, comparison with
the sphere, plane, punctured plane, and torus gives the displayed groups.
Thus the first two outer actions over $\mathbb Q$ are trivial.  For
$\mathbb G_m/\mathbb Q$ the outer action is the cyclotomic character, whose
abelian image cannot be faithful on $G_{\mathbb Q}$.  Explicitly the splitting
field of $x^3-2$ has Galois group $\mathfrak S_3$: Eisenstein's criterion at
$2$ proves irreducibility and the discriminant is $-108$, a nonsquare.
Hence $G_{\mathbb Q}$ has nontrivial commutator subgroup, killed by its
cyclotomic character.

For $(1,0)$ take $E:y^2=x^3-x$ over $K=\mathbb Q(i)$.
The endomorphism $[i](x,y)=(-x,iy)$ is defined over $K$.
Complex uniformization makes the integral lattice a free rank-one
$\mathbb Z[i]$-module, so the full Tate module is free rank one over
$\widehat{\mathbb Z}[i]$.  Equivalently one can choose a basis on which $[i]$
has matrix $J=\left(\begin{smallmatrix}0&-1\\1&0\end{smallmatrix}\right)$.
All Galois matrices commute with $J$.  Solving $AJ=JA$ over any commutative
ring gives $A=aI+bJ$; its invertible solutions commute with each other.
Thus the entire adelic representation has abelian image.  On the other hand
the $\mathfrak S_3$ splitting field of $x^3-2$ intersects $\mathbb Q(i)$
trivially: its only quadratic subfield is $\mathbb Q(\sqrt{-3})$.
Its compositum with $K$ is therefore an $\mathfrak S_3$-extension of $K$.
Consequently $G_K$ is nonabelian and the representation on
$\Delta_E=\widehat{\mathbb Z}^{,2}$ is not faithful.  The code verifies the
discriminant, the finite commutator calculation, and the centralizer
calculation modulo several integers, including $2$ and $3$; the ring
calculation above is the proof for all moduli.

For $(0,3)$ the peripheral generators are $x,y,(xy)^{-1}$, each primitive
in abelianization.  For $(1,1)$ the only peripheral subgroup is generated
by $[x,y]$, which has zero image in abelianization.  A bare identification
of their free profinite groups is therefore not an identification of curves
with their loop data.
\end{proof}

\statement{D.21}{Proved}{Orientation, real fields, and residue characteristics}
The subgroup in D.17 must be the geometric orientation-preserving group.
If a full extended mapping class group containing orientation reversal is
used instead, the analogous faithfulness assertion fails already at the
tripod.  The geometric centre statements remain valid over $\mathbb R$,
but the number-field Galois slimness statement does not extend to
$G_{\mathbb R}$.  None of D.4, D.8, D.16, or D.17 has a residue-characteristic
exclusion at $2$ or $3$ in its stated characteristic-zero setting.
\begin{proof}
Complex conjugation on the complex tripod is an orientation-reversing
homeomorphism preserving its punctures.  Its outer class belongs to the
image of the extended mapping class group and is nontrivial: on the
abelianization $\widehat{\mathbb Z}^{,2}(1)$ it acts by $-1$.
Therefore quotienting by that extended image kills a nontrivial element
of $G_{\mathbb Q}$.  Moreover the extended image need not be normal in
$N^{\mathrm{arith}}$: at the unordered tripod its image under the
isomorphism of D.17 is the subgroup generated by complex conjugation.
If it were normal it would be a nontrivial finite normal subgroup of
$G_{\mathbb Q}$, contradicting D.14.  One cannot silently form an ordinary
group quotient by it.  Taking its normal closure would still kill complex
conjugation.

Over $\mathbb R$ the arithmetic quotient is $C_2$, whose centre is itself.
The correct real section map for a hyperbolic curve is
$\pi_0(U(\mathbb R))\simeq\operatorname{Sec}(U/\mathbb R)$, by
\citet[Introduction, Theorem]{BrescianiVistoli2020}; it is not a bijection
from individual real points.  The outer action of $C_2$ is nevertheless
faithful for a hyperbolic real curve.  If the proper completion has positive
genus, conjugation reverses its orientation and is anti-symplectic on
$H_1$; identity action would contradict its nondegenerate intersection
pairing over $\mathbb Q_\ell$.  If the genus is zero and $n\geq3$, its action
on puncture homology is minus the permutation of the punctures.  An
involution on at least three punctures has at least two orbits, so the
reduced permutation representation has a nonzero $+1$ eigenspace.
The action on that subspace after the minus sign is nontrivial.

A smooth complex algebraic curve is orientable.  The fundamental group of
a nonorientable topological surface is therefore not the geometric group
of one of these curves.  A real curve can have a nonorientable quotient by
conjugation; its arithmetic datum is the geometric oriented group together
with the involution and its extension, not the fundamental group of the
quotient surface alone.  The orientation double cover does not identify
these different arithmetic extensions.

All the local fields in D.4 and D.8 have characteristic zero.  Horizontal
ramification along a deleted point is consequently tame even if its
ramification index is divisible by the residue characteristic; the analytic
unit-root threshold in D.7 handles that divisibility.  The proof of D.8
explicitly uses the correct dyadic Hensel bound.  The universal-monodromy
and full profinite theorems in D.16--D.17 are also characteristic-zero
statements.  In \S36 the orbifold orders $2$ and $3$ do prevent transporting
an argument through tame coarse moduli in characteristics $2,3$; this
requires retaining the stack and is not an exception to the theorems over
$\mathbb Q$ or $\mathbb Q_p$ stated here.
\end{proof}

\statement{D.22}{Proved}{Uniqueness of a Galois-invariant tripod loop structure}
Let $K$ be a number field and equip
$\Delta=\pi_1((\mathbb P^1\setminus\{0,1,\infty\})_{\overline K})$
with its usual outer $G_K$-action.  Among loop structures of type $(0,3)$,
meaning the three conjugacy classes
$\overline{\langle a\rangle},\overline{\langle b\rangle},
\overline{\langle (ab)^{-1}\rangle}$ for a profinite free basis $(a,b)$,
the usual loop structure is the unique one invariant under this action.
This is a precise instance of the first sentence of the conjecture in
\S34, scan 483, independent of its false normalizer explanation.
\begin{proof}
Let such a triple be invariant.  After a finite extension $L/K$, each
conjugacy class is individually fixed.  For any one of its topological
generators $z$, the normalizer of $I=\overline{\langle z\rangle}$ in
$\Pi_{U_L}$ projects onto $G_L$.  Its conjugation action on $z$ is
$z\mapsto z^{\lambda(h)}$ for some
$\lambda(h)\in\widehat{\mathbb Z}^{\times}$.
The true Galois action on $\Delta^{\mathrm{ab}}$ is scalar multiplication
by $\chi_K$, since all three actual punctures are rational.  The image of
$z$ in $\Delta^{\mathrm{ab}}$ is a primitive vector: this is evident for
$a,b$ and for $-a-b$.  Thus
$\lambda(h)[z]=\chi_K(p(h))[z]$ implies
$\lambda(h)=\chi_K(p(h))$.  In particular the geometric part of the
normalizer acts trivially.  D.3, applied over $L$, identifies $I$ with
one of the actual inertia subgroups: it is contained in one and is maximal
procyclic, being a free-basis factor or its automorphic image.
The three candidate conjugacy classes are distinct, as their abelianized
lines are distinct.  Hence they are exactly the three actual classes.
Finite base extension does not change the geometric inertia subgroups,
so the result holds over $K$.
\end{proof}

\statement{D.23}{Open}{What invariant loop data and admissible moduli splittings still require}
For an arbitrary affine hyperbolic $U/K$ over a number field, consider the
assertion that every profinite surface loop structure of its type on
$\Delta_U$ which is invariant under $\rho_U(G_K)$ equals the geometric
cuspidal loop structure.  The cyclotomic version is D.3 and the tripod case
is D.22.  We have not established this bare invariance assertion in the
remaining types, and do not identify it with the false assertion of D.2.
The first type not covered by our elementary argument is $(1,1)$: its
peripheral commutator has zero abelianization, so the primitive-vector
argument forcing the cyclotomic exponent does not apply.
Precisely the missing implication in that argument is that invariance of
an alternative peripheral subgroup forces the normalizer action required
by D.3 after an open restriction of $G_K$.

Independently, the proposed description of admissible partial splittings
of $N_{g,n}$ in \S34, folios 247--247 verso, scans 487--488, is not a
consequence of the Hom theorem or the proper section manuscript.
Once $N^{\mathrm{arith}}$ is identified with an arithmetic moduli-stack
fundamental group in a type where the congruence kernel vanishes, these
splittings concern points of that stack.  Classifying all of them as
curve-moduli points is a moduli-stack section problem; neither D.5 nor
D.7 supplies that classification.  In particular the empty list of
extra conditions at $(1,1)$ in scan 487 is not proved sufficient by the
2026 proper-curve result.  D.10 is the exact consequence obtained on the
specified high-genus modular-curve open subgroup.  In higher genus, even
the identification with the full arithmetic stack fundamental group must
retain the congruence-kernel qualification of D.17.

These Open labels state the exact assertions left unproved by this audit;
they are not claims that every formulation of inertia recognition or every
special section theorem is unknown in the literature.  Our searches and
checks used Nakamura's cyclotomic criterion, Mochizuki's relative Hom theorem,
the full Hoshi--Mochizuki monodromy theorem, and the two new proof packs.
An assertion left conditional or unverified has not been refuted merely
because our proof stops at an identified implication.

\clearpage\section{The fundamental conjecture: three different comparison problems}\label{sec:E}

\subsection{The change of definition in \S29}

We use the conventions of Part A.  In particular, $\pi_{g,n}$ is the
\emph{profinite} surface group, $P_{g,n}$ is the pure orientation-preserving
mapping class group, and $\Sigma_{g,n}$ is the closed image of the full
orientation-preserving mapping class group, with permutations of the marked
points allowed, in
$O_{g,n}=\operatorname{Out}_{\mathrm{lac}}(\pi_{g,n})$.
The suffix ``lac'' retains the conjugacy classes of peripheral procyclic
subgroups and their common cyclotomic exponent.  All normalizers, images,
quotients and commutator closures in this section are topological.  An
identification with an automorphism group is an identification of continuous
automorphisms.  Choices of a geometric marking identify fundamental groups;
changing paths conjugates the resulting actions.  The distinction between a
canonical outer action and a chosen genuine action will be retained.

\begin{center}
\begin{tabular}{L{.20\linewidth}L{.45\linewidth}L{.26\linewidth}}
Our notation & Definition & Manuscript source\\ \hline
$N^{\mathrm{norm}}_{g,n}$ &
$N_{O_{g,n}}(\Sigma_{g,n})$ & \S26, fol.~164, scan~324, (14)\\
$\Gamma^{\mathrm{norm}}_{g,n}$ &
$N^{\mathrm{norm}}_{g,n}/\Sigma_{g,n}$ & same passage; fol.~168, scan~332, (27)\\
$M^{\mathrm{arith}}_{g,n}$ & closed image of the arithmetic universal-curve
fundamental group in $\operatorname{Aut}_{\mathrm{lac}}(\pi_{g,n})$ &
\S29, fols.~187--188, scans~370--371, (9)--(13)\\
$N^{\mathrm{arith}}_{g,n}$ & exterior image of $M^{\mathrm{arith}}_{g,n}$ & same passage\\
$\Gamma^{\mathrm{arith}}_{g,n}$ &
$N^{\mathrm{arith}}_{g,n}/\Sigma_{g,n}$ & same passage\\
$G_{\mathbb Q}$ & $\operatorname{Gal}(\overline{\mathbb Q}/\mathbb Q)$ &
his $\mathbb I\Gamma_{\mathbb Q}$\\
\end{tabular}
\end{center}

The symbols $N$ and $\mathbb I\Gamma$ do not keep their definitions across
these passages.  Scan~370 announces the change to an arithmetic construction;
on scan~371 we read the instruction to denote the \emph{image} of the
arithmetic homomorphism by $M$, and the arrow
from $G_{\mathbb Q}$ to the resulting quotient is explicitly called
surjective.  The latter conjecture asks for injectivity.  We checked these
two scans, including the words surrounding the formulas.  The missing
orientation superscript in one occurrence on scan~371 is recorded in the
readings register; it does not license identifying the two normalizers.

\statement{E.1}{Proved}{The early and revised conjectures are inequivalent statements}
For every hyperbolic type $2g-2+n>0$ there are continuous homomorphisms
\[
 G_{\mathbb Q}\twoheadrightarrow\Gamma^{\mathrm{arith}}_{g,n}
 \hookrightarrow\Gamma^{\mathrm{norm}}_{g,n}.
\]
The first is surjective by definition.  Its injectivity is the revised
conjecture of \S29, fol.~188, scan~371.  The conjecture of \S26,
fol.~168, scan~332, asks that the composite be an isomorphism; it additionally
asserts that every element of the normalizer quotient is arithmetic.
The stronger assertion $N^{\mathrm{norm}}_{g,n}=O_{g,n}$ appearing in that
passage is a further assertion, not an equivalent formulation of either one.
\begin{proof}
The geometric subgroup is normal in the arithmetic fundamental group.
Its image $\Sigma_{g,n}$ is consequently normalized by the arithmetic image.
Thus $N^{\mathrm{arith}}_{g,n}\subseteq N^{\mathrm{norm}}_{g,n}$, and both
quotients are by the same closed subgroup.  The exact sequence of the
arithmetic moduli stack gives the first surjection.  Surjectivity of the
composite is exactly equality of the two images modulo $\Sigma_{g,n}$,
equivalently equality of the two groups $N$.  No assertion about the rest of
$O_{g,n}$ follows.  This also explains why the discussion of admissible
splittings on \S33, fol.~242, scan~479, cannot be read using the original
normalizer definition without an additional argument.
\end{proof}

\statement{E.2}{Known}{The revised conjecture is true for every hyperbolic type}
The natural map
\[
 G_{\mathbb Q}\xrightarrow{\ \sim\ }\Gamma^{\mathrm{arith}}_{g,n}
\]
is an isomorphism of profinite groups for every hyperbolic $(g,n)$.
It is compatible with the arithmetic maps between types wherever those
maps are defined.  This includes $(0,3)$, $(0,4)$, $(0,5)$, $(1,1)$ and
$(1,2)$, as well as proper curves of genus at least two.  It does not require
the congruence subgroup property in unknown genera.
\begin{proof}
Part D, Statement D.17, gives the precise application of
\citet[Corollary~6.4]{HoshiMochizuki2011}: the arithmetic Galois image
has trivial intersection with the geometric universal-monodromy image.
After a finite field extension admitting a curve of the specified type,
this proves injectivity on an open subgroup of $G_{\mathbb Q}$.  The remaining
kernel is a finite normal subgroup of the slim group $G_{\mathbb Q}$, and
is therefore trivial.  Statement E.1 supplies surjectivity.  The use of the
actual image $\Sigma$, rather than an unproved identification of that image
with the full profinite completion, is essential.  Naturality follows from
the arithmetic fundamental-group diagrams.  A continuous bijection between
profinite groups is a homeomorphism.
\end{proof}

For $(0,3)$, the arithmetic moduli stack with unlabelled points is
$B\mathfrak S_3$ over $\mathbb Q$; the labelled stack is $\operatorname{Spec}
\mathbb Q$.  Belyi faithfulness supplies the required injectivity here:
the three-branch-point theorem is \citet[\S4, Theorem~4]{Belyi1980},
and faithfulness is its unnumbered corollary on printed page~256.
Consequently the impression recorded on scan~479 is correct for the
\emph{revised} group.  Lead~1 is false as an identification of that group
with an independently defined $\widehat{\mathrm{GT}}$: surjectivity onto
the revised arithmetic quotient is already built into its definition.

\statement{E.3}{Known}{The abelianization of the revised group}
For every hyperbolic type,
\[
 (\Gamma^{\mathrm{arith}}_{g,n})^{\mathrm{ab}}
 =\Gamma^{\mathrm{arith}}_{g,n}/
   \overline{[\Gamma^{\mathrm{arith}}_{g,n},\Gamma^{\mathrm{arith}}_{g,n}]}
 \xrightarrow{\ \sim\ }\widehat{\mathbb Z}^{\times}
\]
is the cyclotomic character.  Thus the consequence anticipated in
\S26--\S29 (in particular scan~340) is a theorem for the revised groups.
\begin{proof}
Apply E.2 and the Kronecker--Weber theorem:
$\mathbb Q^{\mathrm{ab}}=\bigcup_m\mathbb Q(\mu_m)$.
The restriction map identifies the indicated quotient with
$\varprojlim_m(\mathbb Z/m\mathbb Z)^{\times}$.  The image of inertia
generators in the universal curve is acted on by this same character.
This argument says nothing about additional abelian quotients of a larger
normalizer group.
\end{proof}

\subsection{The genus-zero comparison, including the isolated tripodal case}

We write $\widehat F_2=\langle x,y\rangle$, $z=(xy)^{-1}$.  A profinite
Grothendieck--Teichm\"uller element is represented by $(\lambda,f)$ with
$\lambda\in\widehat{\mathbb Z}^{\times}$ and
$f\in\overline{[\widehat F_2,\widehat F_2]}$, such that
\[
 x\longmapsto x^\lambda,\qquad
 y\longmapsto f^{-1}y^\lambda f
\]
defines an automorphism and satisfies Drinfeld's symmetry, hexagon and
pentagon equations.  The automorphism requirement is part of the definition
in the profinite setting.  We denote by $\widehat{\mathrm{GT}}_0$ the group
with symmetry and hexagon only.  The pentagon is a condition in the
fundamental group of $\mathcal M_{0,5}$, not an extra relation in
$\pi_1(\mathcal M_{0,4})$ with no reference to its embeddings.  Our convention
is that $P_{0,m}=\pi_1(\mathcal M_{0,m}(\mathbb C))$.

\statement{E.4}{Known}{The exact two-level theorem}
Let $\operatorname{Out}^{\#}(\widehat P_{0,m})$ consist of outer
automorphisms which send each standard boundary generator $x_{ij}$ to a
conjugate of $x_{ij}^{\lambda}$ for a single
$\lambda\in\widehat{\mathbb Z}^{\times}$, and which commute with the
outer relabelling action of $\mathfrak S_m$.  Then
\[
 \operatorname{Out}^{\#}(\widehat P_{0,4})
       \cong\widehat{\mathrm{GT}}_0,
 \qquad
 \operatorname{Out}^{\#}(\widehat P_{0,m})
       \cong\widehat{\mathrm{GT}}\quad(m\geq5).
\]
The forgetful maps identify the latter groups for all $m\geq5$, and the
map to the four-point group is the usual injection
$\widehat{\mathrm{GT}}\hookrightarrow\widehat{\mathrm{GT}}_0$.
\begin{proof}
These are precisely the Main Theorem (a),(b) and its Corollary in
\citet{HarbaterSchneps2000}; the four-point identification is their
Theorem~2, and the five-point identification their Theorem~3.
Their Proposition~7 proves the pentagon extension criterion, and
Proposition~8 proves injectivity of forgetting.  The common exponent and
the relabelling centralizer occur in their definition of
$\operatorname{Out}^{\#}$ and cannot be dropped from the statement.
\end{proof}

There is also a version with chosen genuine braid-group automorphisms:
the objects are $\widehat B_j$ and the subgroups generated by
$\sigma_1^2,\sigma_2,\ldots,\sigma_{j-1}$, with the standard inclusions
and doubling maps.  Automorphisms act trivially on the permutation
quotients and commute with those maps.  The tower through three strands
has automorphism group $\widehat{\mathrm{GT}}_0$; through four strands,
and at every further finite truncation, it has
$\widehat{\mathrm{GT}}$.  These are the tower statements recalled and
used in \citet[\S0.3]{HarbaterSchneps2000}.  They specify a concrete
meaning of ``two levels''.  Arbitrary independent outer automorphisms of
two groups are not such a tower action.

\statement{E.5}{Known}{The intrinsic five-point and higher automorphism theorem}
For $m\geq5$,
\[
 \operatorname{Out}(\widehat P_{0,m})
       \cong\widehat{\mathrm{GT}}\times\mathfrak S_m.
\]
The symmetric factor is the geometric relabelling action; the other factor
preserves all generalized fibre subgroups.  Both factors and their
commutation are part of the theorem.
\begin{proof}
Take $(g,r)=(0,3)$ and $n=m-3\geq2$ in
\citet[Corollaries~2.6(i),(iii) and~2.8(i),(ii)]{HoshiMinamideMochizuki2022}.
Their configuration group $\Pi_n$ is $\widehat P_{0,m}$; their
$\operatorname{Out}^{gF}$ is the factor displayed here.  Their
Definition~2.7 identifies $\mathrm{GT}^{\Sigma}$ with the classical
$\widehat{\mathrm{GT}}$ when the set of primes $\Sigma$ is all primes.
In particular the result does not presuppose that a given automorphism
preserves inertia.  It reconstructs the relevant distinguished subgroups.
\end{proof}

\statement{E.6}{Proved}{A normalizer lemma for the universal curve}
Let $\pi$ be a centre-free profinite group, let
$S^!\triangleleft S\subseteq\operatorname{Out}(\pi)$ be closed subgroups,
and let $E$ be the inverse image of $S^!$ in $\operatorname{Aut}(\pi)$.
If $S\subseteq N\subseteq\operatorname{Out}(\pi)$ and $N$ normalizes
$S$ and $S^!$, conjugation induces an injection
\[
 N/S^!\hookrightarrow\operatorname{Out}(E)
\]
whose image normalizes the image of $S/S^!$.
\begin{proof}
Choose a representative $a\in\operatorname{Aut}(\pi)$ of an element of
$N$.  Conjugation by $a$ preserves $E$.  Replacing $a$ by an inner
automorphism of $\pi$ changes this automorphism of $E$ by an inner
automorphism of $E$, since $\operatorname{Inn}(\pi)\subseteq E$.
Suppose conjugation by $a$ on $E$ equals conjugation by $e\in E$.
Then $e^{-1}a$ centralizes $\operatorname{Inn}(\pi)$.  For every
$u\in\pi$ we have
$\operatorname{inn}((e^{-1}a)(u))=\operatorname{inn}(u)$; centre-freeness
gives $(e^{-1}a)(u)=u$.  Hence $a=e$, and its exterior class belongs to
$S^!$.  This proves injectivity.  The normalizer assertion follows by
conjugating the subgroup $S$ before passage to the quotient.
\end{proof}

\statement{E.7}{Proved}{The early normalizer quotient in genus zero}
With the full permutation convention of \S19 and \S26,
\[
 \Gamma^{\mathrm{norm}}_{0,3}\cong\widehat{\mathrm{GT}}_0,
 \qquad
 \Gamma^{\mathrm{norm}}_{0,n}\cong\widehat{\mathrm{GT}}
       \quad(n\geq4).
\]
The identifications respect the arithmetic embeddings after the usual
geometric markings.  The arithmetic subgroup in every case is
$\Gamma^{\mathrm{arith}}_{0,n}\cong G_{\mathbb Q}$.
\begin{proof}
First suppose $n\geq4$.  The profinite Birman sequence, and genus-zero
congruence subgroup property from Parts B and C, identify the inverse
image of $\Sigma^!_{0,n}$ in $\operatorname{Aut}(\pi_{0,n})$ with
$E=\widehat P_{0,n+1}$.  A normalizer element preserves $\Sigma^!$ because
it preserves the peripheral permutation map.  E.6 embeds
$N^{\mathrm{norm}}/\Sigma^!$ into
$\operatorname{Out}(E)=\widehat{\mathrm{GT}}\times\mathfrak S_{n+1}$.
The image of $\Sigma/\Sigma^!$ is the point stabilizer
$\mathfrak S_n\subset\mathfrak S_{n+1}$.  Its normalizer is itself:
its unique fixed letter must be preserved by any normalizing permutation.
Consequently $N^{\mathrm{norm}}/\Sigma$ embeds in
$\widehat{\mathrm{GT}}$.

Conversely, the standard genus-zero tower action of
$\widehat{\mathrm{GT}}$ preserves the kernel of
$\widehat P_{0,n+1}\to\widehat P_{0,n}$ and commutes, outerly, with
the $\mathfrak S_n$ action.  That kernel is $\pi_{0,n}$.  Restriction
therefore produces an exterior automorphism of $\pi_{0,n}$ preserving
its peripheral classes and normalizing $\Sigma_{0,n}$.  An inner change
of a representative of the automorphism of $E$ changes this restriction
by $\Sigma^!$, and hence has no effect modulo $\Sigma$.
Applying E.6 shows that this construction is inverse to the injection.

For $n=3$, $\Sigma_{0,3}\cong\mathfrak S_3$ and
$\Sigma^!_{0,3}=1$.  Every automorphism of $\mathfrak S_3$ is inner.
Thus its normalizer in $O_{0,3}$ is the product of $\mathfrak S_3$ with
its centralizer, with trivial intersection.  The centralizer acts
trivially on the three peripheral labels.  Under
$\pi_{0,3}\cong\widehat P_{0,4}$, it is exactly
$\operatorname{Out}^{\#}(\widehat P_{0,4})$: the relabelling action of
$\mathfrak S_4$ on four-point moduli has image $\mathfrak S_3$.
Now apply E.4.  Naturality of the arithmetic monodromy gives the last
assertions.
\end{proof}

Replacing the full group $\Sigma_{0,3}$ by its pure subgroup $1$ changes
the problem: the pure normalizer quotient is then the entire
label-preserving inertia automorphism group of $\widehat F_2$.
It cannot be called $\widehat{\mathrm{GT}}$ merely because the same free
group occurs.  This distinction is especially relevant to retrospective
accounts which suppress the exclamation-mark convention of \S19.

\statement{E.8}{Open}{The remaining profinite arithmetic equality}
The assertion
\[
 G_{\mathbb Q}\longrightarrow\widehat{\mathrm{GT}}
       \quad\hbox{is surjective}
\]
is the extra arithmetic assertion needed to identify the early normalizer
quotients in types $(0,4)$ and $(0,5)$ with $G_{\mathbb Q}$.
The map is a continuous closed embedding by Belyi faithfulness and the
Galois compatibility of the tower.  The isolated $(0,3)$ assertion instead
asks for surjectivity onto $\widehat{\mathrm{GT}}_0$; it contains, in
addition, the requirement that every such automorphism satisfy the
pentagon.  We do not infer either surjectivity from E.2.

Our source audit did not verify a primary proof of a \emph{strict}
profinite inclusion $\widehat{\mathrm{GT}}\subsetneq
\widehat{\mathrm{GT}}_0$.  Some historical discussions assert such
strictness; we have not substituted a rational or a finite-shadow
calculation for the missing profinite counterexample.  Accordingly we
record the isolated-level equality with its precise additional condition,
rather than declare it false on that basis.  This is a limitation of this
audit, not a claim that no relevant theorem exists.  The invalid step in
Lead~1 is already settled independently by E.1--E.2.

\subsection{The two elliptic types and the meaning of an all-genus tower}

\statement{E.9}{Known}{The two-point elliptic automorphism theorem}
There is an isomorphism
\[
 P_{1,2}\cong B_4/Z(B_4),\qquad
 \operatorname{Out}(\widehat P_{1,2})\cong\widehat{\mathrm{GT}}.
\]
Here $P_{1,2}$ fixes both marked points individually.  The assertion is
\citet[Theorem~B, Theorem~4.3 and Corollary~C]{MinamideNakamura2022}.
It does not identify the outer automorphism group of
$\widehat{\mathrm{SL}_2(\mathbb Z)}$ with
$\widehat{\mathrm{GT}}$, nor the unrestricted group with permuted
elliptic markings with its pure counterpart.  For the full group in type
$(1,2)$ there is a central involution exchanging the two punctures.
The inertia-preserving outer group and the unrestricted outer group are
respectively $\widehat{\mathrm{GT}}$ and
$C_2\times\widehat{\mathrm{GT}}$ in
\citet[Proposition~3.9]{Boggi2026Automorphisms}.

\statement{E.10}{Proved}{The one-point elliptic normalizer}
There is an isomorphism
\[
 \Gamma^{\mathrm{norm}}_{1,1}\cong\widehat{\mathrm{GT}},
\]
compatible with the arithmetic inclusion.  This statement concerns the
normalizer in $\operatorname{Out}_{\mathrm{lac}}(\pi_{1,1})$, not all
outer automorphisms of the elliptic mapping class group.
\begin{proof}
The profinite Birman sequence and genus-one congruence theorem identify
the inverse image of $\Sigma_{1,1}$ in
$\operatorname{Aut}(\pi_{1,1})$ with $E=\widehat P_{1,2}$.
E.6 and E.9 give an injection of the asserted normalizer quotient into
$\widehat{\mathrm{GT}}$.

We verify the converse, including the peripheral condition.  Put
$\zeta_j=(\sigma_1\cdots\sigma_{j-1})^j$.  The homomorphism
$\beta:B_4\to B_3$, given by
$\sigma_1,\sigma_2,\sigma_3\mapsto\sigma_1,\sigma_2,\sigma_1$,
sends $\zeta_4$ to $\zeta_3^2$.  It induces the elliptic forgetful map
\[
 E=\widehat{B_4/\langle\zeta_4\rangle}
   \longrightarrow
 \widehat{B_3/\langle\zeta_3^2\rangle}
   =\widehat{\mathrm{SL}_2(\mathbb Z)},
\]
with kernel $\pi_{1,1}$.  Compose with the quotient to
$\widehat{B_3/\langle\zeta_3\rangle}$.  Its kernel $K$ is
$C_2\mathbin{\widehat *}C_2\mathbin{\widehat *}C_2$, and it is
characteristic in $E$; these are the Cardano--Ferrari calculation and
Proposition~3.2(iii) of \citet{MinamideNakamura2022}.
The subgroup $\pi_{1,1}$ is the kernel of the unique homomorphism
$K\to C_2$ which is nonzero on each conjugacy class of involutions.
Every finite-order element of this profinite free product is conjugate
into a factor.  Thus this kernel is characteristic in $K$, hence in $E$.

Under the Birman identification its peripheral element is, up to inverse
and conjugacy, $(\sigma_1\sigma_2)^6=\zeta_3^2$ in $E$: this is the
two-chain relation, in which the boundary of a one-holed torus also
bounds the disc containing the two marked points.  The standard
Grothendieck--Teichm\"uller braid action sends $\zeta_3$ to
$\zeta_3^\lambda$ and $\zeta_4$ to $\zeta_4^\lambda$; consequently
it preserves this peripheral subgroup with exponent $\lambda$.
Restriction to the characteristic subgroup $\pi_{1,1}$ therefore
defines an element of $O_{1,1}$ normalizing $\Sigma_{1,1}$.
The ambiguity from an inner automorphism of $E$ is exactly $\Sigma_{1,1}$
after exterior restriction.  E.6 shows that restriction and conjugation
are inverse.  This proves the claim.  The full-twist formulas used here
are Drinfeld's standard braid action, also stated explicitly in
\citet[Theorem~2.1]{LochakNakamuraSchneps2026}.
\end{proof}

\statement{E.11}{Proved}{An exact reduction for the two-point normalizer}
Put $E=\widehat P_{1,3}$, and let
$\pi=\ker(E\to\widehat P_{1,2})$ be the kernel of forgetting the third
point.  Let $A\subseteq\operatorname{Out}(E)$ consist of the classes
which preserve $\pi$, whose restrictions preserve its two peripheral
conjugacy classes with a common exponent, and which normalize the
outer $\mathfrak S_2$ action permuting the first two marks.  Then
\[
 \Gamma^{\mathrm{norm}}_{1,2}\cong A/\mathfrak S_2.
\]
In particular
$G_{\mathbb Q}\cong\Gamma^{\mathrm{arith}}_{1,2}
\hookrightarrow A/\mathfrak S_2$.
\begin{proof}
Use E.6 with $S^!=\Sigma^!_{1,2}$ and $S=\Sigma_{1,2}$.
The pullback of $S^!$ in $\operatorname{Aut}(\pi)$ is $E$ by the
Birman sequence and the genus-one congruence theorem.  Conversely, a
representative $b$ of a class in $A$ restricts to an automorphism
$a$ of $\pi$.  For $e\in E$ the identity
$b(eue^{-1})=b(e)a(u)b(e)^{-1}$, $u\in\pi$, implies that
$b$ on $E\subseteq\operatorname{Aut}(\pi)$ is conjugation by $a$.
The normalizer and peripheral conditions are therefore exactly the
conditions that its outer class belong to $N^{\mathrm{norm}}_{1,2}$.
Changing $b$ by an inner automorphism changes $a$ by $\Sigma^!$.
Quotienting by the remaining permutations proves the formula.
\end{proof}

This is an exact reduction, not an identification of the unrestricted
group $A/\mathfrak S_2$ with $\widehat{\mathrm{GT}}$.
The theorem E.9 concerns the base group $P_{1,2}$; the normalizer problem
also retains its universal curve, whose group is $P_{1,3}$.
Applying E.9 to the wrong group would erase precisely this distinction.

\statement{E.12}{Known}{Refined actions with quilt and clutching data}
Let $\mathcal G_{\mathrm{NS}}$ denote the subgroup called
$\mathrm I\!\Gamma$ in \citet[Definition~1.1]{NakamuraSchneps2000},
defined by their additional relations $(\mathrm{III}')$ and $(\mathrm{IV})$,
with their Kummer cocycle $\rho_2$.  Then
\[
 G_{\mathbb Q}\hookrightarrow\mathcal G_{\mathrm{NS}}
       \hookrightarrow\widehat{\mathrm{GT}}.
\]
The first inclusion and subgroup assertion are their Theorem~1.2.
Their Theorems~1.3--1.4 give genuine actions after choosing a quilted
pants decomposition, canonical outer actions after forgetting that
choice, and compatibility with inclusions of subsurfaces cut out by
pants.  The change-of-quilt ambiguity is an explicitly described inner
automorphism.  The group $\Lambda$ of
\citet{HatcherLochakSchneps2000} is a related subgroup of the
$\lambda=1$ group, defined using the genus-one relation $(R)$; it is not
the full arithmetic Galois group.  Their Theorem~A (Lemma~7 of \S3)
proves the subgroup assertion; Theorem~B (Theorem~4 of \S3) gives the
pants-dependent genuine action and choice-independent outer action;
Theorem~C (Theorem~5 of \S3) embeds it in the compatible outer
automorphism group of their full subsurface tower.  These statements require the specified
extra compatibility data.  No equality between these named refined
subgroups and $\widehat{\mathrm{GT}}$ follows from E.9 or E.10.

\statement{E.13}{Known}{The current outer-tower results and their precise scope}
We distinguish two functors on orientable hyperbolic surfaces with
boundary and embeddings which may fill punctures:
\[
 S\longmapsto\widehat{P(S,\partial S)},\qquad
 S\longmapsto\check P(S,\partial S).
\]
Their arrows are the induced continuous homomorphisms modulo inner
automorphisms of the target.  For these specific functors,
\citet[\S2.8, Theorem~6.23 and Corollary~6.24]{Boggi2026Automorphisms}
proves
\[
 \operatorname{Aut}(\check{\mathfrak T}^{\mathrm{out}})
       \cong\widehat{\mathrm{GT}},\qquad
 \widehat{\mathrm{GT}}\hookrightarrow
       \operatorname{Aut}(\widehat{\mathfrak T}^{\mathrm{out}}).
\]
We use the explicitly dated version~6, 22 April 2026, as a preprint
theorem; we did not audit its proof in full.  Equality in the second
display remains its Conjecture~6.21.  The full profinite congruence
subgroup property would imply that equality.

These are statements about natural automorphisms of a group-valued
outer functor.  They are not, without a comparison theorem, statements
about all base-point paths, quilt moves, finite-cover maps or all
degeneration bimodules in a finer tower.  In particular they do not
silently remove the named relations in E.12.  The independent
\citet[Theorem~1.1]{LochakNakamuraSchneps2026} gives explicit actions of
all $\widehat{\mathrm{GT}}$ on the full profinite groups in types
$(g,0)$ and $(g,1)$; its introduction explicitly distinguishes these
actions from actions respecting the additional Lego rules.

The functorial actions in E.13 restrict to the universal-curve kernels
and preserve their peripheral twists.  Since genus one has the
congruence subgroup property, E.11 consequently yields an inclusion
$\widehat{\mathrm{GT}}\hookrightarrow\Gamma^{\mathrm{norm}}_{1,2}$
for that action.  The missing reverse inclusion in our reduction is the
assertion that every element of $A$ preserves the nonseparating inertia
classes in $E=\widehat P_{1,3}$.  The inertia-preserving automorphism
theorem does not itself prove this automatic-inertia assertion.

\statement{E.14}{Open}{The remaining automatic-inertia step in our elliptic reduction}
Every $a\in A$ from E.11 preserves the conjugacy classes of the closed
procyclic subgroups generated by nonseparating Dehn twists in
$\widehat P_{1,3}$.  If this holds, the inertia-preserving theorem
\citet[Theorem~A(ii)]{Boggi2026Automorphisms}, together with the
normalizer calculation for $\mathfrak S_2\subset\mathfrak S_3$, gives
$\Gamma^{\mathrm{norm}}_{1,2}\cong\widehat{\mathrm{GT}}$.
We leave this as the precise unproved step in this audit; we do not
claim that the general automatic-inertia problem has been settled.

For clarity, the comparisons obtained here are as follows.  Entries
describe the groups, not arbitrary choices of representatives of outer
actions.
\begin{center}
\begin{tabular}{L{.13\linewidth}L{.24\linewidth}L{.25\linewidth}L{.27\linewidth}}
Type & Revised arithmetic quotient & Early normalizer quotient & Related tower theorem\\ \hline
$(0,3)$ & $G_{\mathbb Q}$ & $\widehat{\mathrm{GT}}_0$ & pentagon extension required\\
$(0,4)$ & $G_{\mathbb Q}$ & $\widehat{\mathrm{GT}}$ & universal curve has type $(0,5)$\\
$(0,5)$ & $G_{\mathbb Q}$ & $\widehat{\mathrm{GT}}$ & stable genus-zero tower\\
$(1,1)$ & $G_{\mathbb Q}$ & $\widehat{\mathrm{GT}}$ & elliptic universal curve and E.10\\
$(1,2)$ & $G_{\mathbb Q}$ & $A/\mathfrak S_2$ as in E.11 & contains $\widehat{\mathrm{GT}}$ via E.13; E.14 unresolved here\\
\end{tabular}
\end{center}
The expected type-changing isomorphisms for the revised quotients are
all E.2.  For normalizer quotients, the genus-zero stable comparisons and
the comparison with $(1,1)$ above are group-theoretic theorems, whereas
equality with the arithmetic subgroup still requires E.8.  These are
three separate assertions.

\subsection{Arithmetic reconstruction uses more morphisms}

\statement{E.15}{Known}{What Pop's theorem reconstructs}
Let $U_0=\mathbb P^1\setminus\{0,1,\infty\}$ and let
$\mathcal V_{0,\mathrm{bir}}$ contain $U_0$, a cofinal system of open
subsets of $\mathcal M_{0,5}$, their inclusions and the restrictions of
the three functions $t_1,t_2,t_2-t_1$ to $U_0$ wherever defined.
Then the canonical map from $G_{\mathbb Q}$ to the compatible outer
automorphism group of the geometric fundamental-group functor of this
system is an isomorphism.  This is the application on pages~3--4 of
\citet{Pop2018Tripod}, from Theorem~2.7(1); its abelian-by-central
pro-$\ell$ counterpart is Theorem~2.6(1).
\begin{proof}
The category satisfies Hypothesis $(H)$ in Definition/Remark~2.2:
the variety has dimension greater than one, the designated functions
generate its function field over the algebraic closure, and the specified
birational system is rigid.  Theorem~2.7 identifies the automorphisms
of the tame fundamental-group functor; in characteristic zero this is
the full geometric \emph{profinite} functor.  The category has infinitely
many localization objects even though only three rational functions are
designated.  Removing those objects destroys the stated hypothesis.
Thus this theorem cannot be specialized to just
$(\mathcal M_{0,4},\mathcal M_{0,5})$ while keeping its conclusion.
\end{proof}

\statement{E.16}{Known}{The 2025 combinatorial reconstruction is not arithmetic surjectivity onto GT}
\citet[Theorem~C(ii)]{HoshiMochizukiTsujimura2025} reconstruct the
$\widehat{\mathrm{GT}}$-conjugacy class of $G_{\mathbb Q}$ as the maximal
closed subgroups satisfying their AA conditions.  The construction uses
the abstract configuration group
$\pi_1(\operatorname{Conf}_2(U_{0,\overline{\mathbb Q}}))$.
The additional COF, RGC, arithmetic Belyi compatibility and ampleness
requirements are indispensable hypotheses, defined in their
Definitions~3.3 and~5.12.  Their Theorem~A constructs field operations;
Theorem~B proves an injectivity assertion under additional hypotheses.
None states that all of $\widehat{\mathrm{GT}}$ has these properties.

We located this result through \citet[\S4]{Collas2026}, and then opened
the original, dated 14 March 2025.  A characterization of a conjugacy
class of arithmetic subgroups answers one formulation of scan~479's
request for an algebraic description.  It is not the equality in E.8,
and does not classify all admissible splittings of the moduli extensions.

\statement{E.17}{Proved}{The modular consequence of the claimed section theorem}
Assume the section theorem for $X_1(N)$ claimed in Corollary~8.5 of
the 2026 manuscript in family~019, with $g(X_1(N))\geq2$.
Let a section of the arithmetic extension for type $(1,1)$ lift to the
open subgroup belonging to the fine modular curve $Y_1(N)$.
Its composite with
$\pi_1(Y_1(N))\to\pi_1(X_1(N))$ is, up to geometric conjugacy, the
section of a rational cusp of $X_1(N)$.
\begin{proof}
The composite is a section because both maps are over $G_{\mathbb Q}$.
The asserted proper-curve section theorem makes it the section of a
rational point.  Mazur's theorem implies that all rational points of
$X_1(N)$ are cusps when its genus is positive; see
\citet[Introduction, Theorems~(7),(7')]{Mazur1977}, and the detailed
conditional argument in Part D, Statement D.10.  The level assumption
here also removes the generic elliptic stabilizer, so that the subgroup
is the fundamental group of a fine curve, not its coarse stack quotient.
\end{proof}
This does not show that the original affine section is conjugate into
a cusp decomposition group.  Nor does it show nonexistence: tangential
sections at rational cusps supply affine sections.  The unrefereed
proper-curve theorem does not treat affine curves, moduli stacks or all
splittings of $N^{\mathrm{arith}}_{1,1}\to G_{\mathbb Q}$.
The hypothesis and the status of family~019 are retained in the results
register; no unconditional section theorem is asserted here.

\statement{E.18}{Known}{Arithmetic rigidity is a different automorphism problem}
For a number field $k$, an odd prime $\ell$, and $m\geq4$,
\citet[Theorem~A]{Nakamura1994} identifies the $G_k$-compatible
automorphisms of the geometrically pro-$\ell$ arithmetic fundamental
group of $\mathcal M_{0,m}$, modulo geometric inner automorphisms, with
$\operatorname{Aut}_k(\mathcal M_{0,m})$.  The centralizer of the
geometric outer Galois image is $\mathfrak S_3$ at $m=4$, and
$\mathfrak S_m$ at $m\geq5$.  This has an \emph{odd-prime} hypothesis;
it is not a theorem that the unrestricted outer automorphism group is
finite.  The arithmetic compatibility condition cannot be deleted to
deduce E.8.  His Theorem~B is the corresponding Lie-algebra centralizer
statement for $m\geq5$.  This is the precise rigidity result used here,
rather than an unspecified assertion of ``Galois rigidity''.

\subsection{The rational comparison and the maps to the pro-\texorpdfstring{$\ell$}{l} tower}

For a field $k$ of characteristic zero, unipotent completion means the
universal pro-unipotent affine $k$-group receiving the discrete group,
or its continuous $\ell$-adic variant when $k=\mathbb Q_\ell$.
It is not the discrete fundamental group retained by arbitrary locally
constant $D(\mathbb Q)$-coefficients, and it is not the entire
pro-algebraic completion.  In the braid tower we complete the pure braid
groups and their path torsors, retaining the permutations and the
cabling, deletion and reassociation operations.  Completing the whole
mapping class group indiscriminately would define a different object.
The grading below is positive word weight; motivic cohomological weight
is its negative double.

Let $\mathfrak t_j$ be the infinitesimal braid Lie algebra on the
$t_{ab}=t_{ba}$, with
$[t_{ab},t_{cd}]=0$ for distinct indices and
$[t_{ab},t_{ac}+t_{bc}]=0$.  Let
$W=\bigoplus_{w\geq1}W_w\subseteq\operatorname{Lie}_{\mathbb Q}
\langle x,y\rangle$ be the solution space
\begin{align*}
 \psi(x,y)+\psi(y,x)&=0,\\
 \psi(x,y)+\psi(y,-x-y)+\psi(-x-y,x)&=0,\\
 \psi(t_{12},t_{23}+t_{24})+\psi(t_{13}+t_{23},t_{34})
 &=\psi(t_{23},t_{34})+\psi(t_{12}+t_{13},t_{24}+t_{34})
    +\psi(t_{12},t_{23}).
\end{align*}
The third equation lies in $\mathfrak t_4$.  The bracket is
\[
 \{\psi,\varphi\}=D_\psi\varphi-D_\varphi\psi+[\psi,\varphi],
 \qquad D_\psi(x)=0,\quad D_\psi(y)=[y,\psi].
\]
The ordinary commutator on the free Lie algebra is not this bracket.

\statement{E.19}{Known}{The precise rational source, target and established injection}
The weight-completed Lie algebra of $\mathrm{GRT}_1$ is
$\widehat W=\prod_w W_w$ with the Ihara bracket, and
$\mathrm{GRT}=\mathbb G_m\ltimes\mathrm{GRT}_1$.
It is the graded automorphism group for the completed parenthesized
chord tower; the corresponding filtered braid tower has automorphism
group $\mathrm{GT}(k)$.  A choice of a $k$-associator identifies these
two groups, noncanonically.  The symmetry and hexagon occur at the
four-point level and the pentagon at the five-point level; the same
operations determine the action at every genus-zero level.

For the de Rham fibre functor on $\mathrm{MT}(\mathbb Z)$, let
$G_{\mathrm{MT}}=\mathbb G_m\ltimes U_{\mathrm{MT}}$ and let
$\mathfrak u_{\mathrm{MT}}^{\mathrm{gr}}$ denote its weight-graded
unipotent Lie algebra, reindexed by positive word weight.  There is an
injective graded Lie homomorphism
\[
 j:\mathfrak u_{\mathrm{MT}}^{\mathrm{gr}}
   \cong\operatorname{Lie}_{\mathbb Q}\langle e_3,e_5,e_7,\ldots\rangle
   \hookrightarrow W.
\]
It integrates after weight completion to a monomorphism of
pro-unipotent groups.  The generators themselves are not canonical.
\begin{proof}
Drinfeld's definitions in \S\S5--6 and Propositions~5.1 and~5.5
give the associator torsor and the two actions
\citep{Drinfeld1990}.  The structure-preserving categorical formulation
is \citet[Proposition~3.4, Theorems~1--3, Corollary~4.1]{BarNatan1998};
its operations must be retained when invoking that formulation.
Differentiating gives exactly the three equations displayed above.
\citet[Theorem~1.1 and its stated consequence, Conjecture~1]{Brown2012}
proves faithfulness of the mixed Tate action.  The containment in the
Grothendieck--Teichm\"uller solution Lie algebra is stated explicitly in
\citet[\S1.4.2]{Brown2021Depth}, using Brown's theorem and Drinfeld's
equations.  The free odd-weight source is the structure of
$\mathrm{MT}(\mathbb Z)$ recalled in \citet[\S1.1]{Brown2012}.
In characteristic zero, exponential and logarithm give the equivalence
between finite-dimensional nilpotent Lie algebras and unipotent groups;
passing through finite weight quotients gives the completed assertion.
\end{proof}

\statement{E.20}{Proved}{The exact conditional rational counterpart}
Assume the all-weight bound of family~008, Corollary~5.4:
\[
 \dim_{\mathbb Q} W_w\leq
 \dim_{\mathbb Q}\operatorname{Lie}\langle e_3,e_5,\ldots\rangle_w
 \quad\hbox{for every }w.
\]
Then the map $j$ of E.19 is an isomorphism in every weight and after
weight completion.  It gives
$U_{\mathrm{MT}}\cong\mathrm{GRT}_1$ and, with the matching Tate
character and grading, $G_{\mathrm{MT}}\cong\mathrm{GRT}$.
Via a chosen associator this is the analogous statement for
$\mathrm{GT}$ on the unipotent genus-zero braid tower.
\begin{proof}
The injection $j_w$ has a source of the dimension on the right.
The assumed opposite inequality forces equality and surjectivity.
Every weight piece is finite dimensional, so the product of these
isomorphisms is an isomorphism for the complete weight filtration;
the bracket in each weight involves finitely many lower weights.
Exponentiation on the finite nilpotent quotients and inverse limit
give the group isomorphism.  The scalar actions agree by their weight
definitions (with the inverse character if the negative motivic-degree
convention is used), so they extend over $\mathbb G_m$.
\end{proof}
The conclusion is \emph{claimed (OpenAI, 2026, unrefereed)} when
Corollary~5.4 is supplied by that manuscript.  The implication just
proved is unconditional logic; its all-weight premise is not certified
by our bounded calculation or by reading a comparator statement.

\statement{E.21}{Known}{The realization chain and what its arrows do}
Fix a prime $\ell$ and compatible tangential markings.  Write
$\mathrm{GT}_{\ell}$ for the group of invertible pairs satisfying the
pro-$\ell$ versions of Drinfeld's equations, with
$\lambda\in\mathbb Z_\ell^{\times}$.  This is generally not a
pro-$\ell$ group, because its cyclotomic quotient need not be one.
There are natural maps
\[
 G_{\mathbb Q}\hookrightarrow\widehat{\mathrm{GT}}
   \longrightarrow\mathrm{GT}_{\ell}
   \hookrightarrow\mathrm{GT}(\mathbb Q_\ell).
\]
The first arrow is a continuous closed embedding.  The second is
continuous and is not injective; we assert no surjectivity.
The third is injective and continuous for the topology on finite weight
quotients.  Neither its image nor the image of $G_{\mathbb Q}$ is
surjective, topologically dense or open in
$\mathrm{GT}(\mathbb Q_\ell)$ with this inverse-limit topology.
Surjectivity or openness of the arithmetic image \emph{inside}
$\mathrm{GT}_{\ell}$ is not concluded here.
\begin{proof}
Drinfeld constructs the completion comparison maps in \S4
\citep{Drinfeld1990}; the continuous pro-$\ell$ to unipotent
comparison is also \citet[Theorem~A.6 and Corollary~A.10]{HainMatsumoto2003}.
Faithful Magnus expansion of the free pro-$\ell$ group gives
injectivity of its map to the characteristic-zero completed free
unipotent group, hence injectivity of the pair map.
The arithmetic pro-$\ell$ action factors through the Galois group
$G_{\mathbb Q,\{\ell\}}$ of extensions unramified outside $\ell$;
this is the factorization recalled and used in
\citet[Introduction and \S8.4]{HainMatsumoto2003}.
Nontrivial inertia at a prime different from $\ell$ therefore lies in
the kernel of the composite.  Belyi injectivity in the first arrow
forces a nontrivial kernel in the second.

The pro-$\ell$ solution group is compact.  Its scalar image lies in
$\mathbb Z_\ell^{\times}$, a proper closed subset of
$\mathbb Q_\ell^{\times}$, proving nonsurjectivity and failure of
density in the full rational group.  Its image is compact and hence
closed.  It cannot be open: any neighbourhood of the identity in the
inverse-limit weight topology contains the kernel of a sufficiently
large finite weight projection.  By Brown's independent odd-weight
elements, such a kernel contains a closed one-parameter subgroup
isomorphic to the additive group $\mathbb Q_\ell$ in an arbitrarily
large weight, and is not compact.  A compact subgroup cannot contain
that closed noncompact subgroup.  The same argument applies to the
compact arithmetic image.
\end{proof}

Hain--Matsumoto is a statement about weighted completion and the
\emph{image} of an arithmetic action, not faithfulness of
$G_{\mathbb Q}\to\mathrm{GT}_{\ell}$.  More precisely their
Theorem~7.3 computes the free graded Lie algebra of the weighted
completion of $G_{\mathbb Q,\{\ell\}}$ with respect to the cyclotomic
character; this initially includes weight one.  Their Theorem~8.4
identifies the graded Lie image with the rationalized graded image of
the Galois action.  Lemma~8.5 eliminates weight one in the outer action;
Remark~8.6 identifies the nonzero odd abelian classes by Soul\'e's
characters.  Their Theorem~9.8 gives the constrained mixed Tate
factorization.  Brown's faithfulness then supplies the free motivic
image, while E.20, conditionally on the upper bound, says that this
image exhausts $W\otimes\mathbb Q_\ell$.  The resulting density is
\emph{Zariski density in each finite-dimensional algebraic quotient}.
It is compatible with the preceding failure of topological density in
the full group of $\mathbb Q_\ell$-points.  No argument here divides by
$2$ or $3$ in a profinite inertia group; the all-prime Hain--Matsumoto
construction and the separate characteristic-two lattice calculation
must not be confused with Nakamura's odd-prime theorem E.18.

\statement{E.22}{Proved}{The needed profinite analogue of the upper bound}
For the closed embedding $\iota:G_{\mathbb Q}\hookrightarrow
\widehat{\mathrm{GT}}$, the following are equivalent:
\begin{enumerate}
\item $\iota$ is surjective;
\item its image is topologically dense;
\item for every open normal subgroup $V$ of
$\widehat{\mathrm{GT}}$, the map
$G_{\mathbb Q}\to\widehat{\mathrm{GT}}/V$ is surjective.
\end{enumerate}
This finite-quotient exhaustion, not a rational dimension inequality,
is the required profinite counterpart of Corollary~5.4.
\begin{proof}
A compact image in a Hausdorff profinite group is closed.  Density is
therefore equivalent to equality.  Density implies surjectivity on
every finite discrete quotient.  Conversely, if a closed subgroup $H$
is proper, a point outside $H$ has an open neighbourhood disjoint from
$H$; shrinking to a coset of an open normal subgroup $V$ gives a coset
missed by $H$ in the finite quotient.  This contradicts the third
condition.
\end{proof}

\statement{E.23}{Refuted}{Equality of unipotent completions does not imply profinite equality}
The inference ``a subgroup induces an isomorphism of characteristic-zero
unipotent completions, therefore it is the whole profinite group'' is
false, already for $1\subset C_2$.
\begin{proof}
Every finite-order element of a characteristic-zero unipotent group is
the identity.  Indeed, write it as $\exp X$ in a faithful upper
unitriangular representation; $(\exp X)^r=1$ implies $rX=0$, hence
$X=0$.  Therefore every homomorphism $C_2\to U$ is trivial for every
unipotent $U$, and the unipotent completions of $1$ and $C_2$ are both
trivial.  Their profinite completions are themselves, and differ.
The supplied program checks the two-element multiplication table and
its associativity; the torsion-free argument is the proof of the
completion assertion.  This counterexample does not say that the
arithmetic inclusion is proper.  It refutes only the proposed transfer
principle.
\end{proof}

\statement{E.24}{Proved}{Conditional abelianizations of GRT}
Assume the freeness assertion of family~008, Theorem~1.1.  Then
\[
 W^{\mathrm{ab}}\cong\bigoplus_{k\geq1}\mathbb Q e_{2k+1},
 \qquad
 (\mathrm{GRT}_1)^{\mathrm{ab}}
       \cong\prod_{k\geq1}\mathbb G_a(2k+1),
 \qquad
 (\mathrm{GRT})^{\mathrm{ab}}\cong\mathbb G_m.
\]
The middle and last abelianizations are in the category of affine
pro-algebraic groups, equivalently by closed commutator subgroups.
The last formula, with the scalar character, is the rational analogue
of Grothendieck's cyclotomic abelianization statement; the middle
formula is not.
\begin{proof}
The abelianization of a free graded Lie algebra is its generating
space.  Passing to the weight completion replaces the direct sum by
the product, and exponentiation identifies an abelian unipotent group
with its additive Lie algebra.  In the full semidirect product, every
positive-weight abelian piece has commutator action
$(h-1)v=(2^{\mathrm{wt}(v)}-1)v$ for the scalar $h=2$.
The coefficient is nonzero and invertible over $\mathbb Q$.
Thus all unipotent weight pieces vanish in the abelianization.
This argument can be applied successively in every finite weight
quotient; inverse limit gives the closed-commutator assertion.  The
remaining quotient is the already abelian group $\mathbb G_m$.
No assertion about the abstract, nonclosed derived subgroup of an
infinite product is being made.
\end{proof}

\statement{E.25}{Known}{The arithmetic meaning of the odd generators}
On the motivic side the odd generating lines are dual to
\[
 \operatorname{Ext}^1_{\mathrm{MT}(\mathbb Z)}
       (\mathbb Q(0),\mathbb Q(2k+1))
       \cong K_{4k+1}(\mathbb Z)\otimes\mathbb Q,
       \qquad k\geq1.
\]
Each line has dimension one.  Their $\ell$-adic realizations are
measured by the nonzero Soul\'e characters in weight $2k+1$.
In particular they are abelian classes in the unipotent or
cyclotomic-kernel Lie algebra, not extra ordinary abelian characters
of $G_{\mathbb Q}$ contrary to Kronecker--Weber.
\begin{proof}
This is the mixed Tate description recalled in
\citet[\S1.1]{Brown2012}, together with
\citet[Theorem~7.3 and Remark~8.6]{HainMatsumoto2003}.
The $G_{\mathbb Q}$-equivariant abelianization in the latter is
$\bigoplus_{k\geq1}\mathbb Q_\ell(2k+1)$.
Taking the full semidirect product with the cyclotomic torus kills
these positive-weight coinvariants, exactly as in E.24.
\end{proof}

\subsection{Audit of the new all-weight upper bound}

\statement{E.26}{Open}{The imported 2026 claim and the extent of our verification}
Family~008, \citet[Theorem~1.1 and Corollary~5.4]{OpenAI2026DD},
claims that $W$ is free on one generator in every odd weight at least
three.  We read \S\S2--5, including the all-weight reduction, and
verified the implication E.20 using the established Brown lower bound.
The following audit records the actual checks and their limits.
We have not independently certified every all-weight rewriting and
deletion argument, so our mathematical use of the upper bound remains
conditional.  ``Open'' here records that outstanding verification in
this deliverable; it does not erase the announced proofs described
below.

\paragraph{Integral reduction (\S2).}
The use of $R=\mathbb Z_{(2)}$ is appropriate.  In a fixed weight the
free ordinary Lie lattice is finite free by its integral Hall--Lyndon
basis.  The solution lattice
$L_w=W_w\cap\operatorname{Lie}_R\langle x,y\rangle_w$ is saturated:
if $2a\in L_w$, then $a\in W_w$ because the equations are rational
linear, and $a$ already lies in the ambient lattice.  Hence its reduction
embeds in the ambient reduction, and
$\dim_{\mathbb F_2}\overline L_w=\dim_{\mathbb Q}W_w$.
We followed this proof of Proposition~2.5.  It does \emph{not} assert
that $\overline L_w$ is the entire mod-two solution space.
The degree-one pentagon leaves $-\alpha t_{12}-\beta t_{34}$,
and the degree-two three-cycle gives $3[x,y]$; thus Lemma~2.1 indeed
gives $W_1=W_2=0$.

Lemma~2.3 uses the standard free-fibre decomposition of the
infinitesimal braid Lie algebra; its explicit derivations give the
required integral torsion-freeness of the targets.  We followed the
presentation comparison.  For Lemma~2.4 we followed the reduction of
the pentagon in the free fibre, the centralizer argument, and the
linearization giving the special-derivation identity and its transverse
derivative.  Its division by the nonzero integer $w+1$ takes place
over $\mathbb Q$ \emph{before} reducing the saturated lattice; it is
not a division by an even integer in $\mathbb F_2$.  This distinction
is necessary for the stated proof to make sense.

\paragraph{Filtered comparison (\S3).}
The substitution
$A=x^2$, $C=[x,y]$, $B=y$ is made in the associative algebra in
characteristic two, with the number of occurrences of $B$ as filtration.
The prefix words $xx,xy,y$ give the asserted injective associative
encoding.  The source and target commutator tables have a common
count-preserving part: correction terms raise the count in the braid
source and lower it in the cyclotomic arrangement target.  The ordering
rule replaces a base letter followed by a fibre letter by the opposite
order plus its commutator, decreasing the stated inversion measure.
The proof needs independent normal words in the \emph{source}; mere
spanning in the target is enough.  We checked this direction of the
argument and selected entries of the tables (3.6)--(3.12).  We did not
rederive every entry or mechanize the complete induction.  Proposition~3.2
is therefore still an imported all-weight step, not a consequence of
our finite computation.

\paragraph{Deletion representation and leading-term injectivity (\S4).}
The nine type-$B_3$ hyperplanes are first defined over $\mathbb Q$.
The plus and minus hyperplanes remain distinct symbols after mod-two
reduction; merging them geometrically in characteristic two would change
the arrangement Lie algebra.  The deletion operator at a letter has
coefficient equal to the valuation on its right edge minus the
valuation on its left edge, with equal-letter differences assigned zero.
We followed Proposition~4.1's logarithmic flatness and residue argument
over the integers.  Our independent implementation of those edge
valuations tests all 13 codimension-two flats, every incident distinguished
hyperplane, and every word of length at most five in the three-letter
alphabet.  All 13,104 integer commutator instances vanish.
This checks signs, hyperplane labels and the repeated-letter convention
in that range.  It does not prove the representation in unbounded length.

We followed how Lemma~4.3 uses the minimal number of $A$'s, terminal
deletion and reversal to constrain a leading polynomial after $A=0$.
Proposition~4.4 then uses the reduced special identity to compare
coefficients, separating the cases of at least two $A/C$ letters and
the small boundary cases.  In particular the $r=2$ expression
$ABB+BBA$ is not the uniform coefficient sum; the proof cannot discard
that case by an informal characteristic-two symmetry argument.
We checked the displayed reductions and case division, but did not
independently reconstruct the general coefficient-exchange induction.

\paragraph{The image bound (\S5).}
The even-index bound is an application of free-Lie elimination.
Lemma~5.2 uses the Ore-extension calculation with
$D(B)=C$, $D(C)=0$, $D^2=0$, giving invariance under
$B\mapsto B+s(C)$.  In Proposition~5.3 the Vandermonde variables
are independent indeterminates; the argument takes place over a
polynomial ring, not by evaluating a supposedly invertible matrix on
the two elements of $\mathbb F_2$.  Divisibility by the successive
differences in the encoding excludes the index-zero elimination
generator.  Lazard elimination then leaves the ordinary free Lie
algebra generated by
$g_k=\operatorname{ad}_C^k(B)$ for $k\geq1$, whose weights are
$2k+1$.  We followed the polynomial divisibility and ordinary-Lie
retraction steps.  The distinction between ordinary Lie elements and
all primitive elements in characteristic two is consistently necessary.
Combining this image bound with the injective projection and the finite
filtration yields Corollary~5.4 exactly as stated.  Proposition~5.5's
leading-bracket calculation is not needed for E.20, since Brown already
provides the compatible Lie injection.

\paragraph{What this audit establishes.}
The integral dimension transfer, low-weight vanishing, conditional
dimension squeeze, completion passage, abelianization calculation and
bounded deletion check have been independently checked here.  The
unbounded commutator-table/normal-form comparison of \S3 and the
coefficient-detection induction of \S4 remain the principal steps not
fully rederived by us.  We did not use the alternative lower-bound
construction in \S\S6--8, its regularized-transport proposition, or
the graph-complex corollary.  They are claimed, steps not checked by us.
The program additionally computes, from the free-Lie generating series,
the predicted weight dimensions through 30.  Those numbers describe
the proposed bound; they are not a computation of $\dim W_w$.

\statement{E.27}{Refuted}{The attached comparator is not a checked Lean proof}
The claim that the attached file
\texttt{lean/DeligneDrinfeld.lean} itself provides a kernel-checked
proof of Theorem~1.1 is false.
\begin{proof}
We read the supplied file.  It defines the solution equations, Ihara
bracket and the claimed graded and completed statement, and ends with
\texttt{theorem main : MainStatement := by sorry}.
The scope note describes a formalization elsewhere, but that description
does not remove the placeholder in this comparator.  We did not fetch,
build or audit a separate Lean proof repository.  Thus this attachment
confirms the shape of the target statement only.  This is a refutation
of a certification claim, not of the mathematical theorem.
\end{proof}

\statement{E.28}{Known}{The status search finds subsequent proof claims}
As of our search on 11 November 2026,
\citet[Theorem~1.1]{Willwacher2026After} presents a simplified proof
of the Deligne--Drinfeld conclusion following family~008; the opened
version is arXiv:2610.15592v4, dated 28 October 2026.
We read its principal statements and reduction, not a complete
independent proof audit.  It is an additional primary preprint source,
so it would be inaccurate to report that no subsequent mathematical
check or proof treatment exists.  Its existence does not turn the
attached comparator into a checked formal proof, and it does not
address the profinite equality E.8.  Our conditional labels for use of
the supplied upper bound remain explicit.

\subsection{Consequences for the requested categorical formulation}

The group acting on a completed degeneration category in Part F must
be specified by its retained operations.  For the marked genus-zero
system with common inertia exponent, relabellings, hole-filling and
clutching, its restriction is the group in E.4.  A proof that these
restrictions determine and lift category automorphisms is additional
work; it is not part of the wording of the Harbater--Schneps theorem.
Part F supplies that comparison.  Once it is supplied, replacing each
pure braid group and tangential path torsor by its unipotent completion
and retaining cabling and reassociation gives the filtered
$\mathrm{GT}$-tower, or, after an associator, the graded
$\mathrm{GRT}$-tower in E.19.  Its motivic arithmetic equality is
precisely the conditional theorem E.20.  The revised profinite
arithmetic quotient is already $G_{\mathbb Q}$ by E.2; equality with
the full geometric tower automorphism group still asks E.8.

Neither the two excluded low-dimensional elliptic and rational
nonhyperbolic types nor orientation-reversing mapping classes can be
inserted into these assertions without changing the hypotheses.
In particular, quotienting by orientation reversal would kill complex
conjugation in a real model.  The orientation-positive geometric
subgroup in E.1--E.2 is therefore essential.  Stack stabilizers at
elliptic curves with extra automorphisms, including those of orders
$4$ and $6$, are retained throughout; passing to a coarse $j$-line
would not preserve the fundamental-group extensions used above.



\clearpage\section{F. One stratified object, its completion, and its automorphisms}
\subsection{F.1. Which category is being completed}

\statement{F.1}{Known}{The input category, with its variance}
Let $n\geq3$, $G_n=P_{0,n}$, and let $\mathcal C_n$ be the set of isotopy classes of essential multicurves on the labelled punctured sphere, including the empty multicurve. For $\sigma\in\mathcal C_n$ put $\Delta_\sigma=\langle t_c:c\in\sigma\rangle\simeq\mathbb Z^{|\sigma|}$. The exit category is
\[
 \operatorname{Ob}C_n=\mathcal C_n,\qquad
 \Hom_{C_n}(\sigma,\tau)=
 \{g\in G_n:g\sigma\supseteq\tau\}/\Delta_\sigma,
 \qquad [h]\circ[g]=[hg].                                      \tag{F.1}
\]
Its nerve is the exit-path $\infty$-category of the boundary-stratified $\overline{\mathscr M}_{0,n}(\mathbb C)$. Put $T_n=C_n^{\op}$. Thus arrows in $C_n$ smooth nodes, arrows in $T_n$ degenerate, and constructible coefficients are covariant functors on $C_n$, or presheaves on $T_n$.
This is \citet[Construction 3.18, Proposition 3.23, Corollary 6.6, and Theorem 6.7]{AF_Jansen2023}, with the pure labelled genus-zero convention of A.
\begin{proof}[Verification of the input used]
We use Jansen's identification, not a new proof of the topology of the stratification. Its group-theoretic composition is well-defined: if $\tau\subseteq g\sigma$, then $g^{-1}\Delta_\tau g\subseteq\Delta_\sigma$. This is exactly the identity needed to replace $h$ by $hd$ and $g$ by $ge$ without changing the composite coset. The stabilizer quotient is
\[
\Aut_{C_n}(\sigma)=\Stab_{G_n}(\sigma)/\Delta_\sigma
 \simeq\prod_v P_{0,\operatorname{val}(v)}.
\]
There are finitely many orbits, indexed by stable labelled trees. We choose one representative $\sigma_s$ per tree $s$ and henceforth work with this finite-object skeleton $C_n^0$. This changes the category by equivalence, not by groupoid completion. We retain its actual composition maps. For $n=4$ there are four objects and for $n=5$ there are twenty-six, with the Hom-sets specified in STRATA; the combinatorial counts and forgetting map are independently checked by the accompanying program.
\end{proof}

\statement{F.2}{Proved}{A strict, universal finite-category completion}
For a category $C$ with a finite specified object set, a finite congruence $R$ is an equivalence relation of finite index on each $\Hom_C(x,y)$ that respects identities and composition. Let $C/R$ retain the same objects and the quotient Hom-sets. The poset of these congruences, ordered by refinement, is cofiltered, and
\[
 C^{\mathrm{fin}}=(C/R)_R\in\operatorname{Pro}(\mathrm{FinCat}),
 \qquad C/R'\longrightarrow C/R\quad(R'\subseteq R)             \tag{F.2}
\]
is our strict profinite completion. Every transition functor is the identity on objects and surjective on each Hom-set. Its materialization $|C^{\mathrm{fin}}|=\varprojlim_R C/R$ is a category object in profinite sets with a finite discrete object set. Continuous functors from it to a finite category are precisely the functors represented by the pro-object. Continuous natural isomorphisms are retained, and can be computed at this level. In particular their automorphism groups agree.
\begin{proof}
Intersection of two finite congruences is again a finite congruence: its classes inject into the product of the two finite quotient sets. Thus the system and its transitions are explicit. Composition on the inverse-limit Hom-sets is the inverse limit of the finite composition maps, and is continuous and associative. A functor from $C$ to a finite category has a finite kernel congruence, including when different objects have the same image, so factors through a quotient in (F.2).

Conversely a continuous map from a compact profinite Hom-set to a finite set factors through finitely many quotient coordinates. There are only finitely many pairs of objects. Taking a common refinement makes all Hom-maps factor at one stage. Because the projections are surjective on each Hom-set, identities and every composable pair lift from that stage; the functorial equations therefore hold already there. This proves the functor assertion. Apply the same argument to the finitely many components and naturality equations of a natural transformation. Equivalently,
\[
 \Hom_{\operatorname{Pro}}((C_i),(D_j))
   =\varprojlim_j\varinjlim_i\Hom(C_i,D_j)
\]
with the corresponding groupoid of functors and natural isomorphisms. Invertibility is imposed by an inverse with its two natural isomorphisms. Thus replacing this strict pro-system by its continuous materialization does not erase transitions; their equivalence follows from the stated finite-object and quotient hypotheses. It would not follow for an arbitrary pro-system presented only in a homotopy category.
\end{proof}

\statement{F.3}{Refuted}{Arbitrary natural transformations may not be inverted}
It is false that quotienting functors by arbitrary natural transformations is harmless for the completed exit category. The appropriate equivalence is natural \emph{isomorphism}, not a zigzag of transformations whose components need not be invertible.
\begin{proof}
Let $[1]$ have objects $0,1$ and its single nonidentity arrow $0\to1$. The functors $[1]\to\{*\}$ and $\{*\}\to[1]$ choosing $1$ have composite naturally related to the identity by the arrow $0\to1$. They become inverse after arbitrary natural transformations are inverted. They are not equivalences of categories, since $0$ and $1$ are not isomorphic. Both categories are already finite, hence already complete under (F.2). Moreover $\Fun([1],\mathrm{Set})$ has three subobjects of its terminal object, whereas $\mathrm{Set}$ has two; the associated coefficient categories are different. The program enumerates these subobjects and checks the transformation. This gives a finite counterexample to the proposed inference. The CATHOT warning about stagewise groupoidification is therefore respected, not applied as an identification of the exit category itself with a groupoid.
\end{proof}

\statement{F.4}{Proved}{Finite representations, split representations, and all continuous representations}
For (F.2), finite-valued continuous functors satisfy
\[
 \Fun_{\cts}(|C^{\mathrm{fin}}|,\mathrm{FinSet})
   \simeq\varinjlim_R\Fun(C/R,\mathrm{FinSet}).                 \tag{F.3}
\]
For arbitrary sets, the analogous colimit denotes representations that factor through \emph{one} finite stage. It need not contain every continuous representation. For a profinite group $H$, all discrete continuous $H$-sets form the filtered-colimit completion of the finite continuous $H$-sets. They are not the actions of the abstract group with its topology forgotten.
\begin{proof}
The first assertion is F.2 applied to the finite category of maps among the finitely many finite values of a given functor; it also treats its natural transformations. For the distinction, take $H=\widehat{\mathbb Z}$ and
\[
 S=\coprod_{m\geq1}\mathbb Z/m!\mathbb Z.
\]
Each orbit has open stabilizer, so this is a continuous action on a discrete set. A finite stage $\mathbb Z/N\mathbb Z$ cannot act through the given action on all summands, since that would force every $m!$ to divide $N$. On the other hand $S$ is the filtered union of its finite invariant subsets. The abstract regular action of $\widehat{\mathbb Z}$ on itself with the discrete topology is not continuous, since its point stabilizer is not open. This proves both distinctions. Finite representations in this paper always mean (F.3).

This verifies the precise relevant content of the attached GALOIS F.4: \citet[notation preceding Definition 2.10, Theorem 3.13, and Corollary 3.20]{AF_Hoyois2018} distinguishes the pro-object functor category from its classifying topos and calls the stagewise objects \emph{split}. The attached statement includes that distinction. We do not refute it; we refute the unqualified inference that every continuous set action factors through one finite quotient. No Kan--Thurston or rational-local-system claim is needed here.
\end{proof}

\subsection{F.2. Transporters, closed stabilizers, and an explicit model}

\statement{F.5}{Known}{The profinite curve complex and multitwist stabilizers}
For genus zero and every $n\geq3$, let $\widehat G_n$ be the full profinite completion. It equals the procongruence completion. The profinite curve complex is defined by embedding an unoriented simple curve in the space of conjugacy classes $\{\gamma,\gamma^{-1}\}$ in $\widehat\Pi_{0,n}$, taking the closure of simple curves, and then taking the closure of the discrete $r$-multicurves in its finite-subset space, for each $r$. Include the empty simplex. It is also the inverse limit of the finite orbit complexes over congruence levels. For a discrete multicurve $\sigma$, its profinite stabilizer is $\overline{H_\sigma}$, where $H_\sigma=\Stab_{G_n}(\sigma)$, and
\[
1\longrightarrow\overline\Delta_\sigma\simeq\widehat{\mathbb Z}^{|\sigma|}
 \longrightarrow\overline H_\sigma
 \longrightarrow\prod_v\widehat P_{0,\operatorname{val}(v)}\longrightarrow1. \tag{F.4}
\]
Furthermore $N_{\widehat G_n}(\overline\Delta_\sigma)=\overline H_\sigma$.
The sources are \citet[Definitions 3.3--3.4, Theorems 3.5--3.6, Theorem 4.2, and Corollary 4.3]{Boggi2020BC}, with the genus-zero congruence theorem. The trivial empty-multicurve case is interpreted separately.
\begin{proof}[Scope and hypotheses]
Boggi's statements are initially about the procongruence group, with a general class of finite quotients. We take the class of all finite groups. Part B verifies equality with the full profinite topology in genus zero, so the substitution here is legitimate. Pure mapping classes on a labelled sphere cannot permute the edges or reverse the sides of a fixed multicurve, since each edge has a distinct labelled split. Thus the finite permutation group in the general stabilizer theorem is trivial, giving (F.4). The multitwist normalizer theorem gives the last equality. Distinct disjoint twists are independent, so no prime, including $2$ or $3$, is excluded from $\widehat{\mathbb Z}^{|\sigma|}$. This verifies the precise results needed below; it does not assume that an arbitrary subgroup of a residually finite group inherits its full profinite topology.
\end{proof}

\statement{F.6}{Proved}{The completion of a stabilizer really is the stabilizer}
For the groups in F.5, the natural map $\widehat H_\sigma\to\overline H_\sigma$ is an isomorphism. Both $H_\sigma$ and $\Delta_\sigma$ are closed in the profinite topology of $G_n$, and $C_n^0\to|(C_n^0)^{\mathrm{fin}}|$ will be faithful.
\begin{proof}
Cutting the sphere along $\sigma$ gives the discrete exact sequence
\[
1\to\Delta_\sigma\to H_\sigma\to
 Q_\sigma:=\prod_vP_{0,\operatorname{val}(v)}\to1.
\]
Its completion is right exact. The natural homomorphism $\widehat\Delta_\sigma\to\widehat H_\sigma$ is injective because its composite into $\widehat G_n$ is the injection in (F.4). The kernel of $\widehat H_\sigma\to\widehat Q_\sigma$ is exactly the closed image of $\widehat\Delta_\sigma$: quotienting by it has the universal finite quotients of $H_\sigma/\Delta_\sigma$. Comparison with (F.4), and equality of the component completions with their procongruence groups, gives an isomorphism of both ends. A direct kernel-and-surjectivity argument, or the short five lemma, gives the asserted middle isomorphism.

The discrete curve complex injects into the profinite one by surface-group conjugacy separability, used in its construction. If $g\in G_n\cap\overline H_\sigma$, then $g\sigma=\sigma$ in the profinite complex, hence already in the discrete complex; thus $g\in H_\sigma$. This proves closedness of $H_\sigma$. Within $H_\sigma$, the subgroup $\Delta_\sigma$ is the kernel of the map to the residually finite group $Q_\sigma$, so is closed there. The full induced topology on $H_\sigma$ just proved makes it closed also in $G_n$. Distinct cosets modulo $\Delta_\sigma$ can therefore be separated in a finite quotient of $G_n$. The transporter formula in the next statement proves faithfulness on every Hom-set.
\end{proof}

\statement{F.7}{Proved}{Finite quotient categories and the transporter completion agree}
For representatives $\sigma_s$ of the stable labelled trees put
$A_{st}=\{g\in G_n:g\sigma_s\supseteq\sigma_t\}$.
For each finite-index normal subgroup $V\lhd G_n$ let $p_V:G_n\to G_n/V$, and define a finite category $C_{n,V}$ by
\[
\operatorname{Ob}C_{n,V}=\{s\},\qquad
\Hom_{C_{n,V}}(s,t)=p_V(A_{st})/p_V(\Delta_{\sigma_s}),
\quad[h][g]=[hg].                                            \tag{F.5}
\]
If $V'\subseteq V$, its transition functor is induced by $G_n/V'\to G_n/V$. These quotient categories are cofinal among all finite congruence quotients of $C_n^0$. Thus
\[
\widehat C_n:=\varprojlim_V C_{n,V},\qquad
\Hom_{\widehat C_n}(s,t)=\overline{A_{st}}/\overline\Delta_{\sigma_s} \tag{F.6}
\]
materializes the strict pro-object of F.2. In the full profinite multicurve model this is the candidate transporter formula from the prompt:
\[
\Hom(\sigma,\tau)=
 \{g\in\widehat G_n:g\sigma\supseteq\tau\}/\overline\Delta_\sigma.
\]
It is a category in profinite sets, and the discrete comparison is faithful for every $n\geq3$.
\begin{proof}
The group-theoretic identity $g^{-1}\Delta_{\sigma_t}g\subseteq\Delta_{\sigma_s}$ in F.1 proves composition in (F.5) well-defined even after inserting elements of $V$, because $V$ is normal. Associativity and identities descend from multiplication. The stated transitions are actual functors, not merely outer homomorphisms.

If $A_{st}$ is nonempty, pure mapping classes preserve every split of a tree. Exactly one subset $\eta\subseteq\sigma_s$ has the splits of $\sigma_t$. Choose $r$ with $r\eta=\sigma_t$. Then
$A_{st}=H_{\sigma_t}r$. Hence its closure is $\overline H_{\sigma_t}r$, and the compact coset quotient gives (F.6). Conversely, a profinite transporter must send that same subset to $\sigma_t$, by the discrete finite set of split labels, so lies in this closed coset. This proves agreement with the proposed formula, including the empty Hom-sets.

To prove cofinality, let $F:C_n^0\to D$ be a functor to a finite category. Its restriction to $\Aut(s)=H_{\sigma_s}/\Delta_{\sigma_s}$ has finite-index kernel. By F.6 this kernel pulls back to a subgroup of $H_{\sigma_s}$ containing $V_s\cap H_{\sigma_s}$ for some finite-index normal $V_s\lhd G_n$. There are only finitely many $s$, so take $V=\bigcap_sV_s$. Suppose $g,g'\in A_{st}$ have the same class in (F.5). For some $d\in\Delta_{\sigma_s}$,
$h=g'd^{-1}g^{-1}\in V\cap H_{\sigma_t}$. The restriction of $F$ kills $h$, and $d$ already represents the identity at $s$, so $F([g'])=F([g])$. Thus $F$ factors through $C_{n,V}$. This is the missing cofinality assertion that could not be deduced from residual finiteness alone.

Each Hom in (F.6) is a closed subset of the profinite coset space $\widehat G_n/\overline\Delta_{\sigma_s}$, and multiplication is continuous. The full object space is the finite disjoint union $\coprod_s\widehat G_n/\overline H_{\sigma_s}$. Its transporter category is continuously equivalent to the skeleton just defined. To see that continuous choices are possible here, use the countable descending basis of open normal subgroups of the finitely generated profinite group and choose compatible transversals at successive finite quotients. Their inverse limits provide continuous transporters on each orbit. The choices are not canonical, but change the model by continuous natural equivalence. Faithfulness follows from the closedness of $\Delta_{\sigma_s}$ proved in F.6.
\end{proof}

\statement{F.8}{Proved}{Forgetting, clutching, and relabelling survive completion}
Every specified forgetful functor $Q:C_{n+1}\to C_n$, and every clutching functor, extends to the strict pro-completions by its universal property. On the transporter model, forgetting is induced by the completed pure Birman homomorphism $\widehat\phi:\widehat G_{n+1}\twoheadrightarrow\widehat G_n$ and the continuous map on profinite multicurves. In degeneration variance it is $\widehat Q^{\op}:\widehat T_{n+1}\to\widehat T_n$. Relabelling remains a pseudo-action with its specified natural isomorphisms; choosing a skeleton does not manufacture a splitting of a mapping class permutation extension.
\begin{proof}
Pulling a finite congruence back along a functor gives a finite congruence. This defines the pro-morphism together with its comparison maps at all stages; composition is respected by functoriality of pullback. On a discrete multicurve forgetting drops peripheral components and identifies curves that become isotopic. The images of its twists lie in the multitwist subgroup of the resulting multicurve, which proves well-definedness on cosets. The continuous extension exists by F.7 and agrees on a dense subcategory, hence agrees with the universal extension. No arbitrary point of the profinite curve complex is assigned an image without this continuity argument.

For clutching, the discrete full replete subcategory consists of the objects whose tree contains the prescribed split. Cutting along it identifies the relevant stabilizer quotient and its transporter Hom-sets with the product categories of the components. Their full induced profinite topologies, proved in F.6, imply that completion is still fully faithful. One can also apply the finite congruence argument to this product: finite functors on finitely many factor Hom-sets are controlled by a common finite quotient. This proves the product identification in genus zero. In $T$-variance it is a cosieve; in $C$-variance it is a sieve.

A change of marking conjugates a group homomorphism by its transporter, which is a natural isomorphism of the corresponding functors. Retaining these isomorphisms and their cocycle identity proves pseudoequivariance for permutations and for repeated forgetting. It is stronger data than a list of homomorphisms modulo unrelated inner automorphisms.
\end{proof}

\subsection{F.3. Geometric comparison and the Galois action}

\statement{F.9}{Known}{Stratified Riemann existence supplies the geometric comparison}
Let $X_n=\overline{\mathscr M}_{0,n}$ over $\mathbb Q$, stratified by its stable labelled trees. Fix $\overline{\mathbb Q}\hookrightarrow\mathbb C$. The finite boundary-constructible \etale sheaves on $X_{n,\overline{\mathbb Q}}$, with their stalk functors at one chosen geometric point of each stratum, are equivalent to
\[
 \Fun_{\cts}(\widehat C_n,\mathrm{FinSet}),                    \tag{F.7}
\]
with their evaluation functors. The finite-category exodromy object reconstructed from these data is therefore equivalent to the strict completion in F.7. At the full stratified $\infty$-categorical level, the precise theorem is
\[
 \operatorname{Gal}(X_{n,\mathbb C}/P_n)
    \simeq\operatorname{Exit}^{P_n}(X_n(\mathbb C))^{\wedge}_{\pi}, \tag{F.8}
\]
where the right-hand completion is in profinite stratified spaces. Here $P_n$ denotes the finite poset of strata, not the mapping class group $P_{g,n}$. This is \citet[Proposition 12.6.4 and Corollaries 12.6.6--12.6.7]{AF_BGH2020}. We take its finite-category, finite-set exodromy part in (F.7); we do not replace its full $\infty$-categorical completion by a limit of ordinary categories without a further comparison.
\begin{proof}[Precise deduction at the level used]
Base extension between algebraically closed characteristic-zero fields preserves finite \etale coverings of the strata and their constructible gluing data. Apply the analytic/\etale constructible comparison of the cited proposition and Jansen's exodromy theorem F.1. Their restriction to finite stalks gives $\Fun(C_n^0,\mathrm{FinSet})$. By F.2 and F.7 every such functor factors through a finite quotient category, with its maps and natural transformations. This yields (F.7), including compatibility with stalks.

For a finite category, the natural transformations between evaluation functors on its finite covariant representations recover its Hom-sets by the covariant representables. Passing through the actual finite quotient system gives the corresponding profinite Hom-sets by inverse limit. Thus the finite representations \emph{with their evaluation functors} recover precisely (F.6). This proves the claimed comparison rather than inferring it from equality of ordinary fundamental groups. Formula (F.8) is the stronger theorem at its own level. It is not a claim that the unstratified compactification is a $K(\pi,1)$, nor that the full point-and-specialization category $\operatorname{Gal}(X)$ equals the object for one fixed stratification. The latter object is the boundary-constructible localization/spectrification denoted $\operatorname{Gal}(X/P_n)$ in the source.
\end{proof}

\statement{F.10}{Proved}{Continuous outer Galois action on the completed category}
Galois descent defines a coherent continuous action
\[
 BG_{\mathbb Q}\longrightarrow\operatorname{Aut}(\widehat C_n)
\]
in the $2$-category of profinite categories and continuous equivalences. On equivalence classes it gives a homomorphism
$G_{\mathbb Q}\to\pi_0\operatorname{Aut}(\widehat C_n)$.
It is compatible with forgetting, clutching, and relabelling, and with the opposite category $\widehat T_n$. On a local normal-crossings inertia group at a rational boundary point it acts, modulo simultaneous geometric conjugation, by
\[
 t_e\longmapsto t_e^{\chi(\sigma)} \quad\text{for every node }e. \tag{F.9}
\]
For $n\geq4$ the outer action on $\widehat C_n$ is faithful.
\begin{proof}
Every $\sigma\in G_{\mathbb Q}$ acts on $X_{n,\overline{\mathbb Q}}$ over $\mathbb Q$ and hence, by inverse pullback $(\sigma^{-1})^*$, on finite boundary-constructible sheaves; this convention gives a left action. The natural associativity of pullback supplies the coherence for products of $\sigma$'s. Each finite sheaf and morphism descends to a number field, so this action is continuous on the finite data. Transfer by (F.7), including the stalks and their transports, gives the displayed outer categorical action. Changing the chosen geometric paths changes it by continuous natural isomorphism. It does not alter its class in $\pi_0$.

The geometric maps forgetting or clutching markings and the relabellings are defined over $\mathbb Q$. Their pullbacks commute coherently with this descent action, proving compatibility. Locally at a rational boundary point there are smoothing coordinates $q_e$ with divisor $q_e=0$. In characteristic zero Kummer covers $q_e^{1/m}$ have inertia $\mu_m$; Galois acts on their roots of unity by $\chi\bmod m$. Taking all $m$ gives (F.9). If parameters, tangential points, or paths are changed, the common local inertia group is conjugated, which explains the stated ambiguity; there is no preferred global generator independent of these choices.

For $n=4$, the open-stratum group is $\widehat G_4\simeq\widehat F_2$ for $\mathbb P^1\setminus\{0,1,\infty\}$, and Belyi faithfulness, proved or cited precisely in D, detects $G_{\mathbb Q}$. For $n>4$, forgetting to four markings is a surjective map on open-stratum groups and intertwines the Galois action. If the upstairs action were inner, its induced downstairs action would be inner, so faithfulness transfers in the needed direction. For $n=3$ the labelled category is a point and this action is trivial. Simple connectedness of the compactification does not contradict these statements: it concerns its unstratified fundamental group, which is a different quotient of the category.
\end{proof}

\statement{F.11}{Proved}{The failure of smoothness persists profinitely}
For forgetting the fifth marking at the stratum $q_{12,34}$, the continuous lifting groupoid in the STRATA smoothness test is equivalent to
$B\widehat{\mathbb Z}$. Its stabilizer is
\[
\ker(\widehat\phi)\cap\overline\Delta_q
  =\{(a,-a):a\in\widehat{\mathbb Z}\}
  =\overline{\langle t_1t_2^{-1}\rangle}.                      \tag{F.10}
\]
Thus the forgetful functor fails the analogous homotopy-smoothness test with continuous finite derived coefficients. Galois preserves this procyclic stabilizer and acts on it by the cyclotomic character, up to simultaneous conjugation.
\begin{proof}
The two distinct twists upstairs generate $\widehat{\mathbb Z}^2$ by F.5. Both map to the same primitive twist downstairs, so the exponent map is $(a,b)\mapsto a+b$. The only lift targets in the relevant test are open objects: a curve disappearing on forgetting $5$ would have split $i5\mid\{1,2,3,4\}\setminus\{i\}$, which occurs in neither of the two specified source splits. This is exactly the discrete split-label argument, and it persists on the profinite complex because the labels form a finite discrete invariant. The kernel $K=\ker\widehat\phi$ acts transitively on the lift objects, since $\widehat\phi(\overline\Delta_q)$ is the full downstairs twist group. Its stabilizer is (F.10), so the action groupoid is equivalent to $B\widehat{\mathbb Z}$.

For every prime $\ell$, its continuous cohomology has
$H^1_{\cts}(\widehat{\mathbb Z},\mathbb F_\ell)
 =\Hom_{\cts}(\widehat{\mathbb Z},\mathbb F_\ell)\simeq\mathbb F_\ell$.
Consequently this pro-homotopy type is not terminal and the required homotopy-cofinality condition fails. We are not claiming that connectedness alone fails, or that ordinary set-valued colimits detect this $H^1$. At the rational normal-crossings point, (F.9) uses one exponent on both branches; hence $t_1t_2^{-1}\mapsto(t_1t_2^{-1})^{\chi(\sigma)}$ under a compatible local lift. This also proves the statement for the corresponding degeneration functor, with variance reversed. The program checks the kernel and its cyclotomic stability in every finite modulus $2\leq m\leq24$; the displayed algebra proves it in the inverse limit.
\end{proof}

\subsection{F.4. Boundary-category automorphisms and the two levels}

\statement{F.12}{Proved}{Restriction to the open stratum detects every categorical automorphism}
Let $n\geq4$. A continuous autoequivalence of $\widehat C_n$ fixing the labelled stratum types is determined up to natural isomorphism by its outer action on $\widehat G_n=\Aut(u)$. Its image consists precisely of the outer automorphisms preserving the conjugacy classes of all boundary multitwist subgroups. Such an automorphism lifts if it preserves the associated profinite multicurve incidence. An inertia-preserving automorphism does preserve that incidence, by the commuting-twist characterization of simplices. These assertions hold for the continuous materialization and the strict pro-object alike.
\begin{proof}
For each object $s$ choose the standard arrow $j_s:s\to u$ represented by $1$. Its orbit is the left $\widehat G_n$-set
$\Hom(s,u)=\widehat G_n/\overline\Delta_{\sigma_s}$.
Its stabilizer therefore recovers the conjugacy class of $\overline\Delta_{\sigma_s}$. An autoequivalence must preserve these classes. Suppose first that its outer action at $u$ is trivial. Conjugate the functor by a natural isomorphism at $u$ to make the action literally the identity. Its map on $\Hom(s,u)$ is then an equivariant automorphism of a transitive profinite $\widehat G_n$-set, hence is given by an element of
$N_{\widehat G_n}(\overline\Delta_{\sigma_s})/\overline\Delta_{\sigma_s}$.
By F.5 this is exactly $\Aut(s)$. A natural change at $s$ therefore makes the map on $\Hom(s,u)$ the identity. Postcomposition with $j_t$ embeds $\Hom(s,t)$ into $\Hom(s,u)$ by the transporter formula. Thus the functor is the identity on every Hom-set. This proves injectivity after natural isomorphism.

Conversely let $\alpha\in\Aut(\widehat G_n)$ preserve the inertia subgroups and labelled strata. A profinite simple curve is determined by its inertia subgroup, and a finite set of distinct curves is a simplex precisely when the corresponding twist subgroups commute. The injectivity and commutation assertions are \citet[Theorem 4.11, Corollary 4.12, and Proposition 7.2]{AF_BoggiFunar2023}; the exact resulting extension to all inertia groups is also \citet[Proposition 3.5]{Boggi2026Automorphisms}. Hence $\alpha$ induces a continuous simplicial map $f$ on profinite multicurves with
$\alpha(\overline\Delta_\sigma)=\overline\Delta_{f\sigma}$ and
$f(g\sigma)=\alpha(g)f(\sigma)$.
The formula $(\sigma,[g],\tau)\mapsto(f\sigma,[\alpha(g)],f\tau)$ is consequently a well-defined continuous functor on the full transporter category. Its inverse is induced by $\alpha^{-1}$. Transferring to the skeleton gives the desired autoequivalence, independent of the transporters up to natural isomorphism. Inner $\alpha$ give naturally isomorphic functors. The injection already proved gives uniqueness of this extension. F.2 proves agreement with pro-object automorphisms.
\end{proof}


\statement{F.13}{Proved}{The decorated two-level boundary tower has automorphism group $\GT$}
Let $\mathcal T^{\partial}_{4,5}$ consist of $\widehat T_4,\widehat T_5$, all five forgetting functors, clutching inclusions with their $\widehat T_3,\widehat T_4$ factors, and coherent relabelling pseudo-actions. Its automorphisms are continuous equivalences fixing labelled stratum types, preserving oriented boundary inertia with one exponent $\lambda\in\widehat\Z^\times$, and commuting with these operations through natural isomorphisms satisfying their coherence equations. Modulo coherent natural isomorphism,
\[
 \pi_0\Aut(\mathcal T^{\partial}_{4,5})\simeq\GT.             \tag{F.11}
\]
The same holds for the strict pro-objects. The coordinate identification uses the usual marking by oriented loops $x,y$ around $0,1$ and a compatible tangential base point. Changing the marking conjugates this description; the outer tower and arithmetic action are independent of these choices.
\begin{proof}
Restriction to the open objects is injective by F.12. At those objects the inertia and relabelling conditions are exactly the definition of $\Out_n^\#$ in \citet[definition in the introduction and Main Theorem]{HarbaterSchneps2000}. Their theorem identifies $\Out_4^\#$ with $\widehat{\mathrm{GT}}_0$ and $\Out_5^\#$ with $\GT$; forgetting gives the standard injection from the latter to the former. Thus the five-point outer automorphism determines the compatible four-point one.

Conversely their $\GT$ action preserves inertia, labelled types, and forgetting homomorphisms. F.12 extends it uniquely, up to natural isomorphism, to both boundary categories. On a divisor $d_{ij}$, forgetting $i$ or $j$ restricts to the identity of its four-point factor after relabelling; its three-point factor is trivial. Compatibility with forgetting therefore identifies the induced action with the four-point action, giving clutching compatibility. At two-node strata the component factors are all three-point categories. F.12 already forces the crossing bimodules and their compositions.

The compatibility isomorphisms are coherent. A natural automorphism of the identity on $\widehat C_n$ has open component in $Z(\widehat G_n)=1$; naturality along $j_s:s\to u$ and injectivity of postcomposition force every other component to be $1$. For forgetting, surjectivity of the open homomorphism gives the same argument, since its image has centralizer $Z(\widehat G_4)=1$. For clutching, the image is the full automorphism group of the target stratum, again centre-free or trivial. Thus compatibility isomorphisms in this two-level diagram are unique and their coherence equations hold. Relabellings are equivalences and satisfy the same uniqueness. Taking opposites preserves the result; F.2 compares the pro-object and continuous versions.
\end{proof}

\statement{F.14}{Proved}{The full genus-zero boundary tower and its arithmetic subgroup}
Let $\mathcal T^{\partial}_0$ contain all $n\geq3$ with the same decorations and all forgetting and clutching functors. Restriction is an isomorphism
\[
 \pi_0\Aut(\mathcal T^{\partial}_0)
 \xrightarrow{\sim}\pi_0\Aut(\mathcal T^{\partial}_{4,5})
 \simeq\GT.                                                \tag{F.12}
\]
F.10 gives the usual faithful arithmetic inclusion $G_{\Q}\hookrightarrow\GT$. Equating its image with all automorphisms is the open equality E.8. The categorical group thus identifies the early normalizer-type group E.7; the revised arithmetic-image group is the subgroup E.2.
\begin{proof}
The Harbater--Schneps Main Theorem and its stability corollaries identify $\Out_n^\#$ with $\GT$ for all $n\geq5$, compatibly with forgetting. Their tower action supplies existence on open strata; F.12 extends it uniquely at every level. The group of a boundary stratum is $\prod_v\widehat G_{\operatorname{val}(v)}$. Each factor is recovered by forgetting all but its markings and one marking in every branch attached to it. This restricts to the corresponding projection. Forgetting compatibility therefore forces the product of factor actions and proves clutching compatibility. All these groups have trivial centres, including the trivial valence-three factors, so F.13 proves coherence. The extension from two levels is unique.

Arithmetic descent and the usual arithmetic action coincide on the four-point open stratum by fundamental-group comparison, hence everywhere by uniqueness. Faithfulness is F.10. E.8 is the additional surjectivity assertion; no profinite equality follows from a rational Lie-algebra calculation.
\end{proof}

\subsection{F.5. Positive genus, isotropy, and the rational comparison}

\statement{F.15}{Proved}{Completion for the elliptic types retains stack isotropy}
For $(g,n)=(1,1),(1,2)$ take Jansen's labelled stable-curve category in exit variance, with a finite skeleton indexed by stable labelled dual graphs. Its strict completion F.2 is a category in profinite sets, faithful on discrete Hom-sets, and equivalent to the profinite-multicurve transporter category for $\widehat P_{g,n}$. At $\sigma$,
\[
 1\longrightarrow\prod_v\widehat P_{g_v,n_v}
 \longrightarrow\Aut_{\widehat C_{g,n}}(\sigma)
 \longrightarrow A_\sigma\longrightarrow1,                  \tag{F.13}
\]
where $A_\sigma$ is the finite image of the discrete stabilizer on the dual graph and half-edges, fixing the original marked legs. No splitting is chosen. These categories carry F.9--F.10's finite constructible-sheaf comparison and outer Galois action, now for moduli \emph{stacks}.
\begin{proof}
Jansen's Corollary 6.6 and Construction 3.18 apply in all genera, retaining graph automorphisms, and give $\Aut(\sigma)=H_\sigma/\Delta_\sigma$. The finite-index subgroup of $H_\sigma$ fixing components and half-edges has the cutting extension by $\Delta_\sigma$ with quotient $\prod_vP_{g_v,n_v}$. Full and procongruence completions agree in genus one and its genus-zero components. Boggi's stabilizer theorem and F.6's short-five-lemma argument identify this subgroup's completion with its closure. Passing through its finite-index extension does the same for $H_\sigma$ and gives (F.13). Multitwist independence and conjugacy separability give closedness of $H_\sigma,\Delta_\sigma$.

Here a transporter may be a finite union of cosets $H_{\sigma_t}r$, since different submulticurves of $\sigma_s$ may have the type of $\sigma_t$. Closedness of $H_{\sigma_t}$ supplies a finite quotient separating these finitely many cosets. Refine to it before the kernel argument in F.7. Equality modulo $V\Delta_{\sigma_s}$ then forces the same transporter component, since $\Delta_{\sigma_s}$ fixes all submulticurves. Hence $g'd^{-1}g^{-1}\in H_{\sigma_t}$ and the cofinality proof applies. This proves the transporter description and faithfulness, retaining graph symmetries.

Take an \etale atlas and its simplicial scheme of overlaps. Finite constructible sheaves satisfy descent. The analytic/\etale comparison on each overlap commutes with all face and degeneracy maps; its descent limit gives the stack comparison. This agrees with \citet[Proposition 13.9.17 and Corollary 13.9.19]{AF_BGH2020}. Galois pullback on the descent data supplies the coherent action and retains finite stabilizers.
\end{proof}

For $(1,1)$ the open group is $\widehat{\mathrm{SL}_2(\Z)}$, with central involution. Its boundary graph has one rational vertex, a labelled leg and a loop: exchanging the half-edges gives $A_\sigma=C_2$, while $P_{0,3}=1$. The boundary object's group is therefore $C_2$. For $(1,2)$ a rational tail carrying both markings leaves an elliptic component with one node, hence a factor $\widehat P_{1,1}$; a nonseparating node instead has a four-point rational normalization and an exchange of node branches. Formula (F.13) retains these stabilizers, including their $2$- and $3$-primary information.

\statement{F.16}{Open}{The precise remaining positive-genus automorphism comparison}
We have not established that the automorphism group of the two elliptic boundary categories, forgetting, and all boundary clutchings is $\GT$. The precise comparison asks whether restriction from the decorated procongruence outer tower E.13 to this diagram is an isomorphism, with identical oriented inertia and relabelling data. Injectivity requires detection at these levels; surjectivity requires extension of every such pair to the full decorated tower. We do not claim that existence of $\GT$ actions in positive genus remains open: E.12--E.13 give current results and their exact morphism collections.
\begin{proof}[Where our argument stops]
Formal completion and descent do not compute this automorphism group. Since $\widehat P_{1,1}$ has centre $C_2$, F.13's uniqueness-of-coherence argument needs additional work. Elliptic presentations also include nonseparating inertia and handle relations, whereas the genus-zero proof uses labelled splits and separating twists. E.10--E.11 concern specified open groups, and do not alone settle coherence and extension across all boundary bimodules. One must compare the precise decorated towers, retaining the central involution. The smallest diagram not settled here is exactly $(1,1)$--$(1,2)$. The label \emph{Open} records an unresolved comparison in this audit, not an assertion that no theorem exists elsewhere.
\end{proof}

\statement{F.17}{Known}{The rational operadic analogue and its exact limitation}
Over a characteristic-zero field $k$, define the rational analogue to be the parenthesized braid groupoids completed relative to their finite permutation quotients, with pure braid groups replaced by their Malcev prounipotent completions. Retain operadic composition, strand deletion and doubling, the braid symmetry, and the distinguished oriented inertia. Their structure-preserving automorphism group is $\mathrm{GT}(k)$. In the parenthesized chord-diagram realization it is $\mathrm{GRT}(k)$; an associator compares the realizations. The four- and five-point constraints are the symmetry equations and pentagon, and the identity-multiplier Lie algebra is the weight completion of $W=\mathfrak{grt}_1$. The exact sources are \citet[\S\S5--6, Propositions 5.1 and 5.5]{Drinfeld1990} and \citet[Proposition 3.4, Theorems 1--3, Corollary 4.1]{BarNatan1998}.
\begin{proof}[The corollary supplied by the cited formulation]
The algebraic completion is universal for homomorphisms from the pure groups to unipotent algebraic groups, while the discrete objects and finite permutation data are retained. The cited theorem computes automorphisms preserving the indicated operations, and identifies their equations on the two lowest nontrivial levels. It is the operadic, algebraically enriched analogue of the decorated tower automorphism calculation F.13. The change from braids to infinitesimal braids depends on an associator and is not canonical. Applying E.19--E.20 gives its motivic corollary: the mixed Tate group injects, and its equality with the full rational automorphism group follows conditionally from the all-weight upper bound in family 008. That last assertion remains explicitly \emph{claimed (OpenAI, 2026, unrefereed)} with the audit in E.26--E.28.

This does not prove that termwise Malcev completion of every Hom-set in the boundary category produces those parenthesized groupoids. A boundary quotient $G/\Delta_\sigma$ is generally a homogeneous space, not a torsor for $G$, and the exit category has noninvertible smoothing arrows. The parenthesized braid categories are groupoids. For instance the arrow from a boundary object to the open object of $C_4$ has no reverse arrow; completing its groups does not create one. Thus those bare categories cannot be identified. A further comparison of their \emph{structure-preserving automorphism groups} would require an algebraically enriched boundary construction and a reconstruction proof analogous to F.12; it is not imported silently from the groupoid theorem. This is the exact scope of the rational corollary supplied here. Locally constant coefficients in $D(\Q)$ do not by themselves implement Malcev completion, and neither rational construction yields $G_{\Q}=\GT$.
\end{proof}

\statement{F.18}{Open}{A stronger higher-completion comparison is not assumed}
One can also ask whether the nerve of the strict finite-category pro-object F.7, as a pro-object of stratified spaces, equals the full profinite stratified $\infty$-category $\operatorname{Gal}(X_{n,\C}/P_n)$. This requires comparison for all finite homotopy types and their constructible gluing, beyond finite sets. We do not establish it here. The first nontrivial test is $n=4$, whose boundary gluing retains a two-dimensional unstratified homotopy class even though its stratum groups are free.
\begin{proof}[Exact limitation]
BGH's Corollary 12.6.7 identifies the right-hand object with profinite stratified completion of the \emph{nerve} of $C_n$. F.2 instead takes strict quotient categories before nerves. Interchanging them requires cofinality or goodness for representations in all $\pi$-finite spaces. F.9's finite representables prove only the $0$-truncated coefficient comparison. All automorphism statements F.12--F.14 concern the strict completion defined here, whose finite-set comparison is proved. We leave the stronger statement unresolved without asserting a counterexample. No unchecked contractibility assertion from STRATA F.10 is used.
\end{proof}


\clearpage\appendix
\section{Verification protocol and limits of the deliverable}

The accompanying register has one row for every numbered statement. Its status attaches to the exact assertion, including its hypotheses. In particular a proved conditional implication does not certify its antecedent. A bibliography entry marked \emph{verified} means that the original source was opened and the cited passage located; it does not mean that every theorem in that source has been re-proved. The report lists the portions actually read and the depth of verification. Mathematical support for a doubtful transcription is recorded separately from visual evidence, and neither is described as a final editorial reading.

\subsection{The three attached 2026 proof families}
\begin{center}\small
\begin{tabularx}{\textwidth}{@{}L{.13\textwidth}L{.25\textwidth}Y@{}}
\toprule
Family & Exact use & Boundary of this audit\\\midrule
008 & E.19--E.28; F.17 & We checked the dimension argument conditional on the all-weight upper bound, read the mod-two lattice and deletion reductions, and supplied a bounded integral computation. The all-weight upper bound is not certified by that computation. The attached Lean comparator contains a placeholder; it is not a kernel-checked proof. Subsequent October preprints are recorded, with their own verification scope.\\
019 & D.7--D.10; E.17 & We followed the deduction from the completion claim to the local section statement, and the conditional global and modular consequences. The patching, analytic convergence and all-valuation uniformization inputs are not promoted to independently verified theorems. Affine and moduli-stack sections are not covered by the proper-curve theorem.\\
031 & D.11--D.13 & We prove preservation of the geometric kernel independently, then deduce the absolute curve Hom theorem conditionally from the claimed Galois open-Hom theorem. We separate that implication from the new local-to-global argument in the supplied preprint.\\\bottomrule
\end{tabularx}
\end{center}

The report gives the precise first unresolved step for each imported claim. We found no verified counterexample to the main theorem of one of these three manuscripts. Failure to complete its proof audit is not a refutation. Conversely, publication of a later proof treatment is reported as evidence of a subsequent claim, without representing an unread argument as checked.

\subsection{Computations and proof boundaries}
The programs in \texttt{code/} use only the Python standard library. Their saved outputs include the finite category and transition checks for F, the stable-tree counts and forgetting map, the finite inertia-kernel calculations, the permutation and finite-linear-group counterexamples of B/C, the finite section/normalizer checks of D, and the bounded integral deletion calculation of E. The infinite assertions are proved in the text or explicitly imported from the specified primary theorem. A finite enumeration never substitutes for density, convergence, an all-weight inequality, or a statement about all open subgroups.

The genus-zero categorical theorem concerns actual pro-objects of finite categories and their continuous materializations, with the stated decorations and natural \emph{isomorphisms}. It is not a theorem about an unspecified inverse limit in a category obtained by inverting all natural transformations. Its comparison with finite constructible sheaves includes the stalk functors. F.18 records the additional higher-completion comparison which the finite-set argument alone does not prove.

\subsection{What remains unresolved}
The profinite arithmetic equality $G_{\Q}=\GT$ remains open. The higher-genus congruence kernel remains in B.3. Other exact unrestricted comparisons left unresolved in this audit are B.7, B.14, C.8, D.23, E.14, F.16 and F.18. Their wording distinguishes an established open problem from a comparison for which our search and argument did not supply a theorem. The 2026-dependent claims are separated in D.7, D.13 and E.26. All these statements remain in the results register; no question has been silently dropped after a reduction to a special case.

We claim no priority for the arguments, counterexamples, constructions, or reformulations in this paper. The research report records the search scope before any possible inference of novelty. This paper, its code and its source audit are AI-produced by GPT-6 Astra (OpenAI, ChatGPT Pro mode).


\clearpage
\bibliographystyle{apalike}
\bibliography{LM3_TEICHMULLER_TOWER_references}
\end{document}
